\documentclass[UTF8, fontset=macnew, a4paper, 12pt]{ctexart}

\usepackage{geometry}
\geometry{top=2.5cm, bottom=2.5cm, left=2.5cm, right=2.5cm}
\usepackage{amsmath, amssymb}
\usepackage{booktabs}
\usepackage{multirow}
\usepackage{tabularx}
\usepackage{graphicx}
\usepackage{caption}
\usepackage{xcolor}
\usepackage{float}
\usepackage{enumitem}
\usepackage{fancyhdr}
\usepackage{setspace}
\usepackage{pifont}
\usepackage[hidelinks]{hyperref}

\graphicspath{{../output/figures/}}
\captionsetup{font=small, labelfont=bf, labelsep=quad}
\setlist{itemsep=2pt, parsep=0pt, topsep=4pt}
\onehalfspacing

\ctexset{
  section/format = \raggedright\zihao{4}\bfseries,
  section/name   = {},
  section/number = \bignum{\value{section}}、,
  subsection/format = \raggedright\zihao{5}\bfseries,
  subsection/number = \chinese{subsection},
  subsection/aftername = 、,
  subsubsection/format = \raggedright\normalsize\bfseries,
}
% 節次採大寫數字（壹貳參…），與內文的「第壹節」及範例論文的體例一致；
% 小節仍用 \chinese（一二三…），對應內文的「之一」「之二」。
\newcommand{\bignum}[1]{\ifcase#1\or 壹\or 貳\or 參\or 肆\or 伍\or
  陸\or 柒\or 捌\or 玖\or 拾\fi}
\renewcommand{\thesection}{\bignum{\value{section}}}
% 交叉引用子節時補上節次，否則 \ref 只印出「三」而無從判斷屬於哪一節。
% \p@subsection 是 LaTeX 寫入 .aux 時加在 \thesubsection 前的前綴。
\renewcommand{\thesubsection}{\chinese{subsection}}
\makeatletter
\renewcommand{\p@subsection}{\thesection 節之}
\makeatother
\setcounter{secnumdepth}{2}   % subsubsection 不編號（標題已自帶（一）（二））

\pagestyle{fancy}
\fancyhf{}
\fancyhead[L]{\small 小人口 Lee-Carter 模型的參數分層診斷與估計}
\fancyhead[R]{\small 碩士論文進度報告（1006 精選版）}
\fancyfoot[C]{\small \thepage}
\renewcommand{\headrulewidth}{0.4pt}

% U+2460~2463（\cn{1}~\cn{4}）不在 xeCJK 認定的 CJK 範圍內，會落到拉丁字型而缺字，
% 故改用 pifont 的圈號。
\newcommand{\cn}[1]{\ding{\number\numexpr171+#1\relax}}
\newcommand{\edag}{e_0^{\dagger}}
\newcommand{\LCe}{\mathrm{LC}e^{\dagger}}

\begin{document}

\begin{center}
{\zihao{3}\bfseries 小人口 Lee-Carter 模型的參數分層診斷與估計}\\[4pt]
{\zihao{5} —— 從資料的分布，決定哪些參數可估、該用哪個估計量}\\[4pt]
{\zihao{5}\textbf{第四版（1006）}：依 EDA 佈局與穩健方法架構重組，收斂研究主軸}\\[12pt]
{\zihao{5} （林子立）\quad 指導教授：余清祥\,教授}\\[4pt]
{\zihao{5} 國立政治大學統計學系}\\[4pt]
{\zihao{5} 中華民國 115 年 9 月 26 日}
\end{center}

\vspace{6pt}
\begin{center}
\begin{minipage}{0.92\textwidth}
\small
\textbf{摘要}\par\vspace{4pt}
我國多數鄉鎮市區的單一性別人數不足五萬，Lee-Carter 模型的參數在此規模下
出現有方向的偏誤。本文指出三個參數因\textbf{不同機制}失效：$\alpha_x$ 源於
零格替代與對數轉換、漂移項源於等權重分解對異質變異的敏感、$\beta_x$ 則受
訊噪比限制而在相變門檻以下不可還原。三者皆可由四項\textbf{不需真值的事前
診斷}辨認，並分別以偏誤減少、資訊加權與經驗貝氏收縮對應處理。\par\vspace{6pt}
\textbf{關鍵詞}：小區域估計、Lee-Carter 模型、偏誤診斷、偏誤減少、
經驗貝氏收縮、預測區間
\end{minipage}
\end{center}

\vspace{10pt}

\section{研究動機與目的}

平均餘命（Life Expectancy）是衡量地區發展與居民健康的重要指標，可協助政府
配置醫療與社會福利資源，也用於保險業的費率計算。其估計須以各年齡的死亡率
為基礎，而我國官方編製的生命表只涵蓋全國與縣市層級，並未涉及鄉鎮市區。
缺漏的原因不在資料不全，而在樣本規模：368 個鄉鎮市區中有 217 個（59\%）
的單一性別人數在兩萬以下、299 個（81\%）少於五萬人，此等規模下年齡別
死亡率的觀測值震盪劇烈，死亡率較低的青少年年齡組甚至經常沒有死亡人數
（余清祥等 2025）。

人數少所帶來的困難，通常被理解為估計值的震盪，亦即變異數過大。然而
Wang et al.（2018）與余清祥等（2021）指出，Lee-Carter（以下簡稱 LC）模型
的年齡參數在人數減少時還會出現有方向的偏誤，且在人數二十萬以下時特別明顯。
這項區別至關重要，因為變異數大可藉由增加資料或重複抽樣察覺，偏誤卻會在
所有樣本上朝同一方向偏移，單看估計結果無從發現，也無法靠加寬區間補救。
換言之，若小人口的問題確實含有偏誤成分，則補救的第一步不是讓估計更穩定，
而是找出偏誤來自模型的哪一個環節。

既有文獻對小人口的對策大致可歸為三類，而三類各自預設了一件本文的資料
未必具備的事。第一類是向外借資訊。Li and Lee（2005）的一致性 LC 模型
以一組相似人口的共同趨勢為基準，各人口只估計其相對偏離；
Wang et al.（2018）與 Chen et al.（2017）循此把縣市資料連結到全國，
Jarner and Kryger（2011）的 SAINT 模型則以一個大型參考人口的死亡率
作為貝氏架構的先驗。這一類作法的前提是存在一個死亡率改善型態可比的
參考人口，而鄉鎮市區之間的社會經濟差異甚大，此一前提在本文的層級上
是否成立本身即需檢驗。第二類是修勻。Eilers and Marx（1996）的
$P$-樣條經 Currie et al.（2004）、Delwarde et al.（2007）與
Camarda and Basellini（2021）引入死亡率模型後已成常規工具，其前提是
死亡率在年齡與年份方向上足夠平滑，使鄰近格的資訊可以互相支援；
但修勻處理的是變異數而非偏誤，本文第肆節將以實測說明此一界限。
第三類是改換估計架構。Brouhns et al.（2002）以卜瓦松概似取代 SVD，
Renshaw and Haberman（2003）與 Currie（2011）把 LC 納入廣義線性模型，
其前提是概似函數的最大值存在——而零格所導致的完全分離
（Albert and Anderson 1984）恰使此一前提失效。本文的取徑屬第三類，
但關注的正是這個前提失效之處。

須補充的是，這三類作法的評估幾乎都以「估計是否穩定」或「樣本外預測誤差
是否較小」為標準。Lee and Miller（2001）與 Booth et al.（2005）對 LC
各變體的系統性比較即屬此類。以預測誤差為標準的比較有其道理，卻無法分辨
誤差來自偏誤或變異數，也無法指出偏誤落在哪一個參數的哪一個年齡；
而如前所述，這個區分決定了補救的方向。本文因而不以預測誤差為主要標準，
改以「偏誤／變異數的分解」與「可事前計算的診斷」為主軸。

要追究這個環節，須先問 LC 模型做了什麼假設。Lee and Carter（1992）以一組
簡潔的雙線性結構描述死亡率的時間變化，三十年來成為死亡率推估的標準工具；
Basellini et al.（2023）在其三十週年回顧中，將此一地位歸因於四項優點：
結構簡潔，單一時間指標即可導出全年齡的預測；時間指標在多數應用中近乎線性，
故可用隨機漫步含漂移項外推；隨機架構可導出機率預測區間；以及純粹外推、
不需主觀判斷，因而具有可重複性。這四項優點皆已在國家級人口上反覆驗證，
因此本文要檢驗的不是 LC 模型是否正確，而是它在何種條件下適用。

四項優點之中，前三項都建立在同一個技術選擇上，即以奇異值分解（Singular
Value Decomposition，以下簡稱 SVD）求解。SVD 之所以方便，在於它是封閉解
且必然存在，但代價是它等價於無加權的最小平方，因而隱含了同質變異
（Homoscedasticity）與常態誤差的高斯架構——每一個年齡與年份構成的格子，
在目標函數裡被賦予相同的權重。

\subsection{資料所呈現的問題}\label{sec:intro-data}

前述代價在國家級人口上幾乎不必付。表~\ref{tab:datadesc} 為本文主分析資料
（全國女性 2001 至 2024 年）的基本統計，528 個格子中單格死亡數最少亦有
37 人、中位數 1{,}201 人，零死亡格為 0，各格所攜帶的資訊量雖有差異但同屬
可用。

\begin{table}[H]
\centering
\caption{主分析資料的分布特性（女性 2001$\sim$2024 年）}
\small
\begin{tabular}{lr}
\toprule
項目 & 數值\\
\midrule
資料維度 & 22 齡組 $\times$ 24 年 $=$ 528 格\\
單格死亡數：最小／中位數／最大 & 37／1,201／16,184\\
零死亡格數 & \textbf{0}\\
2024 年暴露數：最小／中位數／合計 & 2,883／552,075／11,871,065\\
死亡數最少者（全期合計） & 5--9 歲組 1,812 人\\
死亡數最多者（單年中位數） & 80--84 歲組 10,321 人\\
\bottomrule
\end{tabular}
\label{tab:datadesc}
\end{table}

但同一張表也顯示了一項在全國資料上即已存在、而在小人口會被放大的性質：
死亡數在年齡上的分布極不均勻，5 至 9 歲組的單年死亡數中位數為 69 人，
而 80 至 84 歲組達 10{,}321 人，相差約 150 倍。同一筆資料之中，不同年齡的
有效樣本數相差兩個數量級，估計精度自不可同日而語，等權重的目標函數因而
已被誤設。至於鄉鎮市區層級則完全不同：以單一性別兩萬人的規模推算，528 格
中平均有 206 格無死亡發生，約四成的格子不含任何資訊，同質變異假設在此
不是略有偏差，而是完全失效。

前述主張若僅靠引文與單一張描述統計表支持，尚不足以說明失效的形式與位置。
本文因此設計四項不需要真值、拿到資料即可計算的診斷，直接在資料上檢驗
該假設；其計算方式與判讀門檻詳見第\ref{sec:eda}，此處先呈現結果
（圖~\ref{fig:eda}），以說明本文問題意識的由來。

\begin{figure}[H]
  \centering
  \includegraphics[width=\textwidth]{figEDA_panel.png}
  \caption{四項事前診斷。橫向比較全國女性實際資料與兩組模擬小人口
  （以全國配適為真值，$D_{x,t}\sim\mathrm{Poisson}(N w_x m^{\mathrm{true}}_{x,t})$）。
  \cn{1}\cn{2}\cn{3} 為純粹的探索性分析；\cn{4} 的噪音上緣需一次預備配適以估計 $\sigma$。}
  \label{fig:eda}
\end{figure}

第一項診斷觀察各年齡單年死亡數的量級，用以判斷等權重是否被誤設；如前所述
全國女性的落差已達 150 倍，而模擬五萬人時最低的年齡組降至 0。第二項診斷
觀察零死亡格落在哪些年齡，模擬五萬人時 1 至 4 歲、5 至 9 歲、10 至 14 歲
三組分別有 74\%、77\%、78\% 的年份死亡數為零，而 50 歲以上各組皆為零，
可見零格並非散布於全表，而是集中於特定年齡——這一點後文將顯示具有決定性，
因為它使偏誤的位置可以事前指認。第三項診斷以對數死亡率沿時間的二階差分
估計噪音變異，三個人口規模的點都貼合理論參考線 $1/D$ 且橫跨四個數量級，
是異質變異的直接證據，也量化了它的嚴重程度。第四項診斷在奇異值的碎石圖上
加一條純噪音所能產生的上緣，最大奇異值穿出這條線，年齡參數才可能還原，
實測全國女性的比值為 6.53，模擬二十萬為 1.19，模擬五萬僅 0.43。

此處另有一個反直覺之處值得先行指出：最大奇異值的絕對值完全不能判讀。
全國女性的最大奇異值為 3.44，而模擬五萬人的最大奇異值反而是 4.23——
小人口的奇異值被噪音灌水，必須與上緣相比才有意義。直覺上會以為奇異值大
即訊號強，實則不然。

\subsection{失效的代價因參數而異}

假設不成立本身仍不足以支持更換方法，畢竟許多假設都不完全成立而方法照常
好用；必須進一步證明失效有代價，而且代價的類型決定了補救的方式。本文的
T-ratio 診斷（第\ref{sec:diag}）顯示，三個參數的病因並不相同。年齡水準
參數 $\alpha_x$ 是偏誤主導，人數一萬時 5 至 9 歲組的 $|T|$ 高達 56.4，
隨人口增大而系統性下降，人數二十萬以上後全面落入 $\pm2$ 之內。年齡敏感度
參數 $\beta_x$ 在小樣本下則是變異數主導，人數二十萬以下時 $|T|$ 不及 0.25，
須至五十萬以上變異數收斂後偏誤才顯現，人數一百萬時 70 至 74 歲組的 $|T|$
為 1.87；小樣本時 $T$ 接近零並非估得好，而是變異數大到把偏誤掩蓋。時間
參數的漂移項性質與 $\alpha_x$ 同為偏誤主導，但來源並不相同：本文第伍節的
階梯式分解中，對真值加入對稱高斯噪音後再做 SVD 的變體完全不涉及對數轉換、
亦不存在零格，卻使 $\alpha_x$ 的最大偏誤由 0.743 降至 0.021，而漂移項的
偏誤仍達 0.347，甚至略大於標準 LC 的 0.305。

三項診斷合起來給出本文的核心判斷：$\alpha_x$ 的偏誤來自對數轉換與零格替代，
漂移項的偏誤來自等權重分解對異質變異的敏感，$\beta_x$ 在小樣本下則受訊噪比
限制，在相變門檻以下根本不可還原。三者的補救方式互不相同，偏誤必須找出來源
加以校正，變異數才可藉收縮處理，而不可還原的部分則應誠實地不予估計。更值得
注意的是，這三種失效都在原始資料中留有可事前辨認的徵兆：第二項診斷指出
$\alpha_x$ 的偏誤會落在哪些年齡，第一與第三項診斷指出加權失衡的嚴重程度，
第四項診斷指出 $\beta_x$ 是否可還原。因此在配適任何模型之前，就能判斷哪些
參數可估、該用哪個估計量，這即是本文的主張。取得鄉鎮市區資料的使用者既無
真值亦無從進行模擬，可事前計算的指標因而比一個固定的人口數門檻實用得多，
也才能解釋為何同為五萬人，有的地區可估而有的不可。

\subsection{研究問題與本文架構}

依上述主張，本文提出四個依序遞進的研究問題。其一，四項事前診斷能否在配適
模型之前，正確指出哪些參數可估、偏誤會落在哪些年齡。其二，三個參數的失效
機制是否確實不同，對數轉換是否為偏誤的唯一來源。其三，依機制選用的估計量
各能改善多少，改善的位置是否與診斷所指一致，以及在什麼條件下失效。其四，
小人口的預測區間誤差來自何處，文獻長期主張的「LC 區間過窄」在此是否成立。

本文編排如下。第貳節定義記號與模型、回顧文獻，並於其中說明本研究的目的
與所欲暸解的四項議題；第參節說明研究資料的來源與分析邏輯、四項事前診斷的
計算方式、偏誤診斷與二因子模擬設計，回應問題一；第肆節以 Rabbi and
Mazzuco（2021）的還原與懲罰式估計為例，論證在高斯架構內修補的限度；
第伍節以階梯式分解辨識偏誤的來源，比較各估計量的改善幅度與適用條件，
並以 M-分位數的連續掃描與逐年穩健化兩項檢驗界定穩健損失的作用範圍，
回應問題二與問題三；第陸節處理區間推論，回應問題四；第柒節逐一回應上述
四個問題，並討論本文與既有文獻的關係、研究限制與後續工作。

各節的鋪陳採同一個順序：先指出該節要處理的問題，再說明所採方法及其所依
的假設，接著給出實證結果，然後檢驗該方法的假設在本文資料上是否成立並
回應節首的問題，最後以小結收束並指出下一節所要處理的缺口。
第\ref{sec:gap}的表~\ref{tab:gap} 把各方法的原始應用場景與本文情境的
前提落差並列，該表右欄的每一項落差，在後續各節都有一個對應的檢驗。

\section{文獻回顧}

本節先定義全文使用的記號與模型（第一小節），回顧其兩條估計路線與小人口
情境的既有對策（第二、三小節），據以說明本研究的目的（第四小節），
最後回顧與該目的直接相關的方法論脈絡（第五小節）。第五小節的編排依
「能不能估、偏誤怎麼處理、變異數怎麼處理、損失函數要不要換、如何判斷
適用哪一種」的順序，與本文第伍節的實驗設計一一對應。

\subsection{記號與 Lee-Carter 模型}\label{sec:notation}

為免混淆，本文先統一以下記號。以 $x=1,\dots,A$ 表示年齡組的索引、
$t=1,\dots,T$ 表示年份的索引，本文的主要設定為 $A=22$ 個五齡組、$T=24$ 年。
純量以細體小寫表示，$D_{x,t}$ 為第 $x$ 個年齡組在第 $t$ 年的死亡數、
$E_{x,t}$ 為對應的暴露數（人年）、$m_{x,t}=D_{x,t}/E_{x,t}$ 為中央死亡率。
向量以粗體小寫表示並一律視為行向量，
\[
\boldsymbol{\alpha}=(\alpha_1,\dots,\alpha_A)^{\!\top},\quad
\boldsymbol{\beta}=(\beta_1,\dots,\beta_A)^{\!\top},\quad
\boldsymbol{\kappa}=(\kappa_1,\dots,\kappa_T)^{\!\top},
\]
亦即附下標的 $\alpha_x$、$\beta_x$、$\kappa_t$ 表示對應向量的第 $x$
（或第 $t$）個分量，不附下標的粗體符號則表示整個向量；以 $\mathbf{1}_A$
表示長度 $A$ 的全一向量，$\mathbf{1}_A/A$ 即各分量皆為 $1/A$ 的均勻向量。
矩陣以粗體大寫表示，$\mathbf{D}=(D_{x,t})$、$\mathbf{E}=(E_{x,t})$、
$\mathbf{M}=(m_{x,t})$。估計值於符號上方加帽，真值於必要時加上標
$\text{true}$；向量的 $L_1$ 與 $L_2$ norm 分別記為 $\|\cdot\|_1$ 與
$\|\cdot\|_2$。

Lee and Carter（1992）的模型設定為
\begin{equation}\label{eq:lc}
\log m_{x,t}=\alpha_x+\beta_x\kappa_t+\varepsilon_{x,t},
\end{equation}
其中 $\alpha_x$ 為第 $x$ 個年齡組的平均死亡水準、$\beta_x$ 為該年齡組對
時間趨勢的敏感度、$\kappa_t$ 為第 $t$ 年的死亡率時間指標。以矩陣寫成
$\log\mathbf{M}=\boldsymbol{\alpha}\mathbf{1}_T^{\!\top}
+\boldsymbol{\beta}\boldsymbol{\kappa}^{\!\top}+\boldsymbol{\varepsilon}$，
可見模型的秩為一——每一年的整條對數死亡率曲線，皆由單一純量 $\kappa_t$
沿固定方向 $\boldsymbol{\beta}$ 平移而得。此一結構在後文解釋數個現象時
會反覆用到。

式~\eqref{eq:lc} 對 $(\boldsymbol{\beta},\boldsymbol{\kappa})$ 的尺度與
平移並不唯一，須施加識別條件，本文沿用原始設定
\begin{equation}\label{eq:ident}
\mathbf{1}_A^{\!\top}\boldsymbol{\beta}=\sum_{x=1}^{A}\beta_x=1,
\qquad
\mathbf{1}_T^{\!\top}\boldsymbol{\kappa}=\sum_{t=1}^{T}\kappa_t=0 .
\end{equation}
式~\eqref{eq:ident} 的第一式在第肆節扮演決定性的角色，它使「朝零收縮」
的懲罰項退化為常數，因而完全失效。在卜瓦松架構下模型改寫為
$D_{x,t}\sim\text{Poisson}\!\left(E_{x,t}\exp(\alpha_x+\beta_x\kappa_t)\right)$，
以最大概似法估計，此時 $\varepsilon_{x,t}$ 不再需要同質變異的假設。

本文另須用到壽命差異 $\edag$，它描述死亡年齡分布的離散程度，定義為
$\edag=\int_0^{\infty} d(a)\,e(a)\,\mathrm{d}a$，其中 $d(a)$ 為死亡年齡
分布的密度、$e(a)$ 為 $a$ 歲的平均餘命，直觀而言即「每一位死者平均損失
的餘命年數」。相對於零歲平均餘命 $e_0$ 描述死亡年齡分布的平均，$\edag$
描述其變異，兩者互補。本文採用簡略生命表（Chiang）的離散近似計算。


\subsection{兩條估計路線：高斯與卜瓦松}\label{sec:routes}

LC 模型的估計自始即與 SVD 綁在一起，而此一選擇有其歷史脈絡。雙線性的
對數乘積結構在 LC 之前即已存在於列聯表分析，Goodman（1979）的
對數乘積模型（Log-multiplicative Model）以交替迭代求解行列效應的乘積項，
與 LC 的結構相同；Lee and Carter（1992）的貢獻在於把它用於死亡率的時間
外推，並以 SVD 取得封閉解。SVD 的吸引力在於解必定存在、計算瞬間完成、
且第一主成分所解釋的變異比例可直接作為模型適切性的指標，這在 1990 年代的
計算環境下具決定性。代價則是它等價於無加權最小平方，因而隱含同質變異。

Lee and Carter（1992）本身並未主張誤差同質，而是在估計之後以第二階段調整
補救：重新求解時間參數，使配適的總死亡數等於觀測值。Basellini et al.
（2023）指出，此後三十年間高斯路線的主要進展，幾乎都是在改變第二階段
「對齊什麼」。Lee and Miller（2001）發現對齊總死亡數不足以修正零歲平均
餘命的偏誤，改為直接對齊該量；然而以單一總量為約束仍嫌寬鬆，
Booth et al.（2002）因而改以卜瓦松離差逐年齡對齊死亡數；循此思路，
Rabbi and Mazzuco（2021）進一步主張平均餘命只描述死亡年齡分布的位置，
改為對齊描述其離散程度的壽命差異。這是一條逐步精緻化的線，每一步都在
修正前一步「對齊的量不夠貼近真正關心的量」。

然而這條線在準確度上已無太多可爭取的空間。Booth et al.（2005）系統比較
各種替代調整後發現，它們對預測準確度的影響可以忽略，LC 模型準確度較差的
主因實為初始值誤差問題與配適期長度。換言之，問題並不在於第二階段對齊了
哪個量，而在第一階段的估計本身——若第一階段的偏誤是結構性的，事後調整
單一的時間參數並無法移除它。

卜瓦松路線因此從源頭更換架構。Brouhns et al.（2002）將死亡數視為卜瓦松
分配的觀測值、以暴露數為抵補項（Offset），以最大概似法估計參數。此一
改寫有三項後果：誤差直接作用於計數上，不需先取對數；各格的權重由
$\mu_{x,t}$ 內建於資訊矩陣，不必另行指定；且配適總死亡數依照模型假設即
等於觀測值，第二階段調整不再必要。Currie（2013）進一步將其納入廣義線性
模型（Generalized Linear Model）架構，使懲罰項與平滑得以在同一套語言下
處理；Currie（2011）則專門處理高齡資料品質不佳時的推估。就應用而言，
卜瓦松路線已成為精算界的主流，Human Mortality Database 與多數國家的官方
推估亦採此架構。

惟文獻在比較兩條路線時多半止於「卜瓦松較佳」，未說明較佳之處何在。這是
因為兩條路線的差異其實綁了三件事：是否對資料取對數，高斯路線取對數因而須
處理 $\log 0$ 與 Jensen 不等式造成的二階偏誤，卜瓦松路線則不取；是否加權，
SVD 是無加權最小平方，而卜瓦松概似的權重內建於資訊矩陣之中；以及分配假設
本身，常態對卜瓦松。三者同時更換，就無從判斷改善由何者貢獻，也無從判斷在
什麼情況下不必更換——例如在國家級人口上，第二項差異是否仍然重要。本文
第\ref{sec:decomp}的階梯式分解即為拆開這三項而設計。

兩條路線各自的代價其實都出自同一件事——對誤差分布所作的假設。
這也意味著存在第三種可能，即不指定誤差分布而以非參數方法估計之；
第\ref{sec:notapply}將說明此一取徑為何在本文的資料量下不可行。

\subsection{小人口的失效機制與既有對策}

上一小節所述的困難在國家級人口上大多可以忍受，到了小人口才成為主要問題，
原因是三個彼此加乘的機制。第一是零死亡格，當某年齡組的期望死亡數低於 1 時
觀測值經常為 0，而 $\log 0$ 無定義，實務上多以固定值（如 0.5）替代，此舉會
系統性地墊高該格的死亡率；關鍵在於此一機制只作用於零格所在的年齡，因而在
資料上留下明確的位置訊息。第二是對數轉換的加權失衡，Basellini et al.
（2023）指出取對數等於給低死亡率較大權重、給高死亡率較小權重，而高死亡率
出現在高齡、該處人口又相對龐大，因此以死亡數計的誤差被放大；Wilmoth
（1993）因而建議改以死亡數為權重的加權最小平方。這一點同時解釋了上一小節
的現象——高斯路線之所以需要第二階段調整，正是因為第一階段的權重配置與
資訊量不符。第三是過度離散（Overdispersion），死亡計數普遍較卜瓦松分配更為
分散，高齡尤然；Renshaw and Haberman（2003）提出過度離散卜瓦松擬概似，
Li et al.（2009）以 gamma 脆弱性導出年齡別過度離散參數並顯示預測區間較
卜瓦松 LC 增寬 18\% 至 180\%，Djeundje and Currie（2010）另在平滑架構下
直接處理離散過度的計數資料。

面對這些機制，最直觀的對策是增加樣本數。此一取向在應用上已累積相當經驗：
Li and Lee（2005）的共同 LC 模型以一組死亡改善趨勢相近的母體為參考，原為
處理國家層級的聯合推估；Wang et al.（2018）建議合併目標母體過去 10 至
20 年的歷史資料，對象即為臺灣縣市；Jarner and Kryger（2011）的 SAINT 模型
以大母體為錨定，應用於丹麥的年金評價；Wiśniowski et al.（2015）與
Cairns et al.（2011）則以貝氏方法達成類似效果。這些作法形式各異，卻共用
同一個關鍵假設，即參考母體與目標母體的死亡型態相近。該假設在臺灣並不安全，
縣市間平均餘命的最大差距即超過 10 歲（Yue et al. 2019），鄰近地區的社經與
人口組成差異甚大；本文因此不採此一路線，亦不以參考母體為研究方向。

假設不可靠時，另一條路是不借外部資訊而直接平滑目標母體自身的資料。
Yue et al.（2019）將修勻（Graduation）與死亡率模型結合，先以部分標準化
死亡比（Lee 2003）或 Whittaker 比值修勻小區域的死亡率，再套用 LC 或
Li-Lee 模型，並以電腦模擬比較七種死亡率情境，發現型態相近時部分標準化
死亡比表現最佳、型態差異較大時則以 Whittaker 比值較為穩健。修勻的對象亦可
是參數而非死亡率本身，其技術基礎為 Eilers and Marx（1996）的懲罰式 B 樣條，
經 Currie et al.（2004）推廣至年齡與時間的二維平滑；Basellini et al.（2023）
將此類作法分為估計前、估計中、估計後三類，Hyndman and Ullah（2007）與
Rabbi and Mazzuco（2021）屬估計前修勻率，Delwarde et al.（2007）在配適
演算法中加入懲罰年齡參數不規則性的項、Currie（2013）擴充至同時懲罰兩組
年齡參數，二者屬估計中修勻參數，Renshaw and Haberman（2003）以三次 B 樣條
修勻已估出的年齡參數，則屬估計後修勻參數。

修勻雖不需外部母體，卻引入了另一個問題：平滑的強度須由資料決定，而選擇的
準則未必瞄準真正關心的量。Delwarde et al.（2007）以交叉驗證選取平滑參數，
衡量的是曲面對觀測值的預測誤差；Rabbi and Mazzuco（2021）則直接以對觀測值
的誤差比較各種修勻法。兩者瞄準的都是配適誤差而非參數的估計誤差，本文
第肆節將顯示後者在臺灣資料上甚至會傾向選擇「不修勻」，因為修勻依定義必然
增加對觀測值的誤差。

\subsection{本研究的目的}\label{sec:goal}

前述對策的共同特徵，是在承認估計有誤差之後設法使它變小，而未追問誤差的
性質。小樣本問題並非無人研究：Chen et al.（2017）以參數式拔靴法系統地
探討人口規模對最大概似估計量有限樣本分布的影響，涵蓋參數之間的相依結構
與概似比檢定的檢定力；國內的 Wang et al.（2018）與 Yue et al.（2019, 2021）
則以 $t$-ratio 為判準，結論為人口不超過 50{,}000 時 LC 模型須謹慎使用，
保守而言建議至少一百萬。惟這些研究回答的是「估計有多不準」，而非「估計
是否系統性偏移、偏移來自哪一個環節」。更值得注意的是，Chen et al.（2017）
的對象是 M7 模型、估計方式自始即為卜瓦松最大概似，而本文第伍節將顯示
卜瓦松最大概似估計量的漂移項幾乎無偏誤，真正產生系統性偏誤的是 LC 模型
預設的奇異值分解；以卜瓦松最大概似為起點的研究設計，恰好看不到本文關切
的現象。

本研究因此把問題換一個方向提出：與其追求更小的誤差，不如先辨識誤差的
來源，再依來源選用對應的估計量。若三個參數的誤差來自不同的環節，則補救
方式亦應不同，而混用單一方法必然在某些參數上失效。更進一步，若這些環節
在原始資料中留有徵兆，則判斷所需的量便可在配適任何模型之前計算——這對
實務尤其重要，因為取得鄉鎮市區資料的使用者既無真值亦無從進行模擬，一個
可在自己資料上計算的指標，遠比一個固定的人口數門檻有用。

依此方向，本文對既有文獻有四項議題想要暸解。議題一是小樣本研究多針對
變異數，對奇異值分解本身的偏誤著墨較少，仍有補充空間的是該偏誤的來源
及其對預測區間的影響。議題二是修勻壓低預測區間寬度的現象已被
Basellini et al.（2023）在討論節明確提出並點名五齡組，卻未提供任何量化
證據。議題三即前一小節所述的選參準則錯置。議題四是時間參數在小人口的
行為尚未被檢驗，此點由 Yue et al.（2019）於結論中明白指出，該研究僅設定
年齡參數的各種情境而未考慮時間參數，然而漂移項直接乘上預測期長度進入
平均餘命，其偏誤對長期預測的影響遠大於年齡參數。

\subsection{參數的可估性與補救}

承上，若要依誤差的來源選用估計量，須依序回答五個問題：在給定的資料量下
參數能否估計、偏誤該如何處理、變異數該如何處理、損失函數是否也該更換、
以及如何判斷手上的資料適用哪一種。本小節依此順序回顧各條脈絡的理論發展
與應用場景。

\subsubsection{（一）可估性的界限：尖峰隨機矩陣}\label{sec:bbp}

前述文獻處理的都是如何估得更好，但更前置的問題是能不能估。此一問題的
理論工具來自高維統計，其發展可追溯至 Johnstone（2001）提出的尖峰共變異
模型（Spiked Covariance Model）：在變數個數與樣本數同階成長的漸近架構下，
純噪音樣本共變異矩陣的最大特徵值收斂至 Tracy--Widom 分布，因而可據以判斷
觀測到的最大特徵值是否確實含有訊號。Baik et al.（2005）隨後證明此一判斷
存在相變：當訊號強度超過某個門檻時最大特徵值脫離噪音的塊體、否則被吸收其中。
Paul（2007）補上特徵向量的漸近角度，Benaych-Georges and Nadakuditi（2012）
則把結果推廣至矩形矩陣，這正是 LC 模型所需的形式——設 $n\times m$ 的噪音
矩陣加上秩一訊號 $\theta uv^\top$、$c=n/m$，該文第 3.1 節式 (9) 與式 (10)
指出，當 $\theta>c^{1/4}$ 時估計出的奇異向量與真值有正的漸近內積且該內積
有閉式，當 $\theta\le c^{1/4}$ 時內積漸近為零。

這個結果的性質與前述文獻截然不同：門檻以下不是估得比較差，而是估出來的
方向與隨機向量無法區分，此時任何估計量的改良都無從發揮，正確的作法是
不予估計。就應用而言，此一架構已廣泛用於訊號處理中判斷陣列訊號的個數、
基因體學中決定主成分的保留數目，以及財務計量中的因子個數檢定；死亡率
模型則尚未見到援用，本文為其一。

相變之所以發生，可由 Zhang et al.（2022）的 HeteroPCA 看得更清楚。該文
指出當各列的噪音變異數不同時，樣本 Gram 矩陣的期望為
$\Sigma_0+\mathrm{diag}(\sigma_1^2,\dots,\sigma_p^2)$，亦即對角線被噪音
變異數系統性灌水而非對角線不受影響，故異質變異越嚴重、主特徵向量偏得
越多，補救因而須作用於奇異向量。Gavish and Donoho（2017）給出的最適
奇異值收縮則提供了一個有用的反例：它在矩陣重建的意義下最適，卻對
$\beta_x$ 完全無作用，因為它作用在奇異值上而非奇異向量——最適化矩陣重建
與最適化參數還原並非同一件事。

須註明本文援用此一架構的限制。上述理論皆為漸近結果，成立於維度與樣本數
同時趨於無窮之下，而本文的設定為 $A=22$、$T=24$，屬相當小的維度；本文
第肆節與第伍節的實測顯示門檻的落點與門檻以上的還原程度都與理論吻合，
但這是經驗上的驗證，而非理論保證。

\subsubsection{（二）偏誤減少：在估計當下無偏}

偏誤主導的參數無法靠加寬區間或增加重複處理，必須在估計的環節上動手。
此一脈絡的起點是最大概似估計量偏誤的漸近展開：Cox and Snell（1968）給出
$O(1/n)$ 偏誤項的一般形式，使得「先求最大概似解、再減去估計出的偏誤」
成為可能，此即事後修正（Corrective）作法。Firth（1993）指出另一種作法
更為根本——與其估完再修，不如直接修改得分方程，使其解自始即無 $O(1/n)$
偏誤，此即預先修正（Preventive）作法；對正則連結的廣義線性模型，該修正
等價於對概似乘上 Jeffreys 事前分配 $|I(\theta)|^{1/2}$。

兩種作法的差別在本文的情境下具決定性。事後修正需要最大概似解先行存在，
而本文所面對的正是它不存在的情形：Albert and Anderson（1984）把 logistic
迴歸的資料型態分為完全分離、準完全分離與重疊三類，並證明前兩者下最大
概似估計量不存在；卜瓦松對數線性模型在零計數時有同構的問題，概似在
$\alpha_x\to-\infty$ 時仍遞增。Firth 的預先修正則使概似在參數空間邊界
不再遞增，Kosmidis and Firth（2021）進一步證明其解在二項反應的廣義線性
模型中必定有限。

就應用而言，此一方法在罕見事件的 logistic 迴歸上已成標準：Heinze and
Schemper（2002）把它引入流行病學與臨床研究，處理某一組別全無事件時的
估計；生態學的物種分布模型、以及樣本數少而事件更少的臨床試驗次群組分析，
亦屬常見場景。這些場景的共同特徵是「某些格子的資訊量趨近於零」，與本文
小人口幼年組的情形結構相同。惟死亡率模型上尚未見到系統性的應用，原因
可能在於國家級資料不會出現零格，而小區域資料的研究者多採取增加樣本數的
路線。

\subsubsection{（三）收縮估計：以偏誤換變異數}

變異數主導的參數則適用收縮，其核心思想是接受少量偏誤以換取變異數大幅
下降，使總均方誤差下降。Stein（1956）首先指出多維常態平均數的樣本平均
在三維以上並非可容許（Admissible），James and Stein（1961）給出具體的
支配估計量，Efron and Morris（1973）則以經驗貝氏重新詮釋之，使收縮強度
有了「由資料估計先驗變異」的操作意義。這一層詮釋對本文尤其重要，因為它
說明收縮量不必以調校參數指定。

同一思想在迴歸上的實作則另成一系。Hoerl and Kennard（1970）的 Ridge 以
$L_2$ 懲罰縮減係數，原為處理多重共線性；Tibshirani（1996）的 LASSO 改用
$L_1$ 懲罰，因其近端算子為軟閾值而能把係數精確壓至收縮目標，同時達成
收縮與變數選擇；惟 LASSO 在變數高度相關時傾向任意只保留其一，
Zou and Hastie（2005）的 Elastic Net 因而混合兩種懲罰；LASSO 對大係數
亦施加同等收縮而產生偏誤，Zou（2006）的自適應 LASSO、Meinshausen（2007）
的鬆弛 LASSO 與 Zhang（2010）的 MCP 相繼提出修正。演算法方面，
Friedman et al.（2010）的座標下降法已成為標準作法，Tseng（2001）與
Lee et al.（2014）分別提供區塊座標下降與近端牛頓法的收斂理論。

此一脈絡進入死亡率模型的時間不算晚，但懲罰的對象與本文不同：
Delwarde et al.（2007）與 Currie（2013）懲罰的皆為年齡參數的差分，即
粗糙度，目的在取得平滑的年齡曲線；本文則依指導方向懲罰年齡參數本身，
目的在辨識哪些年齡組顯著偏離參考水準。兩者的解釋並不相同——前者回答
「年齡曲線應該多平滑」，後者回答「哪些年齡的改善速度顯著偏離平均」。

更重要的是，這一整條脈絡有一個共同的操作負擔：懲罰強度必須由資料選出，
而前一小節所述的選參準則錯置問題於此同樣無法迴避。此處值得引入一條字面上
與本文問題完全吻合、卻尚未被死亡率文獻援用的脈絡——小區域估計。
Fay and Herriot（1979）的模型正是為了估計美國人口普查中小地方的所得而提出，
其設定與本文高度平行：各小區域有一個直接估計值與一個已知的抽樣變異數，
真值則圍繞某個迴歸預測值分布；以受限最大概似（Restricted Maximum
Likelihood，以下簡稱 REML）估計先驗變異成分後，收縮強度即為訊噪比的函數，
不需交叉驗證或格點搜尋。Rao and Molina（2015）為該領域的標準參考。就應用
而言，該架構已是各國統計局編製小地區指標的主流工具，涵蓋所得、失業率、
貧窮率等，處理的正是「每個小區域的樣本都太小、但區域之間可以互相借力」
這一類問題——與鄉鎮市區生命表的結構完全相同。

\subsubsection{（四）穩健損失與分位數迴歸}

前三條脈絡都在「平方損失或概似」的架構內改良，第四條則更換損失函數本身。
Huber（1964）提出 $M$-估計量與 $\varepsilon$-污染模型，其設定是絕大多數
觀測來自模型、少數來自別處，目標是使少數觀測無法主導估計；Hampel（1974）
以影響函數刻畫單一觀測的影響力，Donoho and Huber（1983）則以崩潰點
（Breakdown Point）刻畫估計量所能承受的污染比例上限。Rousseeuw（1984）的
最小中位數平方與最小截尾平方把崩潰點推到 50\%，代價是效率極低；
Müller and Neykov（2003）把截尾概似的崩潰點理論推廣至廣義線性模型。

與 $M$-估計平行發展的是分位數迴歸。Koenker and Bassett（1978）以檢查函數
$\rho_\tau$ 取代平方損失，使迴歸得以刻畫條件分位數而非條件平均；
Koenker（2005）為該領域的標準參考，Koenker（2017）回顧其四十年的發展。
此後的擴充包括 Zou and Yuan（2008）的複合分位數迴歸——讓不同分位共用
斜率而只讓截距平移，以降低單一分位估計的變異——以及
Machado and Santos Silva（2005）針對計數資料的抖動處理。
Newey and Powell（1987）的 expectile 則以非對稱平方損失取代檢查函數，
保留可微性因而可用迭代加權最小平方求解。

以上兩條脈絡看似平行，實則可併為一族，而這一點對本文的設計具決定性。
Breckling and Chambers（1988）指出，若把 Huber 的損失乘上一個依殘差正負
而異的非對稱權重，即得一個由兩個參數索引的損失函數
\begin{equation}\label{eq:mq}
\rho_{\tau,c}(u)=2\,w_\tau(u)\,\psi_c(u),\qquad
w_\tau(u)=\begin{cases}\tau & u>0\\ 1-\tau & u\le 0\end{cases},\qquad
\psi_c(u)=\begin{cases}u^2/2 & |u|\le c\\ c|u|-c^2/2 & |u|>c,\end{cases}
\end{equation}
其中 $\tau$ 控制非對稱、$c$ 控制穩健程度。此即 M-分位數（M-quantile）。
其兩個極限恰好是前述兩條脈絡的端點：$c\to0$ 時
$\rho_{\tau,c}$ 正比於 $w_\tau(u)|u|$，即 Koenker and Bassett（1978）的
檢查函數；$c\to\infty$ 時 $\rho_{\tau,c}=2w_\tau(u)u^2/2$，即
Newey and Powell（1987）的非對稱最小平方，而 $\tau=0.5$ 時更退化為
一般最小平方。換言之，\textbf{「要不要穩健」這個二選一的問題，在
式~\eqref{eq:mq} 中變成「要多穩健」這個連續刻度上的定位問題}。
Catania and Luati（2020）以積分 Hogg 函數構造另一個從中位數連續過渡到
平均數的單參數族，並指出把調參設在穩健尺度估計的一個低比例上即可兼顧
穩健與效率，說明這種以刻度取代開關的想法具有一般性。

M-分位數之所以值得注意，還因為它的主要應用場域正是本小節第（三）目
所述的小區域估計。Chambers and Tzavidis（2006）以 M-分位數取代
Fay--Herriot 的隨機效果架構，理由是 M-分位數可直接給出各區域自己的迴歸
係數而不必假設區域效果的分布；Bianchi et al.（2018）補上了估計與檢定的
理論。因此收縮與穩健損失這兩條脈絡在小區域估計文獻中已經交會過一次，
只是尚未被帶進死亡率模型。

損失函數的對稱性則是另一個需要檢查的前提。Huber 的 $\psi_c$ 對兩側施加
相同的截點，其隱含的假設是污染來自兩側。Xu and Chen（2018）因而在
Huber 損失中引入一個非對稱度量，使兩側可以帶不同的權重，並給出該度量的
估計方式。此一修改與本文直接相關：如後文第\ref{sec:mq}所示，小人口的
扭曲並非兩側對稱，而是在不同年齡朝相反方向偏移。

計算面則決定了哪些損失函數在雙線性模型中實際可用。檢查函數不可微，
Koenker and Bassett（1978）以線性規劃求解，Koenker and Park（1996）進一步
處理「對參數為非線性」的情形並提出內點法——這正是 LC 模型所屬的情形，
因為 $\alpha_x+\beta_x\kappa_t$ 對 $(\beta,\kappa)$ 是雙線性而非線性。
另一條路是 Hunter and Lange（2000）的 MM 演算法：以 $|u|+\epsilon$ 取代
$|u|$ 建立一個處處可微的二次優化函數，每一步化為加權最小平方，目標函數
單調不增，$\epsilon$ 隨迭代縮小即收斂至原問題。本文採後者，理由有二：
雙線性結構使兩個區塊各自都有加權最小平方的封閉解，無須呼叫線性規劃求解器；
而同一段程式碼在 $c$ 取不同值時即涵蓋式~\eqref{eq:mq} 的整個族，
包含高斯極限，使各變體之間的比較不受演算法差異污染。

死亡率模型上的援用並不多。Hyndman and Ullah（2007）在函數型主成分架構下
提出穩健版本，Santolino（2020）則直接以分位數迴歸估計 LC 模型的參數，
逐一估計各分位的 $\beta^\tau$ 與 $\kappa^\tau$。須指出 Santolino（2020）
自承參數隨 $\tau$ 並無明顯規律，此點與每個年齡--年份格僅有一個觀測、
識別完全依賴雙線性結構跨格借力有關，本文第伍節之五將回到這個問題；
而本文第\ref{sec:mq}將指出另一層原因——$\tau$ 的最適值在不同年齡朝
相反方向移動，故單一 $\tau$ 的參數本就不應呈現規律。

此處須預先點出一項與本文情境直接相關的區分。古典穩健統計所設想的是
「污染」，即少數觀測來自模型之外；但零死亡格的性質並非如此——一格觀測到
零並不表示該格被弄錯，而是它只告訴我們死亡數為零，不足以定出死亡率。
依 Heitjan and Rubin（1991）的用語，這屬於粗化資料（Coarse Data）而非
污染資料，兩者所需的處置方向相反。本文第柒節之四將以實測結果說明此一
區分的後果。

\subsubsection{（五）判斷準則}

上述四者各有適用的情境，因此還需要一組準則，用以判斷手上的資料落在哪一
情境。既有研究給出的準則多為人口數門檻，如 Wang et al.（2018）的
50{,}000 與一百萬，惟門檻由模擬得出，未必適用於年齡結構或配適期長度不同
的資料。本文第參節因此把前述四個層次各自翻譯成一個可由資料直接計算的量：
零死亡格的分布對應偏誤減少的需求，死亡數的跨年齡落差與噪音變異的跨度
對應加權的需求，最大奇異值與純噪音上緣的比值對應可估性的界限，而零格
比例是否超過 50\% 則對應穩健損失的作用範圍。四者皆不需要真值，前三者
甚至不需要配適模型。

\subsubsection{（六）前提不符的三條取徑}\label{sec:notapply}

上述四條脈絡是本文選用方法的依據，但文獻中還有三條看來更直接的取徑，
本文未予採用。三者被排除的理由並非效果不佳，而是其成立所需的前提在小
人口的資料結構下不成立。把理由寫清楚有兩個用處：一是說明本文的方法選擇
不是任意的，二是這三條取徑之所以不成立，恰好各自標示出小人口問題的一項
特徵，因而本身即是值得深入的起點。

第一條是對「離群年份」穩健化。Hyndman and Ullah（2007）是死亡率文獻中
最直接的穩健化嘗試，其作法分兩步：先以懲罰式樣條把每一年的
$\log m_{x,t}$ 在年齡方向修勻，再以穩健主成分分析取代一般的 SVD，
依每一年整條曲線的積分平方誤差判斷該年是否離群並降低其權重。此一設計
所設想的離群值是戰爭、流感大流行或地震等使\textbf{整個年度}偏離趨勢的
事件，該文所舉的例子即法國的 1914 至 1918 年與 1943 至 1945 年。
其隱含前提因此是「少數異常年份對多數正常年份」。但小人口的噪音來自
每一格各自獨立的卜瓦松抽樣，每一年同樣嘈雜，並不存在可供對比的多數，
故逐年降權無從發揮。須強調這與本文第伍節之五所檢驗的逐格截尾並不相同：
截尾的單位是格，此處的單位是年，若小人口的噪音恰在某些年份同向累積，
逐年穩健化仍可能有效，因此值得單獨檢驗而非由前者類推。本文第\ref{sec:hu}
即以實測給出結論。同一取徑的更進一步版本是 Hörmann et al.（2015）的
動態函數型主成分，它不只用同期的共變異數，還用各階落後的共變異數以構造
最適濾波器；就 LC 而言這是有吸引力的，因為 SVD 把各年份視為可交換，
而 $\kappa_t$ 實則為隨機漫步。但動態主成分需要長期共變異數算子可估，
本文 $T=24$ 且噪音變異與訊號同一量級，落後共變異數的估計量本身即不可靠。

第二條是繞過高斯與卜瓦松、改走半參數的第三條路線。死亡率在本質上是
乘性的——$m_{x,t}=\exp(\alpha_x+\beta_x\kappa_t)\cdot\varepsilon_{x,t}$
——而第\ref{sec:routes}已說明，高斯路線的代價來自取對數、
卜瓦松路線的代價來自等離散假設。Zhang, Lin, and Yang（2022）針對乘性線性迴歸
提出以核密度估計誤差分布的最大非參數概似估計，其模擬顯示在誤差分布為
多峰、非對稱或重尾時優於最小平方與最小乘積相對誤差估計。就前提而言
這正對應本文的困難，因為取對數後的誤差分布恰是非對稱的。但該方法要求
誤差密度可估，而本文每個年齡只有 24 個觀測；更根本的是，在零格集中的
幼年組，誤差分布在零處有一個原子而非連續密度，核密度估計並非為此設計。
此一取徑真正的價值因而在於提醒：兩條路線的差異可以歸結為對誤差分布的
假設，而該假設本身是可以放鬆的——只是放鬆的代價需要資料量來支付，
而資料量正是本問題所缺的。

第三條是改由世代而非年期的角度建模。Basellini and Camarda（2025）指出
以年期為軸的模型難以刻畫真實出生世代的經驗，因而提出把人口學的約束嵌入
修勻架構的世代約束 P-樣條，其樣本外驗證優於世代版的 LC 變體與由
age--period 預測反推 Lexis 對角線的慣行作法。就死亡率預測的一般問題而言
這是重要的推進，但移到小人口時方向相反：Lexis 圖上的一條世代對角線所
涵蓋的死亡數比一個年齡--年份格更少，稀疏問題因而加重而非減輕。
不過該文「把人口學的約束寫進估計而非事後檢查」這一想法與本文第肆節的
結論一致，值得保留為後續的方向。

\subsection{既有應用場景與本文前提的對照}\label{sec:gap}

前一小節按「能不能估、偏誤怎麼處理、變異數怎麼處理、損失函數要不要換、
如何判斷」的順序回顧了五條脈絡。但把方法列出來並不等於說明它們為何可用，
因為每一條脈絡的方法都是在特定的資料結構上發展出來的，其效果的保證也
只在那個結構的前提下成立。本小節因此把各方法的\textbf{原始應用場景}與
\textbf{本文情境的前提落差}並列，作為本文方法選擇的依據，也作為後續各節
「假設檢驗」的檢驗清單——表~\ref{tab:gap} 右欄的每一項落差，在第伍節與
第陸節都有一個對應的實證檢驗。

\begin{table}[H]
\centering
\caption{各方法的原始應用場景與本文情境的前提落差}
\footnotesize
\renewcommand{\arraystretch}{1.25}
\begin{tabular}{p{2.5cm}p{3.6cm}p{3.7cm}p{4.3cm}}
\toprule
方法（文獻） & 原始應用場景 & 該場景成立的關鍵前提 & 本文情境的落差\\
\midrule
\multicolumn{4}{l}{\textit{（一）可估性}}\\
尖峰隨機矩陣、BBP 相變（Johnstone 2001；Baik et al. 2005） &
陣列訊號處理的訊源個數、財務的因子個數、基因體學的主成分保留數 &
維度與樣本數同時趨於無窮且比值收斂，噪音同變異數 &
本文 $A=22$、$T=24$ 皆為固定小維度，門檻只能當作量級的參考而非臨界值\\
HeteroPCA（Zhang et al. 2022） &
單細胞定序等各列噪音變異數差異極大的高維矩陣去噪 &
行列維度皆大，噪音跨列異質但格內獨立 &
維度小，且噪音變異數與死亡數成反比、跨年齡跨越三個數量級，異質程度遠大於原設定\\
\midrule
\multicolumn{4}{l}{\textit{（二）偏誤減少}}\\
Cox and Snell（1968）事後修正 &
指數族模型的 $O(n^{-1})$ 偏誤修正 &
最大概似估計值有限且存在 &
零格導致完全分離，$\hat\theta$ 可為 $-\infty$，修正無從施加\\
Firth（1993）預先修正 &
二元邏吉斯迴歸的完全分離（Albert and Anderson 1984；Clogg et al. 1991）、罕見事件研究（Heinze and Schemper 2002） &
線性預測子，每個參數有足夠的有效觀測數 &
$\alpha_x+\beta_x\kappa_t$ 對參數為雙線性，訊息矩陣隨 $\beta,\kappa$ 而變；每個 $\alpha_x$ 只有 $T=24$ 個觀測\\
\midrule
\multicolumn{4}{l}{\textit{（三）變異數與收縮}}\\
Fay and Herriot（1979） &
美國普查小地方（人口數千以下）的人均所得估計，用於歲收分配 &
直接估計量不偏、其抽樣變異數已知，區域效果為常態 &
本文的直接估計量 $\hat\beta_x$ 本身\textbf{有偏}而非僅有噪音，收縮所交換的不是偏誤換變異數\\
Ridge 與 LASSO（Hoerl and Kennard 1970；Tibshirani 1996） &
共線性與高維變數選擇 &
懲罰強度可由交叉驗證選出，且設計矩陣固定 &
懲罰對象為 $\beta_x$ 而設計矩陣含待估的 $\kappa_t$，交叉驗證需切分格而破壞雙線性結構\\
\midrule
\multicolumn{4}{l}{\textit{（四）穩健損失}}\\
Huber（1964）$M$-估計 &
位置參數的 $\varepsilon$-污染估計；訊號處理與電力系統的離群值抑制（Zoubir et al. 2018） &
污染為少數、對稱、與觀測的資訊量無關 &
零格集中於幼年組且比例可達 95\%，非少數；扭曲為單邊；被扭曲者恰為 $\beta_x$ 最大者\\
LTS 與 LMS（Rousseeuw 1984）；截尾概似（Müller and Neykov 2003） &
迴歸診斷的離群值偵測；廣義線性模型的污染資料 &
崩潰點理論要求逐觀測損失具次緊緻性；被丟棄者為壞資料 &
卜瓦松對數概似在 $z=0$ 時不具次緊緻性；被丟棄者為攜帶趨勢訊號的格\\
分位數迴歸（Koenker and Bassett 1978；Koenker 2005） &
薪資分布、兒童生長參考曲線（Wei et al. 2006）、風險值 &
每個共變數水準上有大量觀測，條件分位數可由資料識別 &
每個年齡--年份格僅一個觀測，識別完全依賴雙線性結構跨格借力\\
M-分位數（Breckling and Chambers 1988；Chambers and Tzavidis 2006；Bianchi et al. 2018） &
英國與義大利的家戶調查，估計地方層級的所得與剝奪指標 &
單位層級的調查資料，每個小區域有多個受訪單位 &
本文僅有聚合後的死亡計數，每格一個觀測；$c=1.345$ 的效率依據為連續對稱誤差\\
\midrule
\multicolumn{4}{l}{\textit{（五）死亡率文獻中的穩健化}}\\
穩健函數型主成分（Hyndman and Ullah 2007） &
法國全國死亡率與澳洲生育率；辨識出的離群年為 1914--1919、1940--1945、1960 &
少數年份因戰爭或疫病而結構斷裂，多數年份乾淨 &
各年同分布、同樣嘈雜，不存在可對比的多數，門檻不被觸發（第\ref{sec:hu}）\\
動態函數型主成分（Hörmann et al. 2015） &
函數型時間序列的最適濾波，利用各階落後共變異數 &
長期共變異數算子可估，序列長度足夠 &
$T=24$ 且噪音與訊號同量級，落後共變異數的估計量本身不可靠\\
LC 分位數模型（Santolino 2020） &
西班牙全國死亡率，以內點法解非線性分位數迴歸 &
全國層級資料，計數大、無零格 &
小人口有零格，且本文發現最適 $\tau$ 隨年齡變號（第\ref{sec:mq}）\\
\bottomrule
\end{tabular}
\label{tab:gap}
\end{table}

表~\ref{tab:gap} 的排列方式本身即說明了本文的策略。落差可分為兩類。
第一類是\textbf{量的落差}，即前提的方向正確而程度不足，例如 BBP 相變的
漸近性與本文的小維度、或 Firth 的觀測數要求與本文 $T=24$ 的落差。
這一類方法仍可使用，但其理論保證退化為量級的參考，必須以模擬確認，
本文第\ref{sec:est}即為此而設。第二類是\textbf{質的落差}，即前提的方向
本身相反，例如 Fay--Herriot 假設直接估計量不偏而本文的直接估計量有偏、
或古典穩健統計假設污染為少數而本文的零格為多數。這一類方法不能以
「加大樣本」或「調整參數」補救，必須改換問題的設定，這也是本文
第\ref{sec:decomp}先做偏誤來源分解、再依來源選方法的理由。

值得注意的是，第二類落差有一個共同的形式：既有方法都假設\textbf{有一個
可信的多數}——多數觀測不偏、多數觀測乾淨、多數年份正常——而小人口所
缺的恰恰是這個多數。這個觀察串起了本文所有的負面結果（第\ref{sec:hu}與
第伍節之五），也解釋了為何有效的三種作法（卜瓦松概似、Firth 懲罰、
資訊加權）都不依賴多數／少數的對比，而是直接改寫每一格的貢獻方式。
本文第\ref{sec:theory}將回到這一點作系統的整理。

\subsection{理論架構：懲罰的三個位置}\label{sec:theory}

前兩小節把方法與其應用前提並列，但仍是一份清單。本小節把它收成一個架構，
作為全文方法選擇的依據，也作為後續各節「假設檢驗」的對象。

\subsubsection{（一）把年齡維度視為小區域}\label{sec:sae}

LC 模型的結構本身就提示了一個歸類。把式~\eqref{eq:lc} 的下標讀一遍：
每一個年齡 $x$ 有自己的水準 $\alpha_x$、自己的斜率 $\beta_x$、
以及自己的噪音變異數 $\sigma_x^2$（第\ref{sec:eda}的診斷\cn{3}已量出
後者跨年齡相差三個數量級），而 22 個年齡靠一個共同的 $\kappa_t$ 互相借力。
這正是小區域估計（Small Area Estimation）中巢式誤差迴歸
（Nested Error Regression）的結構：若把年齡視為「小區域」，則每個區域
只有 $T=24$ 個觀測，樣本量不足以單獨估計該區域的參數，必須借用區域之間的
共同結構。

這個歸類不是修辭。小區域估計文獻在過去四十年裡處理的正是「每個區域的
樣本都太小、但區域之間可以互相借力」這一類問題，其發展脈絡與本文所需
高度重合。Fay and Herriot（1979）以區域層級的經驗貝氏收縮起頭，
Rao and Molina（2015）為該領域的標準專著；
Chambers and Tzavidis（2006）以 M-分位數取代隨機效果的分布假設，
使各區域可以有自己的迴歸係數；而 Lahiri and Salvati（2023）把這一脈
推到最貼近本文的形式——其模型\textbf{同時}容許區域別的迴歸係數與
區域別的變異成分，以穩健的區域別估計方程求解，並以參數式拔靴法與
Jackknife 估計均方誤差。本文目前以三件不同的工具處理這三件事
（Fay--Herriot 收縮處理 $\beta_x$ 的變異、HeteroPCA 處理 $\sigma_x^2$
的異質、M-分位數處理損失的形狀），而該文是在一個模型裡處理。

須立刻指出這個類比的斷裂處，因為那既是本文的新意，也是本文最大的弱點。
小區域估計的模型中，共變數是\textbf{已觀測}的；而 LC 模型的
$\kappa_t$ 是\textbf{待估的潛在變數}，必須與 $(\alpha_x,\beta_x)$ 聯合
求解。Lahiri and Salvati（2023）的漸近理論建立在共變數已知之上，
因此無法直接套用。把小區域估計的工具移到雙線性低秩模型上，等於要求
「借力的基準」本身也由同一批資料估出，這一層在該文獻中尚無對應的處理。
本文因此只援用其\textbf{問題的歸類}與\textbf{方法的分類}，不援用其
漸近結果，並把此點列為研究限制。

另須與既有的死亡率文獻劃清一件事。Wang et al.（2018）已把小區域估計
列為關鍵詞，但該文所稱的小區域是\textbf{地理區域}，借力的方向是由鄉鎮
向縣市與全國；本文所稱的小區域是\textbf{年齡組}，借力的方向是在同一個
人口內部由年齡向年齡。兩者是同一個統計架構在不同維度上的應用，可以並存，
而本文只處理後者。

\subsubsection{（二）懲罰的三個位置}

在上述歸類之下，第\ref{sec:est}所比較的各估計量其實是同一個數學物件的
不同版本：在目標函數上加一項，使估計量不再是最大概似或最小平方。差別只在
\textbf{加在哪裡}。表~\ref{tab:penpos} 把三個位置與其代價並列。

\begin{table}[H]
\centering
\caption{懲罰的三個位置：各自買到什麼、付出什麼}
\small
\renewcommand{\arraystretch}{1.2}
\begin{tabular}{p{2.1cm}p{3.0cm}p{3.1cm}p{2.9cm}p{1.5cm}}
\toprule
加在哪裡 & 代表方法 & 買到什麼 & 付出什麼 & 需調參\\
\midrule
概似 & Firth（1993）的 Jeffreys 懲罰 &
估計量的存在性，以及 $O(n^{-1})$ 的不偏 &
少量變異數 & 否\\
參數 & Fay and Herriot（1979）的經驗貝氏收縮 &
變異數 & 偏誤（收縮量由變異成分定） & 否\\
參數 & Ridge、LASSO、Elastic Net &
變異數，並可作變數選擇 & 偏誤，且強度須由外部準則指定 & \textbf{是}\\
損失形狀 & M-分位數 $\rho_{\tau,c}$ &
對單邊扭曲的抵抗力 & 效率（$c$ 愈小愈多） & \textbf{是}\\
\midrule
\multicolumn{5}{l}{\textit{以下不是懲罰，而是改變各格的相對重要性}}\\
權重 & 資訊加權、HeteroPCA &
異質變異所造成的主向量旋轉 & 無（若權重正確） & 否\\
\bottomrule
\end{tabular}
\label{tab:penpos}
\end{table}

這個分類有兩個用處。其一，它使「該用哪個方法」變成「偏誤落在哪一層」
這個可由診斷回答的問題：$\alpha_x$ 的偏誤源於概似被零格扭曲，故懲罰
應加在概似上；$\beta_x$ 的困難源於變異數過大，故懲罰應加在參數上；
而異質變異不是偏誤也不是變異數的問題，是權重的問題。
其二，它劃出一條界線：三種懲罰都需要目標函數裡還有可用的訊息。
第\ref{sec:bbp}的 BBP 相變即此界線——訊噪比低於門檻時，主奇異向量
與真值漸近正交，任何加在目標函數上的項都無法把訊息變出來。
\textbf{懲罰能重新分配資訊，不能創造資訊。}

\subsubsection{（三）組合：文獻已經走到哪裡}

上述三個位置並非互斥，而近十年的文獻主要在做的正是組合。

穩健損失與參數懲罰的組合已有教科書級的處理。
Zoubir et al.（2018，第 3 章）把 Lasso 與 Elastic Net 的次梯度條件、
最小絕對離差 Lasso（LAD-Lasso）與 Rank-Lasso 整理在一起，並給出
\textbf{迴歸係數與尺度的聯合懲罰式 $M$-估計}及其演算法；
這對本文有直接的意義，因為 M-分位數的截點 $c$ 必須以殘差的尺度為單位，
而本文目前的實作是每次迭代以 MAD 重估尺度，屬權宜之計而非聯合估計。
在應用端，Ranalli et al.（2023）以 Lasso 與 Elastic Net 懲罰的 M-分位數
迴歸分析交通與氣象對空氣品質的影響，同時達成變數選擇、不同 M-分位數上的
係數比較、以及對離群值的抵抗，並給出解析與拔靴兩種變異數估計。
這一篇說明了本文第\ref{sec:limits}的懲罰式工作與第\ref{sec:qlc}的
穩健損失工作並非兩條無關的線，而是同一族的兩個成分。

穩健程度的刻度則不只一條。本文以 Huber 的截點 $c$ 作為刻度，
其作用是\textbf{截斷}大殘差；Hu et al.（2021）另提出
$L_p$-分位數，以損失的\textbf{冪次} $p$ 為刻度——$p=1$ 為分位數迴歸、
$p=2$ 為 expectile，$1<p<2$ 為兩者之間——並處理高維情形下的懲罰與
變數選擇。兩條刻度在數學上不同，哪一條更適合本文的離散計數資料
是一個尚未回答的問題。

另有兩項組合直接對應本文已知的缺口。其一是相依性：本文逐年齡配適，
等於假設 22 個年齡的殘差獨立，而它們共用 $\kappa_t$，顯然不獨立。
Merlo et al.（2022）把廣義估計方程推廣到帶 Huber 損失的 M-分位數
（generalized M-quantile estimating equations），以工作相關矩陣納入
相依結構而保留穩健損失，正是處理這一層。其二是離散標的：
M-分位數的效率論證建立在連續、對稱的誤差上，而死亡數是計數。
Dawber et al.（2022）提出多項 expectile 迴歸並用於 611 個地方勞動市場區的
就業狀態估計，是這批文獻中唯一處理離散標的者。
至於不確定性的量化，Bugallo et al.（2026）針對 M-分位數小區域模型的
均方誤差估計提出偏誤校正，可與本文第\ref{sec:interval}的區間工作對照。

須誠實指出一件事：上述文獻中有六篇的作者群重疊，屬 Pisa 與 Southampton
的 M-分位數小區域估計學派。引用時應意識到這是一個學派內部的連續發展，
而非來自多方的獨立驗證。

\subsubsection{（四）加權的理論：最適權重不是變異數的倒數}\label{sec:wtheory}

表~\ref{tab:penpos} 把加權列為與懲罰並列的一項，而加權在本文的實測中
效果顯著——第\ref{sec:decomp}的變體 C 僅把等權重改為以死亡數加權，
漂移項偏誤即由 0.305 降至 0.141。但「該用什麼權重」在更換損失函數之後
並非自明，而 Koenker（2005，第 5 章）給出了明確的答案，且該答案與
直覺相反。

對最小平方而言，最適權重是觀測變異數的倒數，此即 Gauss--Markov 定理的
加權版本。但 Koenker（2005）第 5.3 節指出，對檢查函數而言，
\textbf{權重應與該分位處的局部密度成正比，而非與標準差的倒數成正比}。
其定理 5.1 給出加權分位數迴歸估計量
$\check{\boldsymbol\beta}_n(\tau)=\arg\min_b\sum_i f_i(\xi_i)
\rho_\tau(y_i-x_i^\top b)$ 的漸近變異數為 $\tau(1-\tau)D_2^{-1}(\tau)$，
其中 $D_2(\tau)=\lim n^{-1}\sum_i f_i^2(\xi_i)x_ix_i^\top$，
並證明該估計量相對於未加權版本為無條件地更有效率。

此一結果對本文有可立即檢驗的後果。本文變體 G 與 K1 所用的權重是
$w_{x,t}=\hat\mu_{x,t}$，而卜瓦松下
$\mathrm{Var}(\log\hat m_{x,t})\approx1/\mu_{x,t}$，故該權重正是
變異數的倒數，即\textbf{最小平方}的最適權重。但依定理 5.1，
檢查函數所需的是局部密度：$\log(D_{x,t}/E_{x,t})$ 在 $\mu$ 不太小時
近似 $N(\log m_{x,t},1/\mu_{x,t})$，其中位數處的密度為
$\sqrt{\mu_{x,t}/2\pi}$，因而正確的權重應與 $\sqrt{\hat\mu_{x,t}}$
成正比。據此可預測：\textbf{$w\propto\sqrt{\hat\mu}$ 應優於
$w\propto\hat\mu$，但只在 $c\to0$（檢查函數）成立；$c\to\infty$
（最小平方）時 $w\propto\hat\mu$ 仍是對的，而中間的 $c$ 則應有中間的
最適冪次。} 此一預測的檢驗列入第\ref{sec:future}的後續工作。

\subsubsection{（五）調參與否：本文的核心判斷}

表~\ref{tab:penpos} 的最後一欄是本文認為最具實務意義的一欄。
需要調參的成員（Ridge、LASSO、Elastic Net、M-分位數的 $\tau$ 與 $c$）
與不需要調參的成員（Firth、Fay--Herriot、資訊加權）在理論上並無高下，
但在本文的情境下有一項不對稱：\textbf{使用者沒有真值，也無法以交叉驗證
選參}——切分年齡--年份格會破壞雙線性結構所依賴的跨格借力，
而這一點第\ref{sec:limits}將以實測說明。

本文的實證結果因此可以濃縮為一句判斷：

\begin{quote}
在資訊不足的資料上，\textbf{懲罰的位置比懲罰的強度重要}，而不需要調參的
成員一律勝過需要調參的成員。
\end{quote}

這句判斷同時給了 Ridge、LASSO 與 Elastic Net 一個明確的定位。
它們在本文中的結果是負面的，但那不是因為懲罰參數這個想法錯了，
而是因為在本問題上選參所需的外部資訊恰好不存在。把它們與 Firth 及
Fay--Herriot 放在同一張表裡比較，負面結果才成為一項有內容的陳述，
而不是一段離題的嘗試。

\subsubsection{（六）與既有方法論的吻合與不符}

最後把本文的架構放回小樣本偏誤與資料污染兩套方法論之中。吻合之處有三。
第一，BBP 相變所預測的「門檻以下不可還原」在本文的 T-ratio 診斷上確實
出現，$\beta_x$ 在人數二十萬以下時變異數主導。第二，Stein（1956）以降的
收縮理論所預測的「變異數主導時收縮有益」亦成立，且 Fay--Herriot 的
免調參性質使其在實測中即追平以真值挑出的最適懲罰強度。第三，
Firth（1993）所針對的「最大概似估計值不存在」確實是本問題的核心，
該文自述其動機即為此，並引 Albert and Anderson（1984）與
Clogg et al.（1991）為例。

不符之處有四項，表~\ref{tab:gap} 右欄已逐一列出，此處補上其理論陳述。

前兩項為低秩模型所特有。\textbf{其一，異質變異在秩一模型中造成的是偏誤，
而非效率損失}——線性模型中忽略異質變異使最小平方失去效率但仍不偏，
而 Zhang et al.（2022）指出各列噪音變異數不同時樣本 Gram 矩陣的期望為
$\Sigma_0+\mathrm{diag}(\sigma_1^2,\dots,\sigma_p^2)$，對角線被不均勻地
灌水，主向量因而被系統性旋轉；本文的變體 E 即此一差異的純粹展示。
\textbf{其二，被扭曲的觀測與其資訊量正相關，而非獨立}——在雙線性秩一模型
中 $|\beta_x\kappa_t|$ 最大的年齡同時也是殘差最大的年齡，故「最不合模型」
與「最具資訊」高度重疊，表~\ref{tab:retain} 即為直接證據。

後兩項關乎零格的性質。\textbf{其三，零格是粗化而非污染}——一格觀測到零
不表示該格被弄錯，而是它只告訴我們死亡數為零；依
Heitjan and Rubin（1991）的用語這屬於粗化資料（Coarse Data），
而兩者所需的處置方向相反：污染需要降低少數觀測的影響力，粗化則需要保留
正確的概似支撐。這統一解釋了為何 Firth 優於加權中位數——並非因為它更穩健，
而是因為它在計數尺度上估計、根本不產生那筆錯誤資訊。
\textbf{其四，截尾類的崩潰點理論在此前提不成立}——
Müller and Neykov（2003）的理論建立在逐觀測損失於 $z>0$ 時為次緊緻之上，
而 $z=0$ 時 $\gamma_0(\theta)=\exp(\theta)$ 在 $\theta\to-\infty$ 時
趨近於零，不具次緊緻性；零格正是小人口的核心特徵，故截尾類的失效不是
實作困難，是理論前提本身不成立。

把這四項壓縮一次，可得一個更簡短的說法：表~\ref{tab:gap} 右欄的質性落差
有一個共同的形式——既有方法都假設\textbf{存在一個可信的多數}。
Fay and Herriot（1979）假設多數區域的直接估計量不偏，Huber（1964）假設
多數觀測來自模型，Rousseeuw（1984）假設多數觀測可保留，
Hyndman and Ullah（2007）假設多數年份正常。而小人口所缺的恰恰是這個多數。
反過來說，本文三項有效的作法——卜瓦松概似、Firth 懲罰、資訊加權——
沒有一項需要區分多數與少數，它們改寫的是每一格的貢獻方式本身。
\textbf{這是本文對方法選擇最可操作的一條啟示：在資訊不足的資料上，
應優先考慮不依賴多數／少數對比的方法。}

綜合上述，本研究的問題宜歸類為\textbf{資訊不足與異質變異下的低秩小區域
估計}，而非資料污染。古典穩健統計在此提供的是兩件可借用的東西——
崩潰點作為方法適用範圍的上界，以及效率代價的量化方式——而非其核心工具。

\section{研究資料與診斷}

\subsection{資料來源與分析邏輯}

\subsubsection{（一）資料來源與處理}

本文資料取自人類死亡率資料庫（Human Mortality Database，以下簡稱 HMD）的
臺灣序列，涵蓋 1970 至 2024 年的分年齡死亡數與暴露數（人年）。選用 HMD 而
非內政部原始檔的理由有二：其一，HMD 已完成暴露數的人年換算與跨年份的
定義一致性檢查，可避免因人口數與人年混用而產生的系統性偏差；其二，其
資料處理流程有公開文件，便於重現。

原始檔為單齡組，本文依官方簡易生命表的慣例聚合為 0、1 至 4、5 至 9、
依此類推至 95 至 99，並將 100 歲以上併為開放組，共 22 個年齡組。併組的
理由是 110 歲以上於 1970 至 2009 年的暴露數為零，若保留則該列在對數尺度上
無定義。主要分析對象為女性 2001 至 2024 年，共 $22\times24=528$ 格；選擇
女性的理由是其死亡率的時間趨勢較為規則，可使模型設定失誤與小樣本效應
分離；選擇 2001 年起的理由則是避免更早期資料登記制度變動所帶來的干擾，
惟此一區間涵蓋 2020 至 2024 年的新冠肺炎衝擊，其影響列為研究限制。

\subsubsection{（二）分析邏輯}

本文的核心問題是小人口下的估計行為，但真實的鄉鎮市區資料並無真值可供
比對，因此採取「以全國資料建構真值、再模擬小人口」的設計。具體而言，
先以全國女性 2001 至 2024 年的配適結果作為真值
$\{\alpha^{\text{true}}_x,\beta^{\text{true}}_x,\kappa^{\text{true}}_t\}$，
令 $E_{x,t}=N\cdot w_{x,t}$、
$D_{x,t}\sim\text{Poisson}(E_{x,t}m^{\text{true}}_{x,t})$，其中 $w_{x,t}$
為全國最後一年的年齡結構、$N$ 為模擬的單一性別人口數。如此一來，任何
估計值與真值的差距都只能來自抽樣與估計程序本身，而非母體差異。

此一設計有兩項必須言明的代價。其一，它隱含「小區域與全國有相同年齡型態」
的假設，真實鄉鎮市區的型態可能本就不同；本文第參節之二的前三項診斷不受
此假設影響（它們只描述資料本身），但第四項診斷的最大奇異值會同時反映型態
差異與噪音。其二，死亡數以卜瓦松生成，過度離散在設計上即不存在，因此
負二項或準概似方法的效益無從測得。

評估則分三個層次，依序對應本文的三類結果。第一層是\textbf{參數}，
包含 $\alpha_x$ 的逐年齡偏誤、$\beta_x$ 與真值的相關係數與平方誤差和、
以及 $\kappa_t$ 的漂移項偏誤；第二層是\textbf{預測區間}，以零歲平均餘命的
區間寬度與涵蓋率衡量；第三層是\textbf{生命表}，以零歲平均餘命的偏誤與
均方根誤差衡量。三層分開評估的理由是它們未必同向——第伍節之七將顯示，
對數尺度上偏誤最大的年齡並非對平均餘命影響最大的年齡。

模擬的人口規模取 $10^4$ 至 $5\times10^6$，重複次數除特別註明者外一律為
100 次、種子固定，以使各表之間可以逐格對照。少數模擬會發散（單次平方
誤差可達中位數的 $10^6$ 倍），故統計量一律採中位數而非平均數。

\subsection{四項事前診斷}\label{sec:eda}

第\ref{sec:intro-data}已呈現四項診斷在本文資料上的結果（圖~\ref{fig:eda}），
本小節說明其計算方式、判讀門檻，以及各門檻所指向之方法的優點與代價。
四者的共同特點是不需要真值：前三項為純粹的探索性資料分析，第四項僅需
一次成本不到一秒的預備配適。程式見 \texttt{R/core/eda\_diagnostics.R}。

\subsubsection{（一）計算方式}

第一項診斷取各年齡組單年死亡數的中位數，以對數刻度繪於年齡上，須看兩件事，
即曲線的最低點是否進入個位數，以及最高與最低相差幾個數量級。第二項診斷取
各年齡組死亡數為零的年份比例，回答的不是「有多少零」而是「零在哪裡」。
第三項診斷以對數死亡率沿時間的二階差分估計噪音變異，二階差分消去線性趨勢，
而 $\mathrm{Var}(\Delta^2 X)=6\sigma^2$，故
$\widehat{\mathrm{Var}}=\mathrm{Var}(\Delta^2\log\hat m)/6$，以此對各年齡的
期望死亡數作雙對數散布圖並加上理論參考線 $1/D$。第四項診斷取中心化對數
死亡率矩陣的奇異值譜，並加上純噪音所能產生的上緣；依第\ref{sec:bbp}，
$n\times m$ 噪音矩陣的最大奇異值收斂至 $\sigma\sqrt{T}\,(1+\sqrt{c})$，
$c=A/T$，而該文式 (9) 另有一個便於操作的性質——在相變門檻上，觀測到的最大
奇異值恰好等於這條上緣，故判讀規則即為最大奇異值是否穿出虛線。

\begin{table}[H]
  \centering
  \caption{四項診斷的數值摘要}
  \label{tab:eda}
  \small
  \begin{tabular}{lrrr}
    \toprule
    診斷項目 & 全國女性 & 模擬 20 萬 & 模擬 5 萬 \\
    \midrule
    \cn{1} 最小單年死亡數中位數        & 69   & 1     & 0     \\
    \cn{1} 死亡數的跨年齡落差（倍）    & 150  & 298   & 72    \\
    \cn{2} 零死亡格比例（\%）          & 0.0  & 4.4   & 18.8  \\
    \cn{2} 零格最集中的年齡組          & ---  & 10--14 & 1--4 \\
    \cn{3} 噪音變異的跨度（數量級）    & 2.0  & 2.3   & 1.7   \\
    \cn{4} 最大奇異值 $s_1$            & 3.44 & 4.96  & 4.23  \\
    \cn{4} 純噪音的上緣                & 0.53 & 4.17  & 9.80  \\
    \cn{4} $s_1\,/\,$上緣              & \textbf{6.53} & \textbf{1.19} & \textbf{0.43} \\
    \midrule
    判讀 & 訊號充分 & 邊緣 & 不可還原 \\
    \bottomrule
  \end{tabular}
\end{table}

第四項診斷需要噪音的尺度 $\sigma$，其估法會直接影響結論，故須特別說明。
本文比較三種事前可得的估法，並以已知結論（人數二十萬以上可還原、五萬為
邊緣、一萬不可還原）檢驗。由觀測死亡數估計者可用但偏保守，若採用宜將門檻
放寬至 1.2；由一次 Firth 卜瓦松預備配適取得 $\hat\mu$ 並令
$\sigma=\sqrt{\overline{1/\hat\mu}}$ 者最為可靠，本文採用此法；而由二階差分
估計者不可使用，人數一萬時會誤判為可還原，因為四成零格被 0.5 墊平之後，
二階差分反映的是替代值的假象而非真實噪音。

\subsubsection{（二）判讀門檻與所指方法的權衡}

\begin{table}[H]
  \centering
  \caption{診斷\cn{4}的判讀門檻與對應建議}
  \label{tab:edarule}
  \small
  \begin{tabular}{llp{7.6cm}}
    \toprule
    $s_1/$上緣 & 判讀 & 建議 \\
    \midrule
    $>2$        & 訊號充分 & 各方法差異小，SVD 亦可 \\
    $1.2\sim2$  & 可還原   & 卜瓦松或加權 SVD（$w=\hat\mu$），加經驗貝氏收縮 \\
    $0.6\sim1.2$& 邊緣     & \textbf{方法選擇的效益最大}；卜瓦松／Firth 加收縮 \\
    $<0.6$      & 不可還原 & \textbf{不要以 SVD 估 $\beta_x$}；$\alpha_x$ 仍可靠 Firth 取得 \\
    \bottomrule
  \end{tabular}
\end{table}

表~\ref{tab:edarule} 最後一列的措辭須加以限定。相變門檻界定的是未正則化的
秩一分解能還原多少，並非所有估計量的上界；本文第伍節將顯示，人數五萬
（比值 0.43）時標準 LC 的 $\hat\beta_x$ 與真值相關僅 0.040，確實與隨機方向
無異，但加上經驗貝氏收縮後可達 0.542。因此該格應讀為「不要以 SVD 估計」，
而非「無論如何都不能估計」。

診斷所指向的方法各有代價，宜在此一併預告，以便閱讀後續各節時能對照。
表~\ref{tab:tradeoff} 整理五項對應關係及其優點與限制；各項的量化證據分別
見第伍節與第陸節，此處只列結論。

\begin{table}[H]
\centering
\caption{診斷所指向的方法及其權衡（量化證據見第伍、陸節）}
\small
\begin{tabular}{p{2.5cm}p{2.4cm}p{4.0cm}p{4.3cm}}
\toprule
診斷所示 & 指向的方法 & 優點 & 代價與限制\\
\midrule
\cn{1} 最小單年死亡中位數 $<5$
 & Firth 偏誤減少
 & 無調校參數；解必定有限；$\alpha$ 最大偏誤降 12.7 倍；大人口亦不劣於卜瓦松
 & 對 $\beta_x$ 的改善有限；其改善集中於幼年組，而該處僅占平均餘命誤差的兩成\\
\addlinespace
\cn{1}\cn{3} 跨年齡落差 $>1$ 個數量級、噪音跨度 $>2$ 個數量級
 & 資訊加權（卜瓦松，或 $w=\hat\mu$ 的加權 SVD）
 & 漂移項偏誤降 2.7 至 13 倍；就平均餘命而言是改善最大的一項
 & 完全不處理 $\alpha$ 的偏誤；權重若取觀測死亡數會與反應變數相關而產生內生性\\
\addlinespace
\cn{2} 零格比例 $<50\%$
 & 檢查函數（中位數 LC）
 & 可改善 $\alpha$ 且不需更換分配假設
 & 超過 50\%（中位數的崩潰點）即失效；且在小人口反而損害 $\beta_x$\\
\addlinespace
\cn{3} $\beta_x$ 為變異數主導
 & Fay-Herriot 經驗貝氏收縮
 & 免調校參數；$\beta_x$ 的相關可由 0.04 升至 0.54
 & 可交換先驗與真實的年齡結構不符；人數增大後偏誤轉為主要成分，固定強度的收縮反而有害\\
\addlinespace
\cn{4} $s_1/$上緣 $<0.6$
 & 不以 SVD 估計 $\beta_x$
 & 避免報告與隨機方向無異的估計值
 & 門檻為漸近理論的近似，而本文 $A=22$ 屬小維度；且僅界定未正則化的分解\\
\bottomrule
\end{tabular}
\label{tab:tradeoff}
\end{table}

綜合本文第伍節的結果，第一項診斷可給出兩個操作門檻：最小單年死亡數中位數
低於 5 時，Firth 為必要而非選配，低於 20 時建議使用；死亡數的跨年齡落差
超過一個數量級時，等權重 SVD 已被誤設，應改用卜瓦松或以 $\hat\mu$ 加權的
SVD。全國女性的 69 與 150 倍分別落在「安全」與「須加權」兩側，這說明即使
在國家級資料上加權仍屬必要，只是零格的問題不存在。

\subsection{參數估計的偏誤診斷}\label{sec:diag}

前一小節的診斷不需要真值，代價是只能指出假設失效與可能的後果，無法量化
偏誤的大小。本小節以電腦模擬補上這一層：令女性 2001 至 2024 年的 LC 配適
為真值，$E_{x,t}=N\cdot w_{x,t}$、
$D_{x,t}\sim\text{Poisson}(E_{x,t}m^{\text{true}}_{x,t})$，$N$ 取 $10^4$ 至
$5\times10^6$ 共九個水準，各重複 1{,}000 次。圖示規格比照 Wang et al.
（2018）第 354 至 355 頁。

除偏誤與標準誤外，本文另計算 T-ratio
\begin{equation}\label{eq:tratio}
T_x=\frac{\text{Bias}(\hat\theta_x)}{\mathrm{s.e.}(\hat\theta_x)} ,
\end{equation}
其意義為系統性偏誤相對於單次估計之抽樣標準差的倍數。$|T_x|>2$ 表示該年齡
的偏誤已超過抽樣變異的兩倍，亦即偏誤而非變異數才是該處的主要誤差來源，
反之 $|T_x|$ 很小則表示變異數主導。此一指標可直接回答「究竟該先處理偏誤
還是先處理變異數」，因而決定了後續應選用哪一類方法。

\begin{figure}[H]
\centering
\includegraphics[width=\textwidth]{figE3_tratio.png}
\caption{$\hat\alpha_x$ 與 $\hat\beta_x$ 的 T-ratio。
虛線為 $\pm2$ 參考線；超出者表示偏誤大於兩倍抽樣標準差。}
\label{fig:E3}
\end{figure}
\begin{table}[H]
\centering
\caption{各人口規模下 $|T|>2$ 的年齡組數與極值（共 22 組）}
\small
\begin{tabular}{rccrc}
\toprule
$N$ & $\alpha$ 之 $|T|>2$ 組數 & $\beta$ 之 $|T|>2$ 組數
 & $\alpha$ 之 $\max|T|$（年齡） & $\beta$ 之 $\max|T|$（年齡）\\
\midrule
$10^4$            & \textbf{9} & 0 & 56.4\ (5--9) & 0.22\ (1--4)\\
$2\times10^4$     & 7 & 0 & 31.0\ (5--9) & 0.21\ (5--9)\\
$5\times10^4$     & 3 & 0 & \phantom{0}9.4\ (5--9) & 0.09\ (1--4)\\
$10^5$            & 1 & 0 & \phantom{0}2.6\ (5--9) & 0.05\ (10--14)\\
$2\times10^5$     & 0 & 0 & \phantom{0}1.2\ (15--19) & 0.05\ (85--89)\\
$5\times10^5$     & 0 & 0 & \phantom{0}1.3\ (5--9) & \textbf{1.47}\ (80--84)\\
$10^6$            & 0 & 0 & \phantom{0}1.1\ (5--9) & \textbf{1.87}\ (70--74)\\
$2\times10^6$     & 0 & 0 & \phantom{0}0.8\ (5--9) & 1.32\ (70--74)\\
$5\times10^6$     & 0 & 0 & \phantom{0}0.5\ (5--9) & 0.76\ (70--74)\\
\bottomrule
\end{tabular}
\label{tab:tratio}
\end{table}

$\hat\alpha_x$ 的 T-ratio 在小樣本時極為可觀，且集中於幼年組。人數一萬時
5 至 9 歲組達 56.4，二萬時為 31.0，五萬時仍有 9.4，至十萬降為 2.6，
二十萬以上後始全面落入 $\pm2$ 之內。此一型態與 Wang et al.（2018）所述
「$\alpha_x$ 的偏誤最為顯著」一致，而偏誤集中於幼年組的原因可由表
\ref{tab:datadesc} 直接讀出：5 至 9 歲組的單年死亡數中位數僅 69 人，在人數
一萬的情境下期望死亡數不足 1，零格與對數轉換的影響於此最為劇烈。值得注意
的是，這個位置不需要模擬就能指出——前一小節的第二項診斷所畫的正是同一件事
的事前版本，兩者指向同一組年齡。

$\hat\beta_x$ 的情形與 $\alpha_x$ 截然不同。其 T-ratio 在小樣本時極小，
人數二十萬以下時最大值不及 0.25，隨人口增大反而上升至一百萬時的 1.87
（70 至 74 歲組），其後又下降至五百萬時的 0.76，呈現明顯的非單調型態。
此一型態不可誤讀：小樣本時 $|T|$ 接近零並不代表 $\hat\beta_x$ 估得好，
而是其抽樣標準差大到使偏誤相形之下微不足道，人數一萬時 $\hat\beta_x$ 與
真值的相關中位數為 $-0.24$，估計值實質上已無資訊。待人數增至五十萬以上、
變異數收斂之後，原先被掩蓋的偏誤才浮現為顯著，此時偏誤方成為誤差的主要
成分。

本小節的結論是三個參數不應以同一種方法處理。$\alpha_x$ 在小樣本下偏誤主導，
須以校正偏誤的手段處理；$\beta_x$ 在小樣本下變異數主導，收縮類方法於此有其
正當性，惟人數五十萬以上後偏誤轉為主要成分，此時固定強度的收縮反而有害，
此即第肆節所觀察到「固定比例的懲罰在人數三十萬以上後反受其害」的成因；
至於漂移項，其偏誤的來源與 $\alpha_x$ 並不相同，須待第伍節的階梯式分解
方能辨識。

\section{在高斯架構內修補的限度}\label{sec:limits}

在更換分配假設之前，須先確認留在高斯架構內的兩類常見修補確實不足，否則
本文後續的工作即無必要。這兩類是年齡方向的修勻與參數方向的懲罰，本節
各以一段交代其界限，詳細的數值見附錄\ref{sec:appmoved}。

修勻方面，本文以 Rabbi and Mazzuco（2021）的 $\LCe$ 為例。在其原文所用
的資料上，該方法確實如原文所述具有較小的配適誤差
（表~\ref{tab:smooth}），但樣本外表現僅在極少數設定下最佳
（表~\ref{tab:panel}），而其預測區間則明顯過窄（表~\ref{tab:e0}）。
機制不難理解：修勻降低的是觀測的變異，而區間寬度所需反映的是參數的
不確定性，兩者方向相反。修勻因此可以讓曲線好看，卻同時使不確定性被
低估——這對編製生命表是危險的，因為使用者拿到的是一個看似精確的數字。
更根本的是，修勻處理的是變異數而非偏誤：第\ref{sec:decomp}將顯示
$\alpha_x$ 的偏誤來自對數轉換與零格替代，這兩者都不因鄰近年齡的資訊
加入而消失。

懲罰式估計方面，本文發現一項結構性的限制。以 LC 的識別條件
$\sum_x\beta_x=1$ 為約束時，對 $\boldsymbol{\beta}$ 施加朝零收縮的懲罰
在最適解處不起作用——因為約束已固定了 $\boldsymbol{\beta}$ 的總和，
懲罰項成為常數，估計值完全不動（表~\ref{tab:shrink}）。把收縮目標改為
一組有意義的參考值之後懲罰才有效（表~\ref{tab:bv}、表~\ref{tab:rate}），
但此時又出現第二個限制：最適的懲罰強度須由真值決定，而使用者沒有真值；
以交叉驗證代替則須切分年齡--年份格，而切分會破壞雙線性結構所依賴的
跨格借力。此外，$L_1$ 懲罰所隱含的稀疏性先驗在此並不正確
（表~\ref{tab:pen}）——$\beta_x$ 沒有理由為零，它描述的是各年齡對時間
趨勢的敏感度。附錄\ref{sec:apppen}記錄實作細節與演算法的收斂性。

兩者合起來指出同一件事：\textbf{高斯架構內的修補能處理變異數，不能處理
偏誤；而偏誤是小人口 LC 的主要問題。} 這也說明為何本文第\ref{sec:est}
所選的三種方法有兩種（卜瓦松概似與 Firth 懲罰）根本不在高斯架構內，
而第三種（資訊加權）雖留在高斯架構內，改變的卻是目標函數的權重而非
平滑度。

\section{偏誤的來源與估計方法的選擇}

第參節已指出 $\alpha_x$ 與漂移項是偏誤主導，第肆節則顯示在高斯架構內修補
無法移除這項偏誤。本節先辨識偏誤來自哪一個環節，再依機制選用估計量並
量化其改善幅度與適用條件。

\subsection{偏誤來源的階梯式分解}\label{sec:decomp}

第貳節指出，高斯與卜瓦松兩條路線的差異綁了三件事：是否對資料取對數、
是否加權、以及分配假設本身。要判斷偏誤由何者造成，須把三者拆開，每次
只移除一個因素。本文據此構造五個估計量所成的階梯（表~\ref{tab:ladder}），
以女性 2001 至 2024 年的配適為真值、人數五萬、重複 400 次。由於少數模擬
會發散（單次平方誤差可達中位數的 $10^6$ 倍），統計量一律採中位數而非
平均數。

\begin{table}[H]
\centering
\caption{偏誤分解的估計量階梯}
\small
\begin{tabular}{llc}
\toprule
代號 & 估計量 & 仍存在的因素\\
\midrule
A & 標準 LC：對 $\log(\max(D,0.5)/E)$ 做 SVD & (1)(2)(3)(4)\\
B & A ＋ Jensen 一階校正 $+1/(2D)$ & (1)(3)(4)\\
C & 以 $D$ 為權重的加權 SVD（Wilmoth 1993） & (1)(2)(4)\\
D & 卜瓦松最大概似（不取對數） & (4)\\
E & \textbf{高斯對照}：對真值 $\log m$ 加均值為零、
    變異數與卜瓦松相當的噪音後做 SVD & \textbf{僅 (4)}\\
\bottomrule
\end{tabular}
\label{tab:ladder}
\end{table}
\begin{table}[H]
\centering
\caption{偏誤來源的階梯分解（重複 400 次，穩健統計）}
\small
\begin{tabular}{llrrrr}
\toprule
$N$ & 估計量 & $\alpha$ 偏誤（中位數） & SSE($\beta$) 中位數
 & $\beta$ 相關 & 漂移項偏誤\\
\midrule
\multirow{5}{*}{$5\times10^4$}
 & A 標準 LC        & $-0.0371$ & 0.409 & $-0.07$ & $0.305$\\
 & B ＋Jensen 校正  & $\phantom{-}0.0115$ & 0.236 & $-0.36$ & $0.199$\\
 & C 加權 SVD       & $\phantom{-}0.0889$ & 0.088 & $\phantom{-}0.03$ & $0.141$\\
 & D 卜瓦松 MLE     & $-0.0107$ & \textbf{0.071} & \textbf{0.41} & \textbf{0.011}\\
 & E 高斯對照       & $\phantom{-}\mathbf{0.0008}$ & 0.921 & $\phantom{-}0.27$ & $\mathbf{0.347}$\\
\midrule
\multirow{5}{*}{$2\times10^5$}
 & A 標準 LC        & $-0.0214$ & 0.061 & 0.51 & $0.175$\\
 & B ＋Jensen 校正  & $\phantom{-}0.0028$ & 0.048 & 0.21 & $0.155$\\
 & C 加權 SVD       & $\phantom{-}0.0322$ & 0.016 & 0.51 & $0.071$\\
 & D 卜瓦松 MLE     & $-0.0008$ & \textbf{0.015} & \textbf{0.67} & \textbf{0.004}\\
 & E 高斯對照       & $\phantom{-}\mathbf{0.0008}$ & 0.149 & 0.47 & $\mathbf{0.259}$\\
\bottomrule
\end{tabular}
\label{tab:decomp}
\end{table}

變體 E 是回答本問題的關鍵。它對真值的對數死亡率加入均值為零、變異數與
卜瓦松相當的高斯噪音後再做 SVD，因而完全不涉及對數轉換、不存在零格，
噪音亦對稱而無重尾。結果 $\alpha_x$ 的最大偏誤由標準 LC 的 0.743 降至
0.021，幾乎歸零；但漂移項的偏誤仍達 0.347，甚至略大於標準 LC 的 0.305。
換言之，若偏誤僅來自對數轉換，變體 E 的漂移項亦應改善，事實並非如此，
因此對數轉換不是唯一的來源。

漂移項的偏誤來自何處，可由變體 C 得到答案。該變體仍在對數尺度上操作、
仍有零格，只是把等權重改為以死亡數加權，漂移項偏誤即由 0.305 降至 0.141、
$\beta$ 的平方誤差和由 0.409 降至 0.088。加權所改變的只有各格在目標函數中
的相對重要性，因此漂移項的偏誤來自等權重分解對異質變異的敏感——用
第\ref{sec:bbp}的語言說，即主奇異向量被噪音旋轉。此一機制與對數無關，這也是
變體 E 的偏誤未能改善的原因。

本小節可歸納為一句話：取對數傷 $\alpha_x$，不加權傷漂移項。兩者是獨立的病，
也需要獨立的藥。

\subsection{五種估計量的比較}\label{sec:est}

依上述診斷，本小節比較五種估計量，每一種各處理一個機制
（表~\ref{tab:estdesign}）。模擬設定同前，重複 100 次、種子固定。

\begin{table}[H]
  \centering
  \caption{五種估計量各自處理哪一個機制}
  \label{tab:estdesign}
  \small
  \begin{tabular}{llccc}
    \toprule
    代號 & 估計量 & 取對數 & 加權 & 收縮 \\
    \midrule
    A & 標準 LC（現行作法）        & 是 & 無     & 無 \\
    B & 加權 SVD（$w=\hat\mu$）    & 是 & $\hat\mu$ & 無 \\
    C & 卜瓦松最大概似             & 否 & 內建   & 無 \\
    D & Firth 偏誤減少             & 否 & 內建   & 無 \\
    E & Firth $+$ Fay-Herriot 收縮 & 否 & 內建   & 有 \\
    \bottomrule
  \end{tabular}
\end{table}

變體 B 須特別說明。Wilmoth（1993）原式以觀測死亡數為權重，但此舉使權重與
反應變數相關——該格恰好觀測到較多死亡時同時被賦予較大的權重，$\hat\alpha$
因而被系統性推高。表~\ref{tab:decomp} 的變體 C 即依原式實作，其 $\alpha$
偏誤為 $+0.0889$，是五個變體中唯一明顯為正者。本小節改以配適期望死亡數
$\hat\mu_{x,t}$ 為權重並隨迭代更新，此為一行程式碼的修改。

\begin{table}[H]
  \centering
  \caption{五種估計量的比較（重複 100 次，取中位數）。$\alpha$ 欄為逐年齡
  偏誤中位數的絕對值，分別取其中位數與最大值；SSE($\beta$) 為相對平方誤差和；
  cor($\beta$) 為與真值的中心化相關。}
  \label{tab:est}
  \small
  \begin{tabular}{llrrrrr}
    \toprule
    $N$ & 估計量 & $\alpha$ 中位 & $\alpha$ 最大 & SSE($\beta$) & cor($\beta$) & 漂移偏誤 \\
    \midrule
    \multirow{5}{*}{\shortstack[l]{$10^4$\\ 零格 206/528}}
      & A 標準 LC        & 0.1612 & 2.3309 & 7.916 & $-0.074$ & 0.3436 \\
      & B 加權 SVD       & 0.1128 & 2.3315 & 3.319 & 0.068 & 0.2436 \\
      & C 卜瓦松 MLE     & 0.0495 & \textbf{13.0317} & 15.525 & 0.118 & 0.1612 \\
      & D Firth          & \textbf{0.0248} & \textbf{0.4267} & 17.797 & 0.116 & 0.2237 \\
      & E Firth $+$ FH   & \textbf{0.0248} & \textbf{0.4267} & 4.984 & 0.125 & 0.2237 \\
    \midrule
    \multirow{5}{*}{\shortstack[l]{$5\times10^4$\\ 零格 100/528}}
      & A 標準 LC        & 0.0611 & 0.8658 & 4.449 & 0.037 & 0.3339 \\
      & B 加權 SVD       & 0.0553 & 0.8723 & 0.904 & 0.201 & 0.1223 \\
      & C 卜瓦松 MLE     & 0.0073 & 0.2260 & 1.756 & 0.365 & 0.0492 \\
      & D Firth          & 0.0088 & \textbf{0.0680} & 1.410 & 0.397 & 0.0485 \\
      & E Firth $+$ FH   & 0.0088 & \textbf{0.0680} & \textbf{0.143} & \textbf{0.542} & 0.0485 \\
    \midrule
    \multirow{5}{*}{\shortstack[l]{$2\times10^5$\\ 零格 29/528}}
      & A 標準 LC        & 0.0231 & 0.1846 & 1.400 & 0.379 & 0.2117 \\
      & B 加權 SVD       & 0.0234 & 0.1397 & 0.305 & 0.527 & 0.0393 \\
      & C 卜瓦松 MLE     & \textbf{0.0018} & \textbf{0.0361} & 0.346 & 0.596 & 0.0267 \\
      & D Firth          & 0.0025 & 0.0398 & 0.333 & 0.605 & 0.0262 \\
      & E Firth $+$ FH   & 0.0025 & 0.0398 & \textbf{0.095} & \textbf{0.697} & 0.0262 \\
    \bottomrule
  \end{tabular}
\end{table}

$\alpha_x$ 的結果最為明確。人數五萬時，逐年齡的最大偏誤由標準 LC 的 0.866
降至 Firth 的 0.068，改善 12.7 倍。人數一萬那一組更值得注意：該規模下
528 格中有 206 格為零，標準 LC 與卜瓦松最大概似的偏誤方向相反而且都極大。
標準 LC 把 5 至 9 歲組的 $\alpha$ 推高 $+2.33$，因為零格以 0.5 替代而該處
的期望死亡數遠低於 0.5，等於人為墊高該格的死亡率；卜瓦松最大概似則把它
推低 $-13.03$，因為零計數下概似在 $\alpha\to-\infty$ 時仍遞增，估計量根本
不存在，此即第貳節所述完全分離現象在卜瓦松上的對應物；Firth 為 $+0.43$，
相對標準 LC 改善 5.5 倍、相對卜瓦松 30 倍。這一組數字修正了一個常見的
簡化說法：「改用卜瓦松路線比較好」在零格眾多時並不成立，單以逐年齡最大
偏誤衡量，此時卜瓦松反而比 SVD 更糟，正確的陳述是「卜瓦松加 Firth 最好」。

變體 B 與 A 的對照則是本節最乾淨的一項證據。人數五萬時把權重加入高斯路線
之後，漂移項偏誤由 0.3339 降至 0.1223、$\beta$ 的平方誤差和由 4.449 降至
0.904，但 $\alpha$ 的最大偏誤紋風不動，0.8658 對 0.8723。因為加權完全沒有
觸及零格與對數轉換。這與前一小節變體 E 的結果恰成互補的一對——一個不取
對數但不加權，$\alpha$ 歸零而漂移項最差；一個取對數但加權，漂移項改善而
$\alpha$ 不動——兩者合起來確立了本文的核心分工。

$\beta_x$ 則須分規模討論。人數一萬時五種估計量的相關係數全在 0.125 以下，
與第四項診斷的「不可還原」判讀一致，此規模下應誠實地不予估計。人數五萬起
情況改變，標準 LC 的 0.037 確實與隨機方向無異，但 Firth 加上 Fay-Herriot
收縮後可達 0.542，平方誤差和由 4.449 降至 0.143；人數二十萬時則由 0.379
提升至 0.697。收縮之所以有效，正因為第參節已指出 $\beta_x$ 在此規模是
變異數主導，而 Fay-Herriot 的收縮強度由估計出的變異成分決定，沒有調校參數，
因而避開了第肆節所指的選參準則錯置問題。

\subsection{改善的位置與事前診斷的對應}

若第參節的主張成立，則估計量的改善應出現在診斷所指的年齡，而非散布全表。
圖~\ref{fig:edabias} 以淺色柱表示各年齡的零死亡格比例、以折線表示各估計量
的 $\alpha$ 偏誤，兩者的峰值完全重疊。

\begin{figure}[H]
  \centering
  \includegraphics[width=\textwidth]{figH_eda_vs_alpha_bias.png}
  \caption{逐年齡的 $\alpha$ 偏誤（折線）與零死亡格比例（淺色柱，相對高度）。
  標準 LC 的偏誤峰值與零格所在的年齡完全重疊；卜瓦松在同處往反方向發散；
  Firth 則近乎水平。}
  \label{fig:edabias}
\end{figure}

若以逐年齡的零格比例與 $|\alpha$ 偏誤$|$ 計算 Spearman 相關，標準 LC 在
人數一萬、五萬、二十萬三個規模下分別為 0.863、0.881、0.740。亦即第
\ref{sec:diag} 小節圖中 T-ratio 的尖峰位置，不需要真值、不需要一千次模擬，
看第二項診斷即可事前指出。

此處有一項判讀上的陷阱須明說。Firth 的 Spearman 相關在人數一萬時為 0.915，
並不低於標準 LC，但這不表示 Firth 沒有改善——Spearman 相關對尺度不變，
衡量的是「偏誤落在哪些年齡」而非「偏誤有多大」，Firth 改善的是同一批年齡
上的量級。值得注意的是 Firth 的相關隨人口增大而下降，三個規模依序為
0.915、0.789、0.484，顯示其殘餘偏誤在大人口下已不再由零格驅動。

\subsection{換損失函數的上界}\label{sec:qlc}\label{sec:mq}\label{sec:hu}

前三小節顯示偏誤的兩個來源分別是對數轉換與等權重分解，而有效的作法是
更換分配假設與加權。本小節檢驗第三條可能的路——更換損失函數——並劃出
它的上界。此處把四項檢驗合併敘述，因為它們的結論指向同一件事；
完整的數值表見附錄\ref{sec:appmoved}。

\subsubsection{（一）問題與既有文獻的前提}

依表~\ref{tab:gap}，穩健損失一脈的前提是污染為少數、對稱、且與觀測的
資訊量無關。本文的情境三項皆不成立，但不成立的方式不同，故須分開檢驗：
丟棄觀測的方法（截尾）與只改變損失形狀的方法（檢查函數、Huber）
在原理上不同，而以「年」為單位的穩健化又與以「格」為單位者不同。

第一項檢驗針對截尾。古典穩健統計處理的是污染，但第\ref{sec:decomp}的
變體 E 已排除此一情境：該變體噪音對稱、期望為零、無重尾、無零格，
卻仍產生 0.347 的漂移項偏誤。既然沒有東西可截，截尾應無從發揮。
實測確認了這點（附錄表~\ref{tab:robust}）：沒有任何截尾變體勝過未截尾的
卜瓦松最大概似；最小截尾平方完全不處理漂移項的偏誤，人數五萬時為 0.3055
而標準 LC 為 0.3144，$\beta$ 的相關係數 0.041 與標準 LC 的 0.040 無異。


截尾失效的機制可由逐年齡的保留率看得更清楚。表~\ref{tab:retain} 於人數
五百萬、模型正確、無零格亦無離群值的條件下統計各年齡被保留的比例；若截尾
是中性的，各年齡的保留率應大致相同。

\begin{table}[H]
\centering
\caption{截尾規則的逐年齡保留率（$N=5\times10^6$，模型正確、無零格、無離群值）}
\small
\begin{tabular}{lrrrr}
\toprule
年齡組 & $\beta_x$ 真值 & MTL 保留率 & LTS 保留率 & 單年期望死亡數\\
\midrule
0        & 0.0496 & 0.98 & 0.88 & 111\\
1--4     & \textbf{0.1066} & 0.99 & \textbf{0.36} & \phantom{0,0}28\\
5--9     & \textbf{0.0697} & 1.00 & \textbf{0.52} & \phantom{0,0}24\\
10--14   & 0.0442 & 1.00 & 0.60 & \phantom{0,0}25\\
15--19   & 0.0578 & 0.99 & 0.71 & \phantom{0,0}50\\
40--44   & 0.0328 & 0.95 & 1.00 & \phantom{0,}485\\
70--74   & 0.0669 & 0.82 & 1.00 & 4,619\\
80--84   & 0.0439 & \textbf{0.71} & 1.00 & \textbf{6,487}\\
85--89   & 0.0319 & \textbf{0.71} & 1.00 & \textbf{7,233}\\
100$+$   & 0.0078 & 0.94 & 1.00 & \phantom{0,}503\\
\bottomrule
\end{tabular}
\label{tab:retain}
\end{table}

最小截尾平方丟掉了 64\% 的 1 至 4 歲組與 48\% 的 5 至 9 歲組，而這兩組
正是 $\beta_x$ 最大的兩個年齡（0.1066 與 0.0697）。機制很直接：
在對數尺度上 $|\beta_x\kappa_t|$ 最大的年齡也是偏離秩一結構時殘差最大的
年齡，因此「最不合模型」的格恰好是「攜帶最多趨勢訊號」的格。
最大截尾概似則偏好丟掉 80 至 94 歲，那是死亡數最多的年齡。
\textbf{兩種截尾規則丟掉的是不同的格，但都是錯的格。}

\subsubsection{（二）方法：把穩健程度由開關改成刻度}

截尾失效不表示所有穩健工具都無用，因為檢查函數只改變損失的形狀而不丟棄
觀測。但與其比較「標準 LC」與「中位數 LC」兩個點，更有信息的是認識到
兩者其實是式~\eqref{eq:mq} 同一族的兩個極限：$\tau=0.5$ 且 $c\to\infty$
時退化為對數尺度的最小平方，$\tau=0.5$ 且 $c\to0$ 時退化為檢查函數。
於是「要不要穩健」變成「$(\tau,c)$ 該取在哪裡」。

求解採 Hunter and Lange（2000）的 MM 演算法：$\rho_{\tau,c}$ 的一階導數
可寫成 $W(u)u$，其中 $W(u)=2w_\tau(u)\min\{1,c/(|u|+\epsilon)\}$，
故每一步化為加權最小平方，雙線性結構使兩個區塊各有封閉解，
無須線性規劃求解器。擾動項 $\epsilon$ 隨迭代縮小。此一實作有一項嚴格的
內部對照：取 $\tau=0.5$、$c=\infty$ 時，所得參數與直接對
$\log m_{x,t}$ 作 SVD 者的最大差異為 $10^{-15}$ 量級，亦即機器精度。
亂數種子與表~\ref{tab:est} 完全相同。

\subsubsection{（三）實證：$\tau$ 的方向與 $c$ 的內部最適}



$(\tau,c)$ 的完整掃描結果見附錄表~\ref{tab:mqc} 與表~\ref{tab:mqtau}，
此處只摘出兩項。其一，穩健程度有\textbf{內部最適}。$\alpha$ 偏誤中位數的最適
截點隨人口規模單調下移，人數一萬與五萬時落在 $c=0.7$、人數二十萬時落在
$c=0$（附錄表~\ref{tab:mqc}）；而在漂移項上，人數二十萬時 $c=0.7$ 的偏誤
為 0.0670，兩個端點分別為 0.1111（中位數）與 0.2117（最小平方），
亦即\textbf{中位數只走到一半}。以兩個點作判斷會漏掉中間這一段。

其二，非對稱參數的效果更大且方向明確（附錄表~\ref{tab:mqtau}）。$\tau=0.6$ 在三個規模下都使
$\alpha$ 偏誤中位數大幅下降，人數一萬時由 0.1427 降至 0.0542；
$\tau=0.3$ 與 $\tau=0.4$ 則一律惡化。以只對負殘差設截點（壓左尾，
Xu and Chen 2018）實作亦得一致的方向，三個規模的 $\alpha$ 偏誤中位數為
0.0566、0.0226、0.0047，是全部掃描中最低的一組。兩種互相獨立的實作都
指向\textbf{受損的是殘差的左側}。

\subsubsection{（四）假設的檢驗：扭曲是否為單邊}

但 $\tau=0.6$ 的 $\alpha$ \textbf{最大}偏誤並未改善（人數一萬時 2.3249
對 $\tau=0.5$ 的 2.3156）。改了大多數年齡卻沒改到最壞的那個，這個型態
本身就是線索。


\begin{table}[H]
\centering
\caption{$\alpha_x$ 偏誤隨 $\tau$ 的變化（$c=1.345$，$N=10^4$，各年齡取重複的中位數）}
\small
\begin{tabular}{lrrrrrrr}
\toprule
年齡組 & 單年期望死亡 & $\tau=0.3$ & $\tau=0.4$ & $\tau=0.5$ & $\tau=0.6$ & $\tau=0.7$ & 最適 $\tau$\\
\midrule
5--9   & \phantom{00}0.1 & \textbf{2.3070} & 2.3105 & 2.3156 & 2.3249 & 2.3461 & \textbf{0.3}\\
20--24 & \phantom{00}0.2 & \textbf{1.1032} & 1.1184 & 1.1424 & 1.1790 & 1.2638 & \textbf{0.3}\\
30--34 & \phantom{00}0.4 & \textbf{0.3981} & 0.4252 & 0.4745 & 0.5831 & 0.6865 & \textbf{0.3}\\
35--39 & \phantom{00}0.5 & \textbf{0.0317} & 0.0852 & 0.1866 & 0.3208 & 0.4504 & \textbf{0.3}\\
\midrule
40--44 & \phantom{00}1.0 & $-0.3402$ & $-0.2245$ & $-0.0762$ & \textbf{0.0311} & 0.1558 & \textbf{0.6}\\
50--54 & \phantom{00}2.0 & $-0.4894$ & $-0.3473$ & $-0.2042$ & \textbf{$-0.0644$} & 0.0724 & \textbf{0.6}\\
65--69 & \phantom{00}6.9 & $-0.2261$ & $-0.1350$ & $-0.0423$ & \textbf{0.0394} & 0.1119 & \textbf{0.6}\\
85--89 & \phantom{0}14.8 & $-0.1303$ & $-0.0711$ & $-0.0271$ & \textbf{0.0223} & 0.0755 & \textbf{0.6}\\
\bottomrule
\end{tabular}
\label{tab:mqage}
\end{table}

$\alpha_x$ 的偏誤在不同年齡\textbf{符號相反}，因而最適的 $\tau$ 朝相反
方向移動。單年期望死亡數不足一人的年齡偏誤為正，最適 $\tau$ 為掃描下限
0.3；一人以上的年齡偏誤為負，最適 $\tau$ 為 0.6。分界落在每年約一人死亡處，
而這正是診斷\cn{1}所量的死亡數量級。

兩個相反的機制解釋了這個符號差異，而兩者都源於取對數。期望死亡數遠小於 1
時多數年份觀測到零，以 0.5 替代必然使該格死亡率被\textbf{高估}
（$0.5/E_{x,t}>\mu_{x,t}/E_{x,t}$），故偏誤為正。期望死亡數大於 1 時
不再有零格，但 $\mathbb{E}[\log D]<\log\mathbb{E}[D]$ 這個 Jensen 不等式
使對數尺度上的觀測被系統性\textbf{低估}，量級約 $-1/(2\mu_{x,t})$。
後者可定量核對：人數二十萬時 85 至 89 歲組的單年期望死亡數為 295.7，
理論值 $-1/(2\times295.7)=-0.0017$，實測 $\tau=0.5$ 的偏誤為 $-0.0010$。

由此可得本節認為最值得強調的結論：\textbf{單一的 $\tau$ 無法同時服務兩個
機制。} 它改善占多數的稠密年齡，卻無法改善零格集中的幼年組。這同時為
Santolino（2020）「參數隨 $\tau$ 並無明顯規律」的觀察提供了解釋——
若最適的 $\tau$ 本身隨年齡變號，則以單一 $\tau$ 估出的參數本就不應呈現
規律。若允許 $\tau$ 隨年齡變動並以診斷\cn{1}決定其方向，應可同時改善
兩組，此點列入第\ref{sec:future}的後續工作。

\subsubsection{（五）以「年」為單位的穩健化：門檻從未被觸發}

最後一項檢驗針對 Hyndman and Ullah（2007）。該文以每一年整條曲線的積分
平方誤差判斷該年是否離群並降低其權重，其設想的離群值是使整個年度偏離趨勢
的事件——該文在法國資料上辨識出的離群年為 1914 至 1919 年、
1940 至 1945 年與 1960 年，前兩段是兩次世界大戰與 1918 年西班牙流感所致。
其前提因此是「少數異常年份對多數正常年份」。本文的噪音來自每一格獨立的
卜瓦松抽樣，各年的年齡結構與期望死亡數相同，故各年偏離秩一結構的程度在
統計上同分布，並不存在可供對比的多數。之所以仍須單獨檢驗，是因為前述
截尾以「格」為單位，若小人口的噪音恰在某些年份同向累積，逐年穩健化仍可能
有效；由格推到年需要格間噪音獨立這個額外前提，而它本身就是檢驗的對象。

實測的結果與預期一致，而一致的程度本身值得注意（附錄表~\ref{tab:hu}）。
平滑降權規則在 24 年之中平均只降權 0.59 至 0.67 年，亦即三次模擬有兩次
一年都不降；所得估計與標準 LC 在小數第三位才開始有差異，漂移項偏誤在
人數一萬時兩者皆為 0.3436。\textbf{不存在離群年，因此沒有東西可以降權。}
硬性剔除則顯示強迫它動手只會變差：剔除 4 年後 $\beta$ 的平方誤差和在三個
規模下一致惡化一成至兩成。

\subsubsection{（六）小結：上界在哪裡，以及為什麼}

把四項檢驗並列，可以看出「穩健」一詞在本問題上有兩種不同的失效方式。
逐格截尾是\textbf{做錯事}——它確實丟掉了格，而丟掉的恰好是攜帶訊號的格；
逐年降權則是\textbf{什麼都沒做}——門檻從未被觸發，因為判斷離群所需的
多數／少數對比並不存在。M-分位數之所以可行，正因為它不判斷任何觀測是否
離群、也不丟棄任何觀測，故不需要這個前提。

但即使把穩健由開關改成刻度，上界並未改變。把 Huber 損失與資訊加權疊加
之後（附錄表~\ref{tab:mqK}），$\beta$ 的平方誤差和較前節最佳的加權中位數
再降兩成至三成，是損失函數這條路上最好的結果；然而其 $\alpha$ 最大偏誤在
人數五萬時仍為 0.7853，而 Firth 為 0.0680（表~\ref{tab:est}），
相差十倍以上。理由是：$(\tau,c)$ 所調整的是損失函數對\textbf{已產生的}
錯誤資訊有多敏感，而 Firth 在計數尺度上估計、根本不產生那筆錯誤資訊。

\textbf{就本文的問題而言，換分配假設優於換損失函數；而換損失函數的上界
不是調參不足，是同一個損失函數必須同時服務兩個方向相反的機制。}
依表~\ref{tab:penpos} 的分類，這是「加在損失形狀上」這個位置的固有界限，
與「加在概似上」所能達到的效果不同層級。

\subsection{從參數到平均餘命}\label{sec:e0}

前述各小節皆以參數本身為評估對象，但本文開頭的問題是能否編製鄉鎮市區
生命表，其標的為平均餘命而非參數。兩者未必同向：Firth 所改善的 $\alpha_x$
偏誤集中於幼年組，而該處的死亡率極低，對平均餘命的影響未必與對數尺度上的
偏誤成比例。本小節因此把估計流程接到生命表，直接檢驗參數的改善能否轉移。

設定同前，重複 100 次、種子與表~\ref{tab:est} 相同。以各估計量配適後的
$\hat m_{x,T}$ 代入簡略生命表求零歲平均餘命，與真值比較；另以
$\hat\kappa_T+10\hat d$ 外推十年，比較點預測的平均餘命。真值的配適末年
$e_0$ 為 84.069 歲、十年後為 85.700 歲。

\begin{table}[H]
\centering
\caption{零歲平均餘命的偏誤與均方根誤差（歲；重複 100 次，取中位數）}
\small
\begin{tabular}{llrrrr}
\toprule
$N$ & 估計量 & 配適偏誤 & 配適 RMSE & 預測偏誤 & 預測 RMSE\\
\midrule
\multirow{5}{*}{$10^4$}
 & 標準 LC     & $-3.621$ & 3.621 & $-5.105$ & 5.105\\
 & 加權 SVD    & $-3.241$ & 3.242 & $-4.019$ & 4.020\\
 & 卜瓦松 MLE  & \textbf{$-0.820$} & \textbf{0.961} & $-2.981$ & 2.981\\
 & Firth       & $-1.327$ & 1.327 & $-2.957$ & 2.957\\
 & Firth $+$ FH & $-1.176$ & 1.257 & \textbf{$-2.667$} & \textbf{2.689}\\
\midrule
\multirow{5}{*}{$5\times10^4$}
 & 標準 LC     & $-1.480$ & 1.480 & $-2.664$ & 2.664\\
 & 加權 SVD    & $-0.394$ & 0.494 & $-0.633$ & 0.683\\
 & 卜瓦松 MLE  & $-0.202$ & \textbf{0.351} & $-0.512$ & 0.609\\
 & Firth       & $-0.282$ & 0.361 & $-0.570$ & 0.661\\
 & Firth $+$ FH & \textbf{$-0.182$} & 0.427 & \textbf{$-0.166$} & \textbf{0.592}\\
\midrule
\multirow{5}{*}{$2\times10^5$}
 & 標準 LC     & $-0.899$ & 0.899 & $-2.011$ & 2.011\\
 & 加權 SVD    & $-0.123$ & 0.212 & $-0.142$ & 0.258\\
 & 卜瓦松 MLE  & $-0.031$ & \textbf{0.200} & \textbf{$-0.086$} & \textbf{0.242}\\
 & Firth       & $-0.050$ & \textbf{0.200} & $-0.098$ & 0.247\\
 & Firth $+$ FH & \textbf{0.000} & 0.214 & 0.072 & 0.314\\
\bottomrule
\end{tabular}
\label{tab:e0bias}
\end{table}

參數的改善確實轉為平均餘命的改善，而且幅度足以影響實務判讀。人數五萬時
標準 LC 低估配適末年的平均餘命達 1.480 歲、低估十年後的點預測達 2.664 歲；
改用卜瓦松最大概似後分別降為 0.202 與 0.512 歲。考慮到我國縣市間平均餘命
的最大差距約為 10 歲（Yue et al. 2019），2.664 歲的偏誤已足以使一個地區
在排序上錯置數個名次，因此這並非僅具統計意義的差距。

值得注意的是配適與預測兩欄的收斂速度並不相同。標準 LC 的配適偏誤隨人口
規模由 3.621 降至 1.480 再降至 0.899，而預測偏誤僅由 5.105 降至 2.664 再
降至 2.011，下降得慢得多。原因是後者由漂移項主導，而第一小節已指出漂移項
的偏誤是等權重秩一分解的結構性產物，不隨人口規模快速消失。這也解釋了為何
人數二十萬時標準 LC 的配適偏誤已降至 0.899 歲、看似可用，其十年點預測卻
仍偏低 2.011 歲。

若必須留在高斯架構，以配適期望死亡數加權的 SVD 可取回大部分的改善：人數
五萬時配適偏誤由 1.480 降至 0.394 歲，約為卜瓦松所達改善幅度的 85\%。
這對實務有直接意義，因為加權 SVD 只需修改權重而不需改寫估計架構。

\subsubsection{（一）誤差來自哪些年齡}

前述改善固然明確，但改善的來源是否即為第三小節所指的幼年組，仍須檢驗。
本文因此逐一把某一年齡組的估計死亡率換回真值，觀察平均餘命的誤差減少多少，
以此衡量該年齡對總誤差的貢獻。

這張表修正了一項容易產生的誤解。$\alpha_x$ 的偏誤集中於幼年組而 Firth 在
該處的改善達 12.7 倍，但就平均餘命而言，30 歲以下各組合計僅貢獻 $-0.3049$ 歲、
占標準 LC 總誤差的 20.6\%；真正的主要來源是 60 至 89 歲，合計 $-0.9300$ 歲、
占 62.8\%——該區間各年齡在對數尺度上的偏誤其實不大（約 $-0.01$ 至 $-0.05$），
但其死亡率比幼年組高出兩到三個數量級。\textbf{換言之，診斷\cn{2}所指出的是
「對數尺度上偏誤最大的年齡」，而非「對平均餘命影響最大的年齡」。}
這並不削弱該診斷的價值，但若標的是平均餘命，則資訊加權的重要性更甚於偏誤
減少——以死亡數加權的估計量之所以在表~\ref{tab:e0bias} 中表現突出，
正是因為加權所強調的高齡組恰好是誤差的主要來源。須註明逐年齡貢獻的合計
（$-1.3043$ 歲）略小於基準誤差（$-1.480$ 歲），差額來自生命表的非線性轉換，
故該表應讀為相對重要性的排序而非嚴格的分解。



\subsection{小結：可估者已估，未估者為何}

本節自偏誤的來源出發，依序檢驗了五類作法，此處把假設與結果並列收束。
階梯式分解確認了第\ref{sec:gap}表~\ref{tab:gap} 所列的兩項落差確實存在：
取對數使 $\alpha_x$ 產生偏誤，而等權重使漂移項對異質變異敏感，且後者
與損失函數的形狀無關——變體 E 在無零格、無重尾、噪音對稱的條件下仍留下
0.347 的漂移項偏誤，這直接否證了「以更穩健的損失即可補救」這一假設。
依此分解選出的三種作法各自對應一項來源：卜瓦松概似與 Firth 懲罰處理
對數與零格所產生的錯誤資訊，資訊加權處理變異數的不均，Fay--Herriot
處理變異數主導的 $\beta_x$。三者皆無調校參數。

反面的結果同樣明確，而且其失效方式恰好驗證了表~\ref{tab:gap} 的前提落差。
逐格截尾（第五小節）丟掉的是 $\beta_x$ 最大、攜帶最多趨勢訊號的幼年組，
因為在雙線性模型中「最不合模型」與「最具資訊」高度重疊，這違反了
古典穩健統計「離群值與資訊量無關」的前提。逐年降權（第\ref{sec:hu}）
則從未被觸發，因為 Hyndman and Ullah（2007）所依的「少數異常年對多數
正常年」在小人口並不存在。M-分位數的連續掃描（第\ref{sec:mq}）雖確實
把損失函數這條路推進了一段——非對稱參數補回 Jensen 項、有限截點兼顧
效率與穩健——但仍停在同一道牆前：$\alpha$ 最大偏誤 0.785 對 Firth 的
0.068。三項負面結果指向同一個結論，即損失函數所能處理的是「對已產生的
錯誤資訊有多敏感」，而非「錯誤資訊要不要產生」。

至於不可估者，本節亦已界定。$\beta_x$ 在人數二十萬以下受訊噪比限制，
落在 BBP 門檻以下，任何估計量都無法還原其形狀，能做的只是以收縮避免
把噪音當訊號。第\ref{sec:e0}進一步顯示，參數層面的改善確實可以轉為
平均餘命的改善（配適期偏誤由 $-1.480$ 降至 $-0.202$ 歲），但改善的
年齡位置與平均餘命所重的年齡並不重合。

至此，點估計的部分已交代完畢：哪些參數可估、以什麼方法估、改善多少、
以及為何有些方法無效。但「能否編製鄉鎮市區生命表」這個問題並未因此
回答完整，因為使用者最終需要的不只是一個數字，而是一個附帶不確定性的
數字。而點估計的偏誤一旦存在，區間的中心就是偏的；區間的寬度又另有
其估計問題。下一節因此轉向區間推論，並沿用同一套作法——先分解誤差的
來源，再依來源選方法。

\section{區間推論}\label{sec:interval}

點估計的偏誤一旦存在，區間的中心就是偏的；區間的寬度又另有其估計問題。
本節沿用與前一節相同的作法——先分解誤差的來源，再依來源選方法——
並得到一個可分離為兩部分的結論。完整的數值見附錄\ref{sec:appmoved}。

\subsection{位置與寬度是兩個獨立的誤差}\label{sec:isrc}

以層層疊加的變體分解區間誤差（表~\ref{tab:idecomp}）後可知，位置的偏誤
直接繼承點估計：漂移項估錯，區間的中心就錯，因此以卜瓦松最大概似取代
SVD 即可校正，這是第\ref{sec:est}的結論在區間上的直接後果。
表~\ref{tab:drift} 進一步顯示漂移項的衰減並非小人口的固有性質，而是
SVD 對異質變異敏感所致。

寬度的偏誤則另有來源。樸素的 $\mathrm{sd}(\Delta\hat\kappa)$ 把測量噪音
誤認為過程變異，因而系統性高估 $\sigma$：人數五萬時其中位數為 1.495，
而真值為 0.679。把 $\kappa_t$ 寫成狀態空間模型——觀測式
$\hat\kappa_t=\kappa_t+u_t$、狀態式 $\kappa_t=\kappa_{t-1}+\delta+\varepsilon_t$，
其中 $u_t$ 的變異數由卜瓦松訊息量 $(\sum_x\hat\mu_{x,t}\hat\beta_x^2)^{-1}$
給定——即可把兩者分離，以 REML 估得的中位數降至 0.575，偏誤由 $+0.816$
縮為 $-0.104$（表~\ref{tab:sigma}）。人數十萬以下時另須加計 Kalman 濾波
所給出的初始值變異數 $P_T$，因為 $\hat\kappa_T$ 本身即有不確定性。

\subsection{兩步校正與「區間過窄」的檢驗}

把位置與寬度兩步合併之後的結果見表~\ref{tab:calib}。


\begin{table}[H]
\centering
\caption{校正前後的預測區間寬度與涵蓋率（名目 80\%）}
\small
\begin{tabular}{rlcc}
\toprule
$N$ & 方法 & 區間寬度 & 涵蓋率\\
\midrule
200,000 & oracle & 1.743 & 0.81\\
        & 現行作法 & 5.081 & 0.75\\
        & 本文校正版 & \textbf{1.724} & 0.63\\
\midrule
1,000,000 & oracle & 1.726 & 0.85\\
          & 現行作法 & 4.667 & 0.95\\
          & 本文校正版 & \textbf{1.721} & \textbf{0.83}\\
\midrule
11,900,000 & oracle & 1.712 & 0.78\\
           & 現行作法 & 2.033 & 0.82\\
           & 本文校正版 & \textbf{1.678} & 0.77\\
\bottomrule
\end{tabular}
\label{tab:calib}
\end{table}

此處也檢驗了一項文獻上的常見主張。LC 區間過窄的說法
（Lee and Miller 2001；Booth et al. 2006）建立在「$\sigma$ 被正確估計、
僅漏計了參數不確定性」這個前提上，而本文的實測顯示小人口的情形相反：
$\sigma$ 本身被高估，故現行作法的區間是\textbf{過寬}而非過窄——人數五萬時
寬度 8.949 而涵蓋率僅 0.27，寬度是不可約下限 2.093 的四倍餘，涵蓋率卻遠低於
名目的 80\%（表~\ref{tab:ival}）。過寬與涵蓋不足並存，是位置偏誤與寬度高估
同時存在的特徵，兩者不可混為一談。這再次印證第\ref{sec:gap}的判斷：
既有結論的前提須逐一核對，不可直接移植。

本文的實務建議因此是：小人口的預測區間應以卜瓦松最大概似取代 SVD 以校正
位置，並以狀態空間 REML 取代樸素的 $\mathrm{sd}(\Delta\hat\kappa)$ 以校正
寬度，人數十萬以下時再加計 Kalman 濾波的初始值變異數。三者皆無調校參數，
計算成本亦可忽略——$\sigma$ 的估計是 23 維三對角矩陣上的一維最佳化，
較動差法所需的多次完整配適快上兩到三個數量級。

\section{結論與討論}

\subsection{對四個研究問題的回應}

\textbf{問題一：事前診斷能否指出可估性與偏誤的位置。} 能，四項診斷各自
對應一項失效機制且皆不需真值。零死亡格的分布正確指出 $\alpha_x$ 偏誤所在的
年齡，逐年齡偏誤與零格比例的 Spearman 相關在人數一萬時達 0.915；
死亡數的跨年齡落差與噪音變異的跨度指出加權失衡的嚴重程度；
最大奇異值與純噪音上緣的比值指出 $\beta_x$ 是否可還原，
而該比值所預測的門檻與 T-ratio 所顯示的轉折位置一致；
零格比例是否逾 50\% 則界定穩健損失的作用範圍。第\ref{sec:TN}顯示門檻的
位置隨配適期長度移動，因此固定的人口數門檻不如可計算的診斷實用。

\textbf{問題二：三個參數的機制是否不同，對數是否為唯一來源。} 不同，
而對數不是唯一來源。階梯式分解給出的分工是「取對數傷 $\alpha_x$，
不加權傷漂移項」：對真值加入對稱高斯噪音後再作 SVD 的變體 E 完全不涉及
對數轉換、亦無零格，使 $\alpha_x$ 的最大偏誤由 0.743 降至 0.021，
而漂移項偏誤仍達 0.347，甚至略大於標準 LC 的 0.305；反之僅把等權重改為
以死亡數加權的變體 C 使漂移項偏誤由 0.305 降至 0.141。
兩者是獨立的病，也需要獨立的藥，而依表~\ref{tab:penpos} 的分類，
前者要把懲罰加在概似上，後者要改權重。

\textbf{問題三：各估計量的改善幅度、位置與失效條件。} $\alpha_x$ 以 Firth
處理，人數五萬時逐年齡最大偏誤由 0.8658 降至 0.0680，改善 12.7 倍；
人數一萬時標準 LC 把 5 至 9 歲組推高 $+2.33$、卜瓦松最大概似因完全分離而
推低至 $-13.03$，Firth 為 $+0.43$，同時修好兩者。漂移項以資訊加權處理，
人數五萬時卜瓦松最大概似的偏誤為 0.0232 而標準 LC 為 0.3144。
$\beta_x$ 以 Fay--Herriot 收縮處理，人數五萬時相關係數由 0.037 提升至
0.542。三者皆無調校參數。

失效條件有五類。其一，收縮的正當性有上界，固定比例的懲罰在人數三十萬
以上反而損害均方誤差。其二，卜瓦松最大概似在零格眾多時反而比 SVD 更糟，
故「改用卜瓦松路線比較好」並不成立，正確的陳述是「卜瓦松加 Firth 最好」。
其三，截尾類全面失效，且失效的機制可以指名（表~\ref{tab:retain}）。
其四，穩健損失的作用範圍由中位數的崩潰點界定，而第\ref{sec:qlc}進一步
顯示其上界不是調參不足，是單一 $\tau$ 必須同時服務兩個方向相反的機制。
其五，以「年」為單位的穩健化在本文的資料上完全無作用，因為不存在離群年。

參數的改善確實轉移到平均餘命，但轉移的來源與直覺不同。人數五萬時標準 LC
低估配適末年的零歲平均餘命達 1.480 歲、低估十年後的點預測達 2.664 歲，
改用卜瓦松最大概似後分別降為 0.202 與 0.512 歲；考慮到我國縣市間平均餘命的
最大差距約為 10 歲，2.664 歲已足以使一個地區在排序上錯置數個名次。
惟第\ref{sec:e0}的年齡別歸因顯示，30 歲以下各組僅貢獻總誤差的 20.6\%，
而 60 至 89 歲合計占 62.8\%。換言之，Firth 改善最顯著之處並非平均餘命
誤差的主要來源；若標的為平均餘命，資訊加權的重要性更甚於偏誤減少。

\textbf{問題四：區間誤差的來源與「區間過窄」的適用性。} 區間的失準來自
兩個可分離的部分——位置繼承點估計的偏誤，寬度則因樸素估計量把測量噪音
誤認為過程變異而被高估。文獻長期主張的「LC 區間過窄」在小人口\textbf{不
成立}，實測顯示區間同時太寬且涵蓋不足：人數五萬時寬度 8.949、涵蓋率 0.27，
而不可約下限為 2.093、名目水準為 80\%。兩步校正後寬度降至 1.610 而涵蓋率
升至 0.47。


\subsection{實務建議}

綜合上述，本文的建議可歸納為一張以事前診斷為入口的策略表
（表~\ref{tab:boundary}）。使用者取得一筆鄉鎮市區資料後，先計算四項診斷，
再依診斷\cn{4}的比值決定 $\beta_x$ 是否估計、依診斷\cn{1}的最小單年死亡數
中位數決定是否必須採用 Firth。須留意的是，策略表依參數的可估性編排，而使用者關心的多為平均餘命；
第\ref{sec:e0}顯示兩者的年齡權重並不相同，若標的為平均餘命，則即使在
$\beta_x$ 不可估的規模下，改用卜瓦松或加權 SVD 仍可使配適的平均餘命偏誤
由 1.480 歲降至 0.2 至 0.4 歲之間，仍屬值得。這樣的入口之所以比人口數
門檻實用，是因為
取得資料的使用者既無真值亦無從進行模擬，而診斷的比值會隨該地區的年齡
結構、死亡率水準與可用配適期自動調整，也才能解釋為何同為五萬人，有的
地區可估而有的不可。


\subsection{討論}

三項需要補充說明的地方。

第一，本文與既有小樣本研究的關係。Wang et al.（2018）與余清祥等（2021）
指出 LC 的年齡參數在人數減少時出現有方向的偏誤，本文確認此一觀察並進一步
指出偏誤的來源不只一個，因而不能以單一方法補救。既有研究所給的人口數門檻
（五萬與一百萬）在本文的診斷架構下成為一個推論結果而非前提，且該門檻的
位置隨配適期長度移動。

第二，診斷所指的年齡與平均餘命所重的年齡並不相同。診斷\cn{2}指出的是
「對數尺度上偏誤最大的年齡」，而非「對平均餘命影響最大的年齡」。
這並不削弱該診斷的價值，但使用者須依標的量決定要處理哪一個誤差來源。
須留意逐年齡貢獻的合計（標準 LC 為 $-1.3043$ 歲）略小於基準誤差
（$-1.480$ 歲），差額來自生命表的非線性轉換——各年齡的效應並非可加，
故該表應讀為相對重要性的排序而非嚴格的分解。

第三，修勻降低觀測的變異因而使配適誤差變小，卻同時使不確定性被低估
（第\ref{sec:limits}與第\ref{sec:interval}），故「配適誤差較小」不足以
支持一個方法用於生命表編製。


\subsection{研究限制}

第一項限制與數值的可比較性有關。本文的模擬分屬數次獨立執行，其真值的
定義並不一致：階梯式分解與區間的來源分解以 SVD 配適為真值（漂移
$-0.3678$），估計量的比較、分位數變體與狀態空間的結果則以卜瓦松配適
為真值（漂移 $-0.4087$）。各表內部一致，但漂移項偏誤與區間寬度不可
跨表比較，例如表~\ref{tab:idecomp} 與表~\ref{tab:ival} 對「現行作法」
在人數五萬的寬度分別為 5.401 與 8.949，即源於此。統一真值的定義是本文
後續必須完成的工作。

第二，本文所有的「小人口」皆由全國資料模擬而來，隱含小區域與全國有相同
年齡型態的假設。真實鄉鎮市區的型態可能本就不同，此時前三項診斷仍然有效
（它們只描述資料本身），但診斷\cn{4}的最大奇異值會同時反映型態差異與
噪音，解釋須更為謹慎。以真實縣市或鄉鎮市區資料驗證，是本文尚未完成的
一步。

第三，診斷\cn{4}的理論依據建立在白噪音之上，而本文的噪音是異質的，此處
以 $\sqrt{\overline{1/\hat\mu}}$ 作為等效尺度乃是近似。門檻的落點與門檻
以上的還原程度都與實測吻合，但嚴格而言需補一組白化後的對照。

第四，重複次數為 100 次，階梯分解為 400 次。由第伍節兩次獨立模擬的比較
可知，重複 100 次時 $\alpha$ 偏誤的蒙地卡羅誤差約一成、$\beta$ 平方誤差和
約兩成，解讀表中數值時應以此為精度參考；涵蓋率的精度更低，其標準誤約為
0.04，故表~\ref{tab:ival} 中 0.47 與 0.54 的差距僅在兩個標準誤上下。

第五，本文的模擬以卜瓦松生成死亡數，因此過度離散的效應在設計上即不存在，
負二項或準概似方法的效益無從測得，須以真實資料檢驗。同理，暴露數在本文
的設計中不隨時間變動，真實資料的人口結構變遷可能使部分結論須要調整。

第六，Fay-Herriot 收縮假設各年齡的 $\beta_x$ 可交換，但真實的 $\beta_x$
有明確的年齡結構（幼年高、中年低、高齡回升）。目前結果顯示此一假設仍
足以改善估計，更貼切的作法需引入相鄰年齡的相關性，等價於回到差分懲罰。

\subsection{後續工作}\label{sec:future}

最直接的延伸來自第\ref{sec:mq}所留下的缺口。該小節已完成 $(\tau,c)$ 的
掃描，並發現最適的 $\tau$ 在不同年齡朝相反方向移動，分界落在單年期望
死亡數約一人處。順此可推知一個具體的作法：\textbf{令 $\tau$ 隨年齡變動}，
在期望死亡數不足一人的年齡取 $\tau<0.5$ 以抵銷零格替代的向上偏誤，
在其餘年齡取 $\tau>0.5$ 以抵銷 Jensen 的向下偏誤，而分界與方向皆由
診斷\cn{1}的死亡數量級決定，不需真值。此一設計的代價是自由度增加，
須以逐年齡的 $\tau$ 是否可由死亡數量級穩定預測來驗證，否則即為過度配適。
若成立，則「$\tau$ 該取多少」也可接回事前診斷，與本文其餘三項診斷的
定位一致。

另一條尚未嘗試且性質不同的路徑是向上對齊到可靠的總量。
Di Fonzo and Marini（2015）為國民所得統計發展的調節（Reconciliation）
方法，處理的是「低頻的年資料可靠、高頻的季資料含噪音」這類問題，
其成長率保持原則在調整時不破壞序列的成長型態。移到本問題上，
鄉鎮市區的死亡數必須加總為縣市與全國的死亡數，而後者的抽樣誤差可以
忽略；成長率保持所要保持的正是死亡率改善的型態，亦即本文最難估好的
漂移項。此一取徑與第\ref{sec:gap}所列的第一類作法（向參考人口借資訊）
不同：它借的不是另一個人口的趨勢，而是自身聚合後的總量，因而不需要
「參考人口的改善型態可比」這個前提。就鄉鎮市區生命表而言，這可能是
比穩健損失更有效的一條，值得列為優先。

其次是統一真值的定義並提高重複次數，使各表可以跨表比較，特別是涵蓋率
那一組。第三是以真實的縣市與鄉鎮市區資料驗證四項診斷；本文的門檻皆由
模擬得出，而診斷的價值正在於使用者能在自己的資料上計算，若能在若干真實
小區域上比對診斷的判讀與實際估計結果，本文的主張才算完整。最後，第\ref{sec:e0}已把流程接上生命表，
但只檢驗了平均餘命的點估計；平均餘命的預測區間是否同樣改善，須把第陸節
的區間方法與生命表的非線性轉換合併處理，此為使本研究能完整回答
「能否編製鄉鎮市區生命表」的最後一步。另可一併檢驗其他生命表函數，
例如壽命差異 $\edag$ 對高齡死亡率的敏感度與平均餘命不同，其誤差的年齡
來源可能又是另一種分布。

\appendix
\renewcommand{\thesection}{附錄\Alph{section}}
\setcounter{section}{0}

\section{懲罰式估計的實作細節}\label{sec:apppen}

本附錄保留懲罰式 Lee-Carter 模型的模型設定、演算法，以及其餘變體的比較。
正文第肆節只取其中與「在高斯架構內修補的限度」直接相關的結果，其餘細節
列於此處，以便日後需要時取用。

\subsection{模型設定}

在卜瓦松架構下加入 $\beta_x$ 的懲罰項：
\begin{equation}
\min_{\alpha,\beta,\kappa}\ -\ell(\alpha,\beta,\kappa)
+\lambda\sum_x w_x\Big[a\,|\beta_x-\beta_{0x}|
+\tfrac{1-a}{2}(\beta_x-\beta_{0x})^2\Big],
\quad \text{s.t. } \sum_x\beta_x=1,\ \sum_t\kappa_t=0,
\end{equation}
其中 $a=1$ 為 LASSO、$a=0$ 為 Ridge、$0<a<1$ 為 Elastic Net；
$w_x$ 為年齡別權重（可取 $1$ 或自適應權重 $1/\sqrt{D_{x\cdot}}$，
後者對應 Zou (2006) 的 adaptive LASSO）；$\beta_0$ 為收縮目標。

\subsection{演算法與收斂性}

本文以區塊座標下降（Block Coordinate Descent）求解：$\alpha$ 有封閉解，
$\kappa$ 以 Newton 法更新後置中，$\beta$ 則以 proximal Newton
（二次近似加近端算子）更新——給定 Lagrange 乘子 $\nu$，各年齡的解為
\begin{equation}
\beta_x=\beta_{0x}+\frac{S\big(W_x(\tilde\beta_x-\beta_{0x})-\nu,\;
\lambda a w_x\big)}{W_x+\lambda(1-a)w_x},
\end{equation}
其中 $S(z,\gamma)=\mathrm{sign}(z)\max(|z|-\gamma,0)$ 為軟閾值算子，
$W_x=\sum_t\hat\mu_{xt}\kappa_t^2$ 為負二階導，$\nu$ 以二分法求解使
$\sum_x\beta_x=1$。每步並以步長折半確保目標函數不下降。

由於不可微的懲罰項僅作用於 $\beta$ 這一個區塊、且在各年齡間可分離
（separable），Tseng（2001）關於區塊座標下降的收斂定理適用，可保證每個
聚點皆為駐點。實測結果（表~\ref{tab:conv}、圖~\ref{fig:conv}）顯示：目標函數單調上升，
$\log_{10}\max_x|\Delta\beta_x|$ 隨迭代次數線性下降，即\textbf{幾何收斂}；
且\textbf{懲罰強度越大收斂越快}（$\lambda=0$ 需 33 次、$\lambda=10^3$ 需 18 次、
$\lambda=10^4$ 僅需 6 次），因為懲罰項改善了 Hessian 的條件數。

\begin{table}[H]
\centering
\caption{迭代收斂情形（女性 2001$\sim$2024 年，$\tau=10^{-11}$）}
\small
\begin{tabular}{rccc}
\toprule
$\lambda$ & 迭代次數 & 收斂 & 最終 $\max_x|\Delta\beta_x|$\\
\midrule
$0$      & 33 & 是 & $8.0\times10^{-6}$\\
$10^3$   & 18 & 是 & $1.2\times10^{-5}$\\
$10^4$   & \phantom{0}6 & 是 & $5.9\times10^{-7}$\\
\bottomrule
\end{tabular}
\label{tab:conv}
\end{table}

\begin{figure}[H]
\centering
\includegraphics[width=0.88\textwidth]{fig1_convergence.png}
\caption{迭代收斂診斷。左：目標函數與最終值的間隙（對數尺度）；
右：參數變動量。兩者皆呈線性下降，即幾何收斂。}
\label{fig:conv}
\end{figure}

\subsection{Elastic Net 的其他變體}

由於前一小節顯示 LASSO 的弱點在偏誤，本文另外試驗三種以「降低收縮偏誤」
為目的的變體，並與 LASSO、Elastic Net 並列（表~\ref{tab:envar}）：

\begin{itemize}
  \item \textbf{自適應 LASSO}（Zou 2006）：以未懲罰估計的偏離量倒數為權重，
        $w_x=1/|\hat\beta_x^{\text{unpen}}-\beta_{0x}|$，使證據較強的年齡組
        受到較輕的懲罰。
  \item \textbf{鬆弛 LASSO}（Meinshausen 2007）：先以 LASSO 選出作用集，
        再於作用集上不加懲罰重新配適，非作用集固定於 $\beta_0$。
  \item \textbf{MCP}（Zhang 2010）：非凸懲罰，其導數隨 $|\beta_x-\beta_{0x}|$
        增大而遞減至零，因而保留稀疏性卻不對大係數施加偏誤。
\end{itemize}

\begin{table}[H]
\centering
\caption{Elastic Net 族各變體的比較（$N=20$ 萬，重複 100 次，各取其最適 $\lambda$）}
\small
\begin{tabular}{lcccc}
\toprule
方法 & 最適 $\lambda/\lambda_{\max}$ & 偏誤$^2$ & 變異數 & MSE\\
\midrule
Ridge（前表） & 0.003 & 0.00270 & 0.00311 & \textbf{0.00578}\\
LASSO & 0.03 & 0.00219 & 0.00484 & 0.00698\\
MCP（$\gamma=3$） & 0.03 & 0.00219 & 0.00485 & 0.00698\\
Elastic Net & 0.10 & 0.00463 & 0.00255 & 0.00716\\
自適應 LASSO & 0.40 & 0.00313 & 0.00532 & 0.00840\\
鬆弛 LASSO & 0.20 & 0.00734 & \textbf{0.00192} & 0.00924\\
\bottomrule
\end{tabular}
\label{tab:envar}
\end{table}

三項觀察。其一，\textbf{沒有任何一個稀疏型變體能勝過單純的 Ridge}。
這與前一小節的診斷一致：真值 $\beta_x$ 並無恰等於 $1/A$ 者，
問題本質是\textbf{稠密}的，因此所有為稀疏性設計的機制都無從發揮。
其二，\textbf{MCP 與 LASSO 幾乎完全相同}（MSE 差異在第四位小數）。
原因是在最適的 $\lambda$ 之下閾值本就很小，MCP「對大係數不加偏誤」的區段
極少被觸發，退化為 LASSO。其三，\textbf{鬆弛 LASSO 的變異數最低卻 MSE 最差}
——它把非作用集\textbf{硬性固定}於 $\beta_0$，等於施加一個過強的等式限制，
偏誤平方因而升至 0.0073。

\subsection{點估計 LASSO 與分位數 LASSO 的演算法差異}\label{sec:algodiff}

本文在兩處用到 $L_1$ 懲罰，但兩者的最佳化結構\textbf{並不相同}，
可用的演算法也不相同。此一區分在文獻中常被「LASSO」一詞掩蓋，
卻直接決定了實作的可行性，故於此釐清。

\subsubsection{（一）兩種問題的結構}

\textbf{點估計 LASSO}（本文第柒節所用）以平滑的損失搭配 $L_1$ 懲罰：
\begin{equation}\label{eq:pointlasso}
\min_{\boldsymbol{\beta}}\ \underbrace{-\ell(\boldsymbol{\beta})}_{\text{可微}}
+\underbrace{\lambda\textstyle\sum_x w_x|\beta_x-\beta_{0x}|}_{\text{不可微，但逐座標可分離}} .
\end{equation}

\textbf{分位數 LASSO}（即第肆節所述 Rabbi and Mazzuco 的二維修勻）則以檢查函數
$\rho_\tau(u)=u(\tau-\mathbf{1}\{u<0\})$ 為損失：
\begin{equation}\label{eq:qlasso}
\min_{z}\ \underbrace{\textstyle\sum_i\rho_\tau(y_i-(Mz)_i)}_{\text{不可微，且耦合所有座標}}
+\lambda\|Dz\|_{L_1}.
\end{equation}

兩者的關鍵差別在於\textbf{不可微的部分落在哪裡}。
式~\eqref{eq:pointlasso} 的不可微項是懲罰，逐座標分離；
式~\eqref{eq:qlasso} 的不可微項是\textbf{損失本身}，而
$\rho_\tau(y_i-(Mz)_i)$ 把所有座標耦合在內積 $(Mz)_i$ 之中，無法分離。

\subsubsection{（二）Tseng（2001）的條件與其違反}

Tseng（2001）的收斂定理要求目標函式可寫成
\begin{equation}\label{eq:tseng}
f(x_1,\dots,x_N)=f_0(x_1,\dots,x_N)+\sum_{k=1}^{N} f_k(x_k),
\end{equation}
其中 $f_0$ 可微、不可微的 $f_k$ 逐區塊可分離。該文並明白指出：

\begin{quote}\small
「若 $f$ 不可微，座標下降法\textbf{可能卡在非駐點，即使 $f$ 為凸}……
因此一般認為該方法不適用於不可微的情形。
\textbf{惟當 $f$ 的不可微部分可分離時則為例外。}」
（Tseng 2001: 475，本文中譯）
\end{quote}

對照可知：式~\eqref{eq:pointlasso} 符合式~\eqref{eq:tseng} 的形式，
故第柒節之三的區塊座標下降有收斂保證；
式~\eqref{eq:qlasso} 則\textbf{不符合}，座標下降並無保證。

\subsubsection{（三）以實際資料驗證卡滯}

上述為理論陳述，本文另以實際資料檢驗其是否真的發生。取女性
2015$\sim$2024 年（$22\times10=220$ 格）建構修勻問題，
對\textbf{同一個}目標函式式~\eqref{eq:qlasso} 施以兩種演算法，
結果如表~\ref{tab:algo}。

為避免「卡滯只是步長設定不良」的疑慮，此處的座標下降採用
\textbf{精確的單座標最小化}——$L_1$ 損失沿單一座標的最小值為加權中位數，
有封閉解，不涉及任何步長選擇。

\begin{table}[H]
\centering
\caption{同一目標函式下兩種演算法的比較（女性 2015$\sim$2024 年）}
\small
\begin{tabular}{llrcl}
\toprule
目標函式 & 演算法 & 目標值 & 是否達最適 & 診斷\\
\midrule
分位數 LASSO & 線性規劃（內點法） & \textbf{30.02} & 是 & ---\\
分位數 LASSO & 座標下降（精確） & \textbf{1244.20} & \textbf{否} & 可改善的單座標移動 $0/220$\\
\midrule
點估計 LASSO & 座標下降（軟閾值） & --- & 是 & KKT 最大違反 $7.1\times10^{-16}$\\
\bottomrule
\end{tabular}
\label{tab:algo}
\end{table}

座標下降在\textbf{第一次掃描後即停止，且目標值等於起始值}——換言之
完全沒有移動，而該點的目標值是線性規劃解的 \textbf{41.4 倍}。

\textbf{卡滯的機制可以明確指認。} 在起始點 $z=\mathbf{0}$ 處，
沿第 $j$ 個座標移動時，保真列貢獻權重 $1$、懲罰列貢獻權重
$\sum|D_{\cdot j}|$；本設定下後者的中位數為 $21.60$、最小值 $5.40$，
\textbf{全部 220 個座標的懲罰權重都大於保真權重}。
由於 $L_1$ 的單座標最小值為加權中位數，權重集中於「不動」的一側時
最小值即為 $0$，故\textbf{每個座標都同意不要移動}。
然而聯合移動可使目標值由 1244.20 降至 30.02——
此即典型的「對角山谷」：沿任一軸皆上坡，沿對角線則下坡。
本文逐一檢驗 220 個座標的八種位移，確認\textbf{無任何單座標移動可改善}，
證實此為結構性的卡滯而非實作缺陷。

對照之下，同一組資料改用平方損失搭配 $L_1$ 懲罰，座標下降
\textbf{兩次掃描即收斂}，KKT 條件的最大違反為 $7.1\times10^{-16}$，
確實達到最適。\textbf{兩者的差別純粹來自損失函數是否可微。}

\subsubsection{（四）對本文實作的意涵}

上述分析解釋了本文兩個模組為何採用不同的演算法，且此一選擇是
\textbf{結構決定的，而非實作偏好}：

\begin{table}[H]
\centering
\caption{本文兩個模組的目標函式與演算法}
\small
\begin{tabular}{lll}
\toprule
模組 & 目標函式 & 演算法\\
\midrule
\texttt{smooth\_l1\_2d()} & $L_1$ 損失 ＋ $L_1$ 懲罰 &
  \texttt{rq.fit.sfn}（Frisch--Newton 內點法）\\
\texttt{lc\_penalized()} & 卜瓦松對數概似 ＋ $L_1$ 懲罰 &
  proximal Newton ＋ 軟閾值\\
\bottomrule
\end{tabular}
\label{tab:modules}
\end{table}

三項延伸意涵值得記錄。

\textbf{其一，第柒節之三的收斂論證因此有了對照組。} 該處主張
「不可微的懲罰項僅作用於 $\boldsymbol{\beta}$ 這一個區塊、且在各年齡間
可分離，故 Tseng（2001）的定理適用」；表~\ref{tab:algo} 顯示，
若把損失換成 $L_1$，同一套演算法便會失效。此一對照使該主張
不再只是引用定理，而有了具體的反面證據。

\textbf{其二，$L_1$ 損失的選擇有其統計理由，非僅為計算方便。}
$\tau=0.5$ 的檢查函數等價於 Laplace 分配的最大概似，對離群值較平方損失
穩健；Schuette（1978）即以此為由主張以線性規劃處理修勻問題，
Dokumentov et al.（2018）沿用之，Rabbi and Mazzuco（2021）亦引述此點。
換言之，穩健性與「必須使用線性規劃」是同一個選擇的兩面。

\textbf{其三，解的性質不同。} 線性規劃的解落在可行域的頂點上，
意味著恰有若干筆殘差為零（曲面「內插」這些觀測值）；
平方損失的解則一般不內插任何一點。此一差異在解釋修勻後曲面的行為時
須留意，亦是 $L_1$ 修勻「較不平滑但誤差較小」（表~\ref{tab:smooth}）的
成因之一。


\section{懲罰前後 $\beta_x$ 的逐年齡比較}

本附錄逐年齡呈現加懲罰與不加懲罰的 $\beta_x$ 估計值差異。

\subsection{實際資料}

表~\ref{tab:appreal} 為女性 2001$\sim$2024 年的估計結果
（LASSO 與 Elastic Net 取 $\lambda/\lambda_{\max}=0.10$、Ridge 取 0.003）。
虛線水準 $1/A=0.0455$ 為收縮目標。

\begin{table}[H]
\centering
\caption{各年齡組的 $\beta_x$ 估計值（女性 2001$\sim$2024 年）}
\scriptsize
\begin{tabular}{lrrrrrrc}
\toprule
年齡組 & 未懲罰 & LASSO & Elastic Net & Ridge
& LASSO $-$ 未懲罰 & Ridge $-$ 未懲罰 & 併入 $1/A$\\
\midrule
0      & 0.0496 & 0.0455 & 0.0455 & 0.0502 & $-0.0042$ & $+0.0005$ & 是\\
1--4   & \textbf{0.1066} & 0.0455 & 0.0611 & 0.0661 & $\mathbf{-0.0612}$ & $-0.0406$ & 是\\
5--9   & 0.0698 & 0.0455 & 0.0455 & 0.0540 & $-0.0243$ & $-0.0158$ & 是\\
10--14 & 0.0443 & 0.0455 & 0.0455 & 0.0485 & $+0.0012$ & $+0.0042$ & 是\\
15--19 & 0.0578 & 0.0455 & 0.0455 & 0.0531 & $-0.0124$ & $-0.0047$ & 是\\
20--24 & 0.0440 & 0.0455 & 0.0455 & 0.0473 & $+0.0015$ & $+0.0032$ & 是\\
25--29 & 0.0484 & 0.0455 & 0.0455 & 0.0494 & $-0.0029$ & $+0.0010$ & 是\\
30--34 & 0.0453 & 0.0455 & 0.0455 & 0.0474 & $+0.0002$ & $+0.0021$ & 是\\
35--39 & 0.0394 & 0.0455 & 0.0455 & 0.0428 & $+0.0061$ & $+0.0034$ & 是\\
40--44 & 0.0328 & 0.0455 & 0.0455 & 0.0368 & $+0.0127$ & $+0.0040$ & 是\\
45--49 & 0.0261 & 0.0387 & 0.0341 & 0.0302 & $+0.0126$ & $+0.0042$ & \\
50--54 & 0.0290 & 0.0389 & 0.0348 & 0.0321 & $+0.0100$ & $+0.0031$ & \\
55--59 & 0.0442 & 0.0455 & 0.0455 & 0.0457 & $+0.0013$ & $+0.0015$ & 是\\
60--64 & 0.0561 & 0.0546 & 0.0570 & 0.0568 & $-0.0015$ & $+0.0006$ & \\
65--69 & 0.0686 & 0.0687 & 0.0691 & 0.0689 & $+0.0001$ & $+0.0003$ & \\
70--74 & 0.0668 & 0.0675 & 0.0674 & 0.0675 & $+0.0006$ & $+0.0006$ & \\
75--79 & 0.0528 & 0.0531 & 0.0537 & 0.0537 & $+0.0004$ & $+0.0010$ & \\
80--84 & 0.0439 & 0.0455 & 0.0455 & 0.0451 & $+0.0016$ & $+0.0012$ & 是\\
85--89 & 0.0319 & 0.0357 & 0.0333 & 0.0333 & $+0.0038$ & $+0.0014$ & \\
90--94 & 0.0215 & 0.0264 & 0.0224 & 0.0236 & $+0.0049$ & $+0.0022$ & \\
95--99 & 0.0135 & 0.0255 & 0.0197 & 0.0194 & $+0.0121$ & $+0.0059$ & \\
100+   & \textbf{0.0078} & 0.0455 & 0.0422 & 0.0284 & $\mathbf{+0.0376}$ & $+0.0206$ & 是\\
\bottomrule
\end{tabular}
\label{tab:appreal}
\end{table}

值得注意的是，\textbf{1--4 歲的未懲罰估計值最大（0.1066）卻被完全併入
$1/A$}。這並非矛盾：LASSO 的軟閾值為 $\lambda w_x/W_x$，而
$W_x=\sum_t\hat\mu_{xt}\kappa_t^2$ 正是該年齡的資訊量。幼兒死亡數極少使
$W_x$ 很小，因此其估計值雖大卻不精確，反而受到最重的懲罰。
\textbf{懲罰的強弱取決於資訊量而非係數大小}，此點與一般迴歸中
LASSO 的行為一致，但在死亡率模型上格外明顯。

存活下來（未被併入 $1/A$）的是 45--54 歲與 60--99 歲，
亦即\textbf{高齡段的改善型態確實顯著偏離平均}，這與人口學的認知相符。

\begin{figure}[H]
\centering
\includegraphics[width=\textwidth]{figA_beta_by_age.png}
\caption{懲罰前後 $\beta_x$ 的逐年齡比較。左上為實際資料；右上與左下為模擬的
平均估計值（含真值）；右下為逐年齡的標準差。}
\label{fig:appbeta}
\end{figure}

\subsection{模擬（已知真值）}

在模擬中可進一步把差異拆成偏誤與標準差（表~\ref{tab:appsim}）。
$N=5$ 萬、$\lambda/\lambda_{\max}=0.20$ 時，未懲罰版的平均絕對偏誤為
0.0069、平均標準差 0.0679；LASSO 則為 0.0123 與 0.0078。
亦即\textbf{標準差的中位數降幅達 75\%，與絕對偏誤增為 1.8 倍之間取得平衡}。

\begin{table}[H]
\centering
\caption{偏誤最大的五個年齡組（模擬，$N=5$ 萬，重複 100 次）}
\small
\begin{tabular}{lrrrrrr}
\toprule
年齡組 & 真值 & 未懲罰平均 & LASSO 平均 & LASSO 偏誤
& 未懲罰 SD & LASSO SD\\
\midrule
1--4   & 0.1107 & 0.1107 & 0.0461 & $-0.0646$ & 0.1666 & 0.0060\\
100+   & 0.0058 & $-0.0012$ & 0.0455 & $+0.0397$ & 0.1008 & 0.0075\\
5--9   & 0.0785 & 0.0924 & 0.0455 & $-0.0330$ & 0.1899 & 0.0000\\
95--99 & 0.0127 & 0.0078 & 0.0366 & $+0.0239$ & 0.0288 & 0.0123\\
45--49 & 0.0260 & 0.0286 & 0.0417 & $+0.0157$ & 0.0321 & 0.0097\\
\bottomrule
\end{tabular}
\label{tab:appsim}
\end{table}

\textbf{懲罰造成的偏誤集中在兩端年齡}（1--9 歲與 95 歲以上），
正是資訊量最稀薄之處；而這些年齡的未懲罰標準差高達 0.10 至 0.19，
本就不具參考價值。換言之，懲罰在這些年齡是以「明確的偏誤」取代
「無法使用的變異」，是否值得須視用途而定：
若後續要計算平均餘命（對 $\beta_x$ 的加權以死亡數為主），
兩端年齡的偏誤影響有限；若要直接解讀各年齡的改善速度，則須特別註明。

\vspace{10pt}
\section{移入附錄的數值結果}\label{sec:appmoved}

本附錄收錄主文因篇幅而未列出的數值表，依主文引用的順序排列。各表的
討論見主文對應各節，此處僅列出結果。

\subsection{配適期長度與人口規模的二因子設計}\label{sec:TN}

30、40、55 年，$N$ 取兩萬、五萬、二十萬與一百萬，各重複 100 次。

評估指標改用 $\hat\beta_x$ 與真值的相關係數而非誤差平方和，原因是在配適期
與人口規模都很小時，時間參數的變化範圍不足以識別 $\beta_x$，估計值會完全
退化為均勻向量 $\mathbf{1}/A$；此時誤差平方和恰為
$\|\mathbf{1}/A-\beta^{\text{true}}\|^2=0.0070$，反而小於配適期稍長但估計值
劇烈震盪時的數值，若只看誤差平方和會誤判為「短配適期較佳」。

\begin{table}[H]
\centering
\caption{$\hat\beta_x$ 與真值的相關係數中位數（未懲罰，重複 100 次）}
\small
\begin{tabular}{rcccc}
\toprule
$T$（年） & $N=2\times10^4$ & $N=5\times10^4$ & $N=2\times10^5$ & $N=10^6$\\
\midrule
10 & 0.247 & 0.156 & 0.170 & 0.120\\
15 & 0.188 & 0.191 & 0.226 & 0.465\\
20 & 0.178 & 0.217 & 0.407 & 0.760\\
30 & 0.323 & 0.593 & 0.785 & \textbf{0.952}\\
40 & 0.508 & 0.725 & \textbf{0.899} & \textbf{0.977}\\
55 & 0.746 & \textbf{0.874} & \textbf{0.966} & \textbf{0.991}\\
\bottomrule
\end{tabular}
\label{tab:TNcor}
\end{table}
\begin{table}[H]
\centering
\caption{退化率：$\hat\beta_x$ 完全塌陷為 $\mathbf{1}/A$ 的比例}
\small
\begin{tabular}{rcccc}
\toprule
$T$（年） & $N=2\times10^4$ & $N=5\times10^4$ & $N=2\times10^5$ & $N=10^6$\\
\midrule
10 & \textbf{0.81} & 0.27 & 0 & 0\\
15 & \textbf{0.53} & 0.04 & 0 & 0\\
20 & 0.22 & 0 & 0 & 0\\
$\ge30$ & 0 & 0 & 0 & 0\\
\bottomrule
\end{tabular}
\label{tab:TNdegen}
\end{table}

結果支持第三種情境。配適期十年時，即使人口達一百萬，相關係數仍僅 0.120，
而配適期五十五年時人口僅兩萬亦達 0.746；配適期二十年以下時，沒有任何
人口規模能達到 0.8 的水準。換言之，配適期的作用不只是提供更多資料，
增加人口無法補償它的不足。表~\ref{tab:TNdegen} 的退化率則說明了短配適期
失效的形式：配適期十年、人口兩萬時有 81\% 的模擬中 $\hat\beta_x$ 完全塌陷
為均勻向量，亦即估計程序根本沒有從資料中取得任何年齡型態的資訊。

此一結果可由第\ref{sec:bbp}的相變理論解釋。訊號強度隨配適期成長的
速度遠快於門檻，而噪音尺度僅隨人口規模的平方根下降，因此等準確度曲線在
兩個因子上並不對稱——這也說明了為何本文第四項診斷的比值必須逐筆資料
計算，而不能以單一的人口數門檻代替。

\subsection{修勻與壽命差異調整}

故目標函數為分段線性，等價於 $\tau=0.5$ 的分位數迴歸而須以線性規劃求解。
第二部分是在估計之後重新求解時間參數，使配適的壽命差異等於觀測值，
取代 Lee and Carter（1992）對齊總死亡數的作法。在女性 2001 至 2024 年
的資料上，該方法確實如原文所述具有較小的配適誤差，表~\ref{tab:smooth} 顯示
$L_1$ 修勻的平均絕對誤差為 0.0413，低於一維與二維樣條的 0.0538 與 0.0759。

\begin{table}[H]
\centering
\caption{三種修勻法的比較（女性 2001$\sim$2024 年）}
\small
\begin{tabular}{lcccc}
\toprule
修勻法 & MAE & MSE & 年齡方向粗糙度 & 時間方向粗糙度\\
\midrule
觀測值 & --- & --- & 0.0204 & 0.0981\\
一維樣條（$L_2$，年齡） & 0.0538 & 0.0218 & 0.0140 & 0.0908\\
二維樣條（$L_2$） & 0.0759 & 0.0247 & \textbf{0.0137} & 0.0025\\
LASSO（$L_1$） & \textbf{0.0413} & \textbf{0.0047} & 0.0179 & \textbf{0.0012}\\
\bottomrule
\end{tabular}
\label{tab:smooth}
\end{table}

然而把比較擴充為性別與配適起始年構成的十二格面板之後，此一優勢並不穩固。
表~\ref{tab:panel} 顯示 $\LCe$ 僅在五格取得最佳的樣本外平均絕對誤差，
其餘七格由 Lee-Miller 奪冠，而且在長配適期之下無一例外。此一型態與
Booth et al.（2005）的結論一致：替代調整方式的效益不僅有限，且會隨可用
資訊增加而消失。

\begin{table}[H]
\centering
\caption{各面板格的樣本外 MAE（$h=5$）；粗體為該格最佳}
\small
\begin{tabular}{lcccc@{\hspace{14pt}}lcccc}
\toprule
起始年 & $\LCe$ & M$\LCe$ & LM & LC & 起始年 & $\LCe$ & M$\LCe$ & LM & LC\\
\midrule
女 2001 & \textbf{.0962} & \textbf{.0962} & .1080 & .1194 & 男 1981 & .1323 & --- & \textbf{.1099} & .1454\\
女 1991 & .0992 & \textbf{.0988} & .1064 & .1197 & 男 1971 & .1124 & --- & \textbf{.1078} & .1353\\
女 1981 & .1194 & .1202 & \textbf{.1057} & .1163 & 總 2001 & .0862 & \textbf{.0841} & .0967 & .1141\\
女 1971 & .1172 & .1152 & \textbf{.1061} & .1300 & 總 1991 & .0945 & \textbf{.0905} & .0956 & .1266\\
男 2001 & .0965 & \textbf{.0942} & .1062 & .1232 & 總 1981 & .1080 & --- & \textbf{.0950} & .1268\\
男 1991 & .1243 & --- & \textbf{.1090} & .1436 & 總 1971 & .1005 & --- & \textbf{.0935} & .1227\\
\bottomrule
\end{tabular}
\label{tab:panel}
\end{table}

更值得注意的是預測區間。表~\ref{tab:e0} 顯示 $\LCe$ 的 2050 年零歲平均餘命
預測區間是六種方法中最寬者。修勻的直觀效果是壓低歷史變異、使區間變窄，
此處卻反其道而行，原因在於壽命差異調整重新求解時間參數之後，其年際變動
反而增大。這一點與 Basellini et al.（2023）所提出、卻未量化的疑慮方向相反，
本文第陸節將以變體分解說明兩者何以並存。

\begin{table}[H]
\centering
\caption{2050 年女性零歲平均餘命的預測與 80\% 預測區間}
\small
\begin{tabular}{lccc}
\toprule
方法 & $e_0$（2050） & 80\% 預測區間 & 區間寬度\\
\midrule
LC      & 87.99 & $[86.10,\ 89.69]$ & 3.591\\
LCP     & 88.02 & $[86.53,\ 89.40]$ & \textbf{2.875}\\
LM      & 87.86 & $[86.25,\ 89.34]$ & 3.090\\
BMS     & 87.95 & $[86.45,\ 89.34]$ & 2.895\\
$\LCe$  & 87.61 & $[85.51,\ 89.49]$ & \textbf{3.977}\\
M$\LCe$ & 87.47 & $[85.45,\ 89.29]$ & 3.842\\
\bottomrule
\end{tabular}
\label{tab:e0}
\end{table}

此處另有一項與本文議題三直接相關的觀察。原文以各修勻法對觀測值的誤差作為
選擇準則，而修勻依定義必然增加對觀測值的誤差，因此在臺灣資料上，該準則
必然傾向選擇「不修勻」。換言之，這個準則衡量的是配適誤差而非參數的估計
誤差，兩者在小人口下並不同向。

\subsection{懲罰式估計的數值結果}

若取收縮目標為零向量（即一般 LASSO 的預設），則在 $\hat\beta_x$ 全為正
（人口規模較大時必然如此）的情形下 $\sum_x|\beta_x|\equiv1$ 為常數，
懲罰項不影響最適解，估計值完全不動。表~\ref{tab:shrink} 顯示 $\lambda$ 由
0 加到 $10^6$，估計值的變動量始終在機器精度之內。

\begin{table}[H]
\centering
\caption{收縮目標的選擇（女性 2001$\sim$2024 年）}
\small
\begin{tabular}{rcccc}
\toprule
& \multicolumn{2}{c}{$\beta_0=\mathbf{0}$} & \multicolumn{2}{c}{$\beta_0=\mathbf{1}/A$}\\
\cmidrule(lr){2-3}\cmidrule(lr){4-5}
$\lambda$ & $\max|\hat\beta-\hat\beta^{\text{unpen}}|$ & 自由參數 &
            $\max|\hat\beta-\hat\beta^{\text{unpen}}|$ & 自由參數\\
\midrule
$0$      & --- & 22 & --- & 22\\
$10^2$   & $9.0\times10^{-17}$ & 22 & 0.0041 & 19\\
$10^4$   & $1.5\times10^{-16}$ & 22 & 0.0612 & 6\\
$10^6$   & $2.1\times10^{-15}$ & 22 & 0.0612 & 0\\
\bottomrule
\end{tabular}
\label{tab:shrink}
\end{table}

本文因此改取均勻向量 $\mathbf{1}/A$ 為收縮目標。此一選擇有三項優點：
$\sum_x(\beta_x-\beta_{0x})=0$ 使懲罰項有效；具明確的人口學意義，即
「各年齡以相同速度改善」，且當 $\lambda\to\infty$ 時解收斂至 $\mathbf{1}/A$；
而 LASSO 的稀疏性在此回答的是「哪些年齡組的改善速度顯著偏離平均」，
正是指導方向所期待的解釋效果。

收縮目標修正之後，懲罰確實有效。表~\ref{tab:bv} 與表~\ref{tab:rate} 顯示
人數五萬時 $\beta$ 的均方誤差可降低 94\%，符合第參節「小樣本下 $\beta_x$
變異數主導」的診斷。但同一組結果亦顯示了限制所在：變異數隨人口規模下降的
速度快於偏誤，兩者在人數三十萬附近交叉，此後固定比例的懲罰反而使均方誤差
上升。這與第參節 T-ratio 的非單調型態是同一件事的兩種表現——收縮的正當性
僅存在於變異數主導的區間，而該區間的上界可由診斷事前判定。

\begin{table}[H]
\centering
\caption{$\beta$ 的偏誤—變異數分解（LASSO，重複 100 次）}
\small
\begin{tabular}{rcccc}
\toprule
$N$ & $\lambda/\lambda_{\max}$ & 偏誤$^2$ & 變異數 & MSE\\
\midrule
50,000  & 0（未懲罰） & 0.00239 & 0.15980 & 0.16058\\
        & 0.10 & 0.00648 & 0.00512 & 0.01155\\
        & \textbf{0.20} & 0.00854 & 0.00158 & \textbf{0.01010}\\
        & 1.00 & 0.01066 & $2.8\times10^{-6}$ & 0.01066\\
\midrule
200,000 & 0（未懲罰） & 0.00024 & 0.01691 & 0.01698\\
        & 0.01 & 0.00015 & 0.01062 & 0.01066\\
        & \textbf{0.03} & 0.00147 & 0.00560 & \textbf{0.00702}\\
        & 0.10 & 0.00707 & 0.00102 & 0.00808\\
\bottomrule
\end{tabular}
\label{tab:bv}
\end{table}
\begin{table}[H]
\centering
\caption{變異數與 MSE 隨暴露數的變化（重複 100 次）}
\small
\begin{tabular}{rcccc}
\toprule
 & \multicolumn{2}{c}{未懲罰} & \multicolumn{2}{c}{LASSO ($\lambda/\lambda_{\max}=0.1$)}\\
\cmidrule(lr){2-3}\cmidrule(lr){4-5}
$N$ & 變異數 & MSE & 變異數 & MSE\\
\midrule
20,000    & 4.3912 & 4.3860 & \textbf{0.0223} & \textbf{0.0286}\\
50,000    & 0.2139 & 0.2148 & \textbf{0.0060} & \textbf{0.0119}\\
100,000   & 0.0526 & 0.0530 & \textbf{0.0021} & \textbf{0.0089}\\
200,000   & 0.0184 & 0.0184 & \textbf{0.0009} & \textbf{0.0081}\\
500,000   & 0.0066 & \textbf{0.0066} & 0.0003 & 0.0078\\
1,000,000 & 0.0030 & \textbf{0.0030} & 0.0001 & 0.0077\\
\midrule
\multicolumn{1}{l}{log-log 斜率} & $-1.77$ & & $-1.27$ & \\
\bottomrule
\end{tabular}
\label{tab:rate}
\end{table}

若純就均方誤差而論，三種懲罰之中勝出的是 Ridge 而非 LASSO
（表~\ref{tab:pen}），亦即稀疏性並非此處正確的先驗。LASSO 的價值在於
可解釋性：它以約兩成的均方誤差，換取「哪些年齡組顯著偏離平均改善速度」
的明確答案。自適應 LASSO、鬆弛 LASSO 與 MCP 三種變體則全數輸給單純的
Ridge，其細節列於附錄。

\begin{table}[H]
\centering
\caption{三種懲罰的比較（$N=20$ 萬，重複 100 次，各取其最適 $\lambda$）}
\small
\begin{tabular}{lcccccc}
\toprule
懲罰 & $a$ & 最適 $\lambda/\lambda_{\max}$ & 偏誤$^2$ & 變異數 & MSE & $\beta$ 相關中位數\\
\midrule
未懲罰 & --- & 0 & 0.00024 & 0.01691 & 0.01698 & 0.711\\
LASSO & 1.0 & 0.030 & 0.00147 & 0.00560 & 0.00702 & 0.703\\
Elastic Net & 0.5 & 0.030 & 0.00158 & 0.00534 & 0.00687 & 0.705\\
Ridge & 0.0 & 0.003 & 0.00270 & 0.00311 & \textbf{0.00578} & \textbf{0.721}\\
\bottomrule
\end{tabular}
\label{tab:pen}
\end{table}

\subsection{截尾類估計量與 HeteroPCA}

無重尾，沒有任何離群值，卻仍產生 0.347 的漂移項偏誤。既然沒有東西可截，
截尾自然無從發揮。

本文以最大截尾概似與最小截尾平方兩種代表性方法驗證此一預期
（表~\ref{tab:robust}），並同時納入 HeteroPCA 作為對照。

\begin{table}[H]
\centering
\caption{穩健估計的兩側：截尾類與 HeteroPCA（重複 100 次，取中位數）}
\small
\begin{tabular}{llrrrrr}
\toprule
$N$ & 估計量 & $\alpha$ 中位 & $\alpha$ 最大 & SSE($\beta$) & cor($\beta$) & 漂移偏誤\\
\midrule
\multirow{7}{*}{\shortstack[l]{$5\times10^4$\\ 零格 102/528}}
 & 標準 LC        & 0.0700 & 0.8658 & 3.538 & 0.040 & 0.3144\\
 & 卜瓦松 MLE     & 0.0119 & 0.2053 & 2.166 & 0.298 & 0.0232\\
 & Firth          & \textbf{0.0060} & \textbf{0.0882} & 1.742 & 0.310 & \textbf{0.0180}\\
 & HeteroPCA      & 0.0700 & 0.8658 & 1.992 & 0.093 & 0.2247\\
 & HeteroPCA $+$ FH & 0.0700 & 0.8658 & \textbf{0.355} & 0.273 & 0.2247\\
 & MTL（保留 90\%） & 0.0276 & 0.4712 & 4.082 & 0.222 & 0.0192\\
 & LTS（保留 90\%） & 0.0190 & 0.7871 & 4.311 & 0.041 & 0.3055\\
\midrule
\multirow{7}{*}{\shortstack[l]{$2\times10^5$\\ 零格 29/528}}
 & 標準 LC        & 0.0148 & 0.1671 & 1.341 & 0.390 & 0.2124\\
 & 卜瓦松 MLE     & 0.0034 & 0.0367 & 0.313 & \textbf{0.615} & \textbf{0.0015}\\
 & Firth          & \textbf{0.0013} & \textbf{0.0271} & 0.299 & \textbf{0.619} & 0.0016\\
 & HeteroPCA      & 0.0148 & 0.1671 & 0.632 & 0.416 & 0.1100\\
 & HeteroPCA $+$ FH & 0.0148 & 0.1671 & \textbf{0.155} & 0.551 & 0.1100\\
 & MTL（保留 90\%） & 0.0066 & 0.0836 & 0.377 & 0.603 & 0.0016\\
 & LTS（保留 90\%） & 0.0081 & 0.2239 & 1.243 & 0.389 & 0.2124\\
\bottomrule
\end{tabular}
\label{tab:robust}
\end{table}

結果有三。其一，沒有任何截尾變體勝過未截尾的卜瓦松最大概似。其二，最小
截尾平方完全不處理漂移項的偏誤，人數五萬時為 0.3055、標準 LC 為 0.3144，
人數二十萬時兩者皆為 0.2124，數值幾乎相同；它確實改善了 $\alpha$（因為
截掉了零格替代所造成的極端對數值），但 $\beta$ 的相關係數 0.041 與標準 LC
的 0.040 無異。這正是本小節開頭所預期的：截尾與「異質變異噪音旋轉主奇異
向量」這個機制正交。其三，最大截尾概似的表現接近卜瓦松最大概似，因為它
本就建立在卜瓦松概似之上，但每一項指標都略差。

\subsection{分位數與 M-分位數的完整掃描}

求解採交替最小化。固定 $\boldsymbol{\kappa}$ 時，各年齡各自是一條簡單的
分位數迴歸，以線性規劃求解；固定 $\boldsymbol{\alpha}$ 與
$\boldsymbol{\beta}$ 時，各年份各自是一維的加權分位數問題，$\tau=0.5$ 時
解為 $\{(\log m_{x,t}-\alpha_x)/\beta_x\}$ 以 $w_{x,t}|\beta_x|$ 為權重的
加權中位數。兩步皆為各自區塊的全域最小，故目標函數單調不增，實際收斂
約需六至七次迭代。模擬的亂數種子與表~\ref{tab:est} 完全相同，故兩表可
逐格對照；表~\ref{tab:qlc} 中標準 LC 一列即為內部對照，其數值與
表~\ref{tab:est} 相同。

\begin{table}[H]
\centering
\caption{分位數 Lee-Carter（重複 100 次，種子與表~\ref{tab:est} 相同）}
\small
\begin{tabular}{llrrrrr}
\toprule
$N$ & 估計量 & $\alpha$ 中位 & $\alpha$ 最大 & SSE($\beta$) & cor($\beta$) & 漂移偏誤\\
\midrule
\multirow{3}{*}{$10^4$}
 & A\quad 標準 LC（對照） & 0.1612 & 2.3309 & 7.916 & $-0.074$ & 0.3436\\
 & F\quad 中位數 LC       & 0.0965 & 2.3020 & 10.864 & $-0.077$ & \textbf{0.3671}\\
 & G\quad 加權中位數 LC   & \textbf{0.0733} & 2.3020 & 4.950 & 0.023 & 0.2986\\
\midrule
\multirow{3}{*}{$5\times10^4$}
 & A\quad 標準 LC         & 0.0611 & 0.8658 & 4.449 & 0.037 & 0.3339\\
 & F\quad 中位數 LC       & 0.0284 & 0.6925 & 5.349 & 0.013 & \textbf{0.3110}\\
 & G\quad 加權中位數 LC   & \textbf{0.0187} & 0.6925 & 1.548 & 0.084 & \textbf{0.1591}\\
\midrule
\multirow{3}{*}{$2\times10^5$}
 & A\quad 標準 LC         & 0.0231 & 0.1846 & 1.400 & 0.379 & 0.2117\\
 & F\quad 中位數 LC       & 0.0048 & 0.1780 & 1.055 & 0.443 & \textbf{0.0696}\\
 & G\quad 加權中位數 LC   & \textbf{0.0032} & 0.1818 & \textbf{0.536} & 0.435 & \textbf{0.0428}\\
\bottomrule
\end{tabular}
\label{tab:qlc}
\end{table}

預測的前半段成立。變體 F 在三個規模下都使 $\alpha$ 的偏誤中位數約降一半
至數倍，人數二十萬時由 0.0231 降至 0.0048，改善近五倍。預測的後半段則
只在小人口成立：人數一萬時 F 的漂移項偏誤為 0.3671，反而略大於標準 LC 的
0.3436；人數五萬時為 0.3110 對 0.3339，改善僅 7\%。但到了人數二十萬，
F 的漂移項偏誤降至 0.0696，是標準 LC 的三分之一——這已不能說是「沒有
改善」。

\begin{table}[H]
\centering
\caption{零死亡格比例隨人口規模的變化（各年齡，重複 200 次）}
\small
\begin{tabular}{lrrrr}
\toprule
年齡組 & $\beta_x$ 真值 & $N=10^4$ & $N=5\times10^4$ & $N=2\times10^5$\\
\midrule
0       & 0.0496 & 0.79 & 0.33 & 0.01\\
1--4    & \textbf{0.1066} & 0.93 & 0.74 & \textbf{0.30}\\
5--9    & \textbf{0.0697} & 0.95 & 0.77 & \textbf{0.35}\\
10--14  & 0.0442 & 0.95 & 0.77 & \textbf{0.36}\\
15--19  & 0.0578 & 0.90 & 0.60 & 0.14\\
50--54  & 0.0290 & 0.14 & 0.00 & 0.00\\
80--84  & 0.0439 & 0.00 & 0.00 & 0.00\\
\bottomrule
\end{tabular}
\label{tab:zerorate}
\end{table}

此一看似不一致的型態，可由表~\ref{tab:zerorate} 得到完整的解釋，而解釋的
關鍵正是中位數的崩潰點。零格替代所造成的扭曲並非隨機散布，而是集中於
幼年組，且其比例隨人口規模急遽變動：人數二十萬時 1 至 4 歲組有 30\% 的
年份死亡數為零，人數五萬時升至 74\%，人數一萬時達 93\%。中位數的崩潰點
為 50\%，亦即當被扭曲的觀測超過半數時，中位數本身就是一個被扭曲的值，
穩健性不復存在。

三個規模的結果恰好落在崩潰點的兩側，因而可逐一對應。人數二十萬時扭曲
比例為 30\% 至 36\%，低於崩潰點，中位數得以忽略它們，5 至 9 歲組的
$\alpha$ 偏誤由 $-0.0330$ 降至 $-0.0006$，幾乎完全消除；由於幼年組正是
$\beta_x$ 最大的年齡（1 至 4 歲組為 0.1066），其對時間參數的影響最大，
修好它們連帶使漂移項偏誤下降三分之二。人數五萬時扭曲比例升至 74\% 至
77\%，已超過崩潰點，同一年齡的 $\alpha$ 偏誤僅由 0.8658 降至 0.6925，
改善兩成；人數一萬時比例達 93\% 至 95\%，改善僅 1\%，漂移項甚至變差。

這個結果為第貳節之四所述的論點提供了直接的量化證據。本研究的問題不是
古典意義的污染——在鄉鎮市區的規模下，被扭曲的並非少數幾格，而是幼年組
的絕大多數格。以污染為設計目標的工具在此失效，不是因為它們不夠好，而是
因為污染比例已超出其設計的適用範圍。換言之，\textbf{穩健損失的作用範圍
有一個可以事前計算的上界，即診斷\cn{2}所量出的零格比例是否超過 50\%}。

加權的必要性則在三個規模下都成立。變體 G 與 F 的差別只在權重，而 G 的
漂移項偏誤在人數五萬時為 0.1591、F 為 0.3110，相差近一倍；$\beta$ 的
平方誤差和亦由 5.349 降至 1.548。這與第一小節「取對數傷 $\alpha_x$、
不加權傷漂移項」的分工完全一致，並顯示檢查函數與資訊加權是可以疊加的
兩件事。

最後須指出兩項限制。其一，變體 F 的 $\beta$ 平方誤差和在人數一萬與五萬時
反而大於標準 LC（10.864 對 7.916、5.349 對 4.449），亦即檢查函數在改善
$\alpha$ 的同時損害了 $\beta$，這是穩健損失的效率代價在本問題上的具體
形式。其二，即使是表現最佳的變體 G，其 $\alpha$ 最大偏誤在人數五萬時
仍為 0.6925，遠高於 Firth 的 0.0680（表~\ref{tab:est}）。兩者的差別在於
處理的層次：檢查函數改變的是損失的形狀，但零格替代所產生的錯誤資訊仍在
資料裡；Firth 則根本不取對數、不需替代零格，從源頭避免了該筆錯誤資訊。
\textbf{就本文的問題而言，換分配假設仍優於換損失函數。}

\begin{table}[H]
\centering
\caption{穩健程度 $c$ 的掃描（$\tau=0.5$，等權重，重複 100 次，取中位數）}
\small
\begin{tabular}{llrrrrr}
\toprule
$N$ & 估計量 & $\alpha$ 中位 & $\alpha$ 最大 & SSE($\beta$) & cor($\beta$) & 漂移偏誤\\
\midrule
\multirow{5}{*}{\shortstack[l]{$10^4$\\ 零格 206/528}}
 & A\quad $c=\infty$（標準 LC）  & 0.1612 & 2.3309 & \textbf{7.916} & $-0.074$ & \textbf{0.3436}\\
 & $c=3.0$                      & 0.1509 & 2.3309 & 8.638 & $-0.073$ & 0.3450\\
 & $c=1.345$（Huber 標準）      & 0.1427 & 2.3156 & 8.233 & $-0.079$ & 0.3592\\
 & $c=0.7$                      & \textbf{0.0995} & 2.3053 & 10.563 & $-0.079$ & 0.3604\\
 & F\quad $c=0$（中位數）        & 0.1015 & \textbf{2.3020} & 9.562 & $-0.063$ & 0.3494\\
\midrule
\multirow{5}{*}{\shortstack[l]{$5\times10^4$\\ 零格 100/528}}
 & A\quad $c=\infty$            & 0.0611 & 0.8658 & \textbf{4.449} & 0.037 & 0.3339\\
 & $c=3.0$                      & 0.0630 & 0.8652 & 5.363 & 0.052 & 0.3353\\
 & $c=1.345$                    & 0.0412 & 0.8006 & 5.302 & 0.042 & 0.3026\\
 & $c=0.7$                      & \textbf{0.0319} & 0.7351 & 4.920 & 0.023 & \textbf{0.2886}\\
 & F\quad $c=0$                 & 0.0393 & \textbf{0.6925} & 5.447 & 0.007 & 0.3150\\
\midrule
\multirow{5}{*}{\shortstack[l]{$2\times10^5$\\ 零格 29/528}}
 & A\quad $c=\infty$            & 0.0231 & 0.1846 & 1.400 & 0.379 & 0.2117\\
 & $c=3.0$                      & 0.0197 & 0.1656 & 1.449 & 0.418 & 0.1784\\
 & $c=1.345$                    & 0.0108 & \textbf{0.1231} & 1.010 & 0.468 & 0.1109\\
 & $c=0.7$                      & 0.0109 & 0.1289 & \textbf{0.900} & \textbf{0.478} & \textbf{0.0670}\\
 & F\quad $c=0$                 & \textbf{0.0050} & 0.1435 & 1.110 & 0.411 & 0.1111\\
\bottomrule
\end{tabular}
\label{tab:mqc}
\end{table}

第一項預測成立，但形式比原先設想的溫和。$\alpha$ 偏誤中位數的最適 $c$
隨人口規模單調下移：人數一萬與五萬時落在 $c=0.7$（0.0995 與 0.0319），
人數二十萬時落在 $c=0$（0.0050），亦即零格比例愈低、愈該把穩健程度
推到底。更值得注意的是\textbf{內部最適的存在}：在三個規模下都有一個
介於兩端之間的 $c$ 同時勝過兩端。人數二十萬時 $c=0.7$ 的漂移項偏誤為
0.0670，而兩個端點分別為 0.1111（中位數）與 0.2117（最小平方），
亦即中位數只走到一半；人數五萬時 $c=0.7$ 的漂移項偏誤 0.2886 亦優於
兩端的 0.3150 與 0.3339。前一小節以兩點作出的判斷因而漏掉了中間這一段，
而漏掉的原因正是把穩健當成開關而非刻度。

第三項預測亦成立，且結果超出前一小節所達到的水準。表~\ref{tab:mqK}
把 Huber 損失與資訊加權疊加，並以前一小節表現最佳的變體 G（加權中位數）
為比較基準。

\begin{table}[H]
\centering
\caption{穩健損失與資訊加權的疊加（$\tau=0.5$，重複 100 次）}
\small
\begin{tabular}{llrrrr}
\toprule
$N$ & 估計量 & $\alpha$ 中位 & $\alpha$ 最大 & SSE($\beta$) & 漂移偏誤\\
\midrule
\multirow{3}{*}{$10^4$}
 & G\quad 加權中位數（前節）        & 0.0733 & 2.3020 & 4.950 & 0.2986\\
 & $c=0$ $+$ 加權（本節重算）       & 0.0901 & 2.3020 & 7.696 & 0.3188\\
 & $c=1.345$ $+$ 加權              & 0.1431 & 2.3151 & \textbf{3.759} & \textbf{0.2579}\\
\midrule
\multirow{3}{*}{$5\times10^4$}
 & G\quad 加權中位數（前節）        & 0.0187 & \textbf{0.6925} & 1.548 & 0.1591\\
 & $c=0$ $+$ 加權（本節重算）       & 0.0353 & \textbf{0.6925} & 1.640 & \textbf{0.1339}\\
 & $c=1.345$ $+$ 加權              & 0.0351 & 0.7853 & \textbf{1.091} & 0.1404\\
\midrule
\multirow{3}{*}{$2\times10^5$}
 & G\quad 加權中位數（前節）        & 0.0032 & 0.1818 & 0.536 & 0.0428\\
 & $c=0$ $+$ 加權（本節重算）       & 0.0054 & 0.1708 & 0.560 & \textbf{0.0204}\\
 & $c=1.345$ $+$ 加權              & 0.0106 & \textbf{0.1267} & \textbf{0.411} & 0.0336\\
\bottomrule
\end{tabular}
\label{tab:mqK}
\end{table}

以 Huber 損失取代中位數並保留資訊加權後，$\beta$ 的平方誤差和在三個
規模下分別由 4.950、1.548、0.536 降至 3.759、1.091、0.411，
改善兩成至三成；漂移項偏誤在人數一萬時由 0.2986 降至 0.2579，
在另兩個規模則與中位數版相當。機制是效率：中位數只用到排序資訊，
Huber 損失在殘差小於截點的範圍內仍使用平方損失，因而保留了那些格的
完整資訊，而 $\beta$ 的估計恰恰依賴跨格的相對大小。$\alpha$ 偏誤中位數
則以中位數版較佳，這與穩健程度的取捨一致：$c$ 放大換到效率、犧牲了
一部分對零格扭曲的抵抗力。
\subsection{非對稱參數與單側截點的完整掃描}

\begin{table}[H]
\centering
\caption{非對稱參數 $\tau$ 與單側截點的掃描（$c=1.345$，等權重）}
\small
\begin{tabular}{llrrrrr}
\toprule
$N$ & 估計量 & $\alpha$ 中位 & $\alpha$ 最大 & SSE($\beta$) & cor($\beta$) & 漂移偏誤\\
\midrule
\multirow{7}{*}{$10^4$}
 & $\tau=0.3$              & 0.3692 & 2.3070 & 9.369 & $-0.040$ & 0.3344\\
 & $\tau=0.4$              & 0.2521 & 2.3105 & 8.551 & $-0.066$ & 0.3487\\
 & $\tau=0.5$              & 0.1427 & \textbf{2.3156} & \textbf{8.233} & $-0.079$ & 0.3592\\
 & $\tau=0.6$              & \textbf{0.0542} & 2.3249 & 9.365 & $-0.075$ & 0.3549\\
 & $\tau=0.7$              & 0.1500 & 2.3461 & 9.463 & $-0.093$ & 0.3589\\
 & 只壓右尾 $c^+\!=\!1.345$ & 0.2361 & 2.3188 & \textbf{7.318} & $-0.053$ & \textbf{0.3310}\\
 & 只壓左尾 $c^-\!=\!1.345$ & \textbf{0.0566} & 2.3309 & 13.065 & $-0.067$ & 0.3636\\
\midrule
\multirow{7}{*}{$5\times10^4$}
 & $\tau=0.3$              & 0.1903 & \textbf{0.7363} & 5.512 & 0.008 & \textbf{0.2857}\\
 & $\tau=0.4$              & 0.1309 & 0.7616 & 5.657 & 0.017 & 0.2880\\
 & $\tau=0.5$              & 0.0412 & 0.8006 & 5.302 & 0.042 & 0.3026\\
 & $\tau=0.6$              & \textbf{0.0219} & 0.8564 & \textbf{4.736} & 0.070 & 0.2878\\
 & $\tau=0.7$              & 0.0972 & 0.9504 & 5.426 & \textbf{0.084} & 0.2954\\
 & 只壓右尾                & 0.1027 & 0.8036 & 4.935 & $-0.009$ & 0.2980\\
 & 只壓左尾                & \textbf{0.0226} & 0.8658 & 5.194 & 0.059 & 0.3308\\
\midrule
\multirow{7}{*}{$2\times10^5$}
 & $\tau=0.3$              & 0.0923 & 0.4145 & 1.502 & 0.322 & 0.1385\\
 & $\tau=0.4$              & 0.0495 & 0.2795 & 1.167 & 0.401 & 0.1238\\
 & $\tau=0.5$              & 0.0108 & 0.1231 & \textbf{1.010} & 0.468 & 0.1109\\
 & $\tau=0.6$              & 0.0232 & \textbf{0.0589} & 1.021 & \textbf{0.502} & \textbf{0.0989}\\
 & $\tau=0.7$              & 0.0587 & 0.2390 & 1.158 & 0.492 & 0.1136\\
 & 只壓右尾                & 0.0292 & 0.3604 & 1.674 & 0.234 & 0.1903\\
 & 只壓左尾                & \textbf{0.0047} & 0.1489 & 1.323 & 0.455 & 0.1601\\
\bottomrule
\end{tabular}
\label{tab:mqtau}
\end{table}
\subsection{逐年穩健化}

實作上依每一年整條曲線的積分平方誤差
$v_t=\sum_x(\log m_{x,t}-\alpha_x-\beta_x\kappa_t)^2$ 判斷該年是否離群，
以 $v_t$ 的中位數與 MAD 作穩健定位與尺度，再以逐年權重重解秩一結構，
反覆至收斂。降權規則有兩種：平滑降權
$w_t=\min\{1,\lambda\,\mathrm{MAD}/(v_t-\mathrm{med})_+\}$，以及硬性剔除
$v_t$ 最大的若干年。

\begin{table}[H]
\centering
\caption{對離群年份穩健化的效果（重複 100 次，取中位數）}
\small
\begin{tabular}{llrrrrr}
\toprule
$N$ & 估計量 & 降權年數 & $\alpha$ 中位 & $\alpha$ 最大 & SSE($\beta$) & 漂移偏誤\\
\midrule
\multirow{4}{*}{$10^4$}
 & A\quad 標準 LC        & ---  & 0.1612 & 2.3309 & 7.916 & 0.3436\\
 & 平滑降權 $\lambda=3$  & 0.67 & 0.1609 & 2.3310 & 7.904 & 0.3436\\
 & 剔除 2 年             & 2    & 0.1597 & 2.3335 & 8.065 & 0.3544\\
 & 剔除 4 年             & 4    & 0.1588 & 2.3366 & 8.523 & 0.3449\\
\midrule
\multirow{4}{*}{$5\times10^4$}
 & A\quad 標準 LC        & ---  & 0.0611 & 0.8658 & 4.449 & 0.3339\\
 & 平滑降權 $\lambda=3$  & 0.59 & 0.0638 & 0.8658 & 4.835 & 0.3366\\
 & 剔除 2 年             & 2    & 0.0588 & 0.8501 & 4.977 & 0.3307\\
 & 剔除 4 年             & 4    & 0.0554 & 0.8658 & 5.671 & 0.3343\\
\midrule
\multirow{4}{*}{$2\times10^5$}
 & A\quad 標準 LC        & ---  & 0.0231 & 0.1846 & 1.400 & 0.2117\\
 & 平滑降權 $\lambda=3$  & 0.66 & 0.0242 & 0.1854 & 1.465 & 0.2160\\
 & 剔除 2 年             & 2    & 0.0242 & 0.1738 & 1.490 & 0.2155\\
 & 剔除 4 年             & 4    & 0.0218 & 0.1661 & 1.622 & 0.2301\\
\bottomrule
\end{tabular}
\label{tab:hu}
\end{table}

\subsection{區間推論的完整結果}

全部已知，區間只反映時間參數未來的創新項，此為任何方法都無法低於的理論
下限，其涵蓋率應接近名目的 0.80，可作為整個實驗的效度檢查。變體 P 以估計
出的時間參數與其變異數代入，即現行實務作法，P 與 O 的差距即為以估計值
取代真值所須付出的權衡。變體 B 在 P 之上以拔靴法納入參數估計的不確定性，
若參數不確定性是區間失效的主因，B 應明顯優於 P。變體 S 則在 B 之上先行
修勻，專門檢驗 Basellini et al.（2023）所提出而未量化的疑慮。預測期取
$h=10$、名目水準 80\%，目標為涵蓋真模型在觀測的時間參數路徑下的零歲
平均餘命。

閱讀時須同時看兩個指標。平均寬度反映區間的大小，涵蓋率反映區間是否擺對
位置，兩者必須合看——寬度過大而涵蓋率仍低，即代表位置錯誤而非寬度不足。

\begin{table}[H]
\centering
\caption{預測區間寬度與涵蓋率的來源分解（名目 80\%）}
\small
\begin{tabular}{rlcc}
\toprule
$N$ & 變體 & 平均寬度 & 涵蓋率\\
\midrule
50,000 & O\quad 僅未來創新項（不可約下限） & 1.743 & 0.85\\
       & P\quad 插入式（現行作法） & 5.401 & 0.38\\
       & B\quad P ＋參數不確定性 & 6.166 & 0.24\\
       & S\quad 修勻 ＋ B & 1.683 & 0.66\\
\midrule
200,000 & O & 1.720 & 0.82\\
        & P & 5.651 & 0.82\\
        & S & 1.515 & 0.63\\
\midrule
1,000,000 & O & 1.738 & 0.77\\
          & P & 4.390 & 0.94\\
          & S & 1.194 & 0.45\\
\midrule
11,900,000 & O & 1.723 & 0.87\\
           & P & 1.950 & 0.88\\
           & S & \textbf{0.529} & \textbf{0.38}\\
\bottomrule
\end{tabular}
\label{tab:idecomp}
\end{table}

結果有三。其一，小人口的預測區間同時太寬且涵蓋不足，寬度與涵蓋率解耦。
人數五萬時現行作法的區間寬達 5.401，而不可約下限僅 1.743，涵蓋率卻只有
0.38。寬度被灌水是因為 $\hat\kappa_t$ 含估計雜訊、使一階差變大而
$\hat\sigma$ 被高估；涵蓋失敗則因漂移項被壓平，使區間擺錯位置。換言之，
小人口的問題是偏誤而非寬度，加寬區間無濟於事。此一結果限定了文獻長期
以來「LC 模型預測區間過窄」的診斷在小人口的適用性。

其二，拔靴法的參數不確定性幾乎沒有貢獻，有時甚至更糟。因為拔靴法是繞著
已經偏掉的配適值重抽，能偵測到的偏誤有限——第伍節的偏誤校正實驗顯示其
偵測比例約為七成，但代價是變異數加倍。

其三，修勻確實壓低區間寬度，而且代價隨人口增大而惡化。人數一千一百九十萬
時變體 S 的寬度僅 0.529、涵蓋率降至 0.38，方向與直覺相反。這為
Basellini et al.（2023）的疑慮提供了量化證據，亦即本文的議題二。

直接的作法是以拔靴法估計 $\mathrm{Var}(\Delta u^*)$ 後相減，但兩個大噪音量
相減並不穩定，是動差法去偏的典型病症。本文改以狀態空間表述，令觀測方程
為 $\hat\kappa_t=\kappa_t+u_t$、$u_t\sim N(0,\tau_t^2)$，狀態方程為
$\kappa_t=\kappa_{t-1}+\delta+\varepsilon_t$、
$\varepsilon_t\sim N(0,\sigma^2)$。此一表述有兩項優勢。其一，$\tau_t^2$
不必估計，給定年齡參數後 $\kappa_t$ 的 Fisher 資訊為
$\sum_x\hat\mu_{x,t}\hat\beta_x^2$，而這一項正是卜瓦松 LC 迭代中 $\kappa$
更新式的分母，故取得它沒有額外成本。其二，$\sigma^2$ 是在概似中被估計
出來的，而非兩個估計量相減的餘數，結構上不可能為負，動差法的病症自動
消失。變異成分的估計採第貳節所述的 REML，因為最大概似在估計變異成分時
未扣除估計固定效果所消耗的自由度，會系統性低估。

\begin{table}[H]
\centering
\caption{$\sigma$ 的三種估計方式（真值 $0.6791$，重複 100 次）}
\small
\begin{tabular}{llrrrc}
\toprule
$N$ & 方法 & 中位數 & 偏誤 & IQR & 可行率\\
\midrule
\multirow{3}{*}{$5\times10^4$}
 & $\mathrm{sd}(\Delta\hat\kappa)$（現行） & 1.4950 & $+0.8158$ & 0.896 & 1.00\\
 & 狀態空間 ML                             & 0.4029 & $-0.2763$ & 0.429 & 1.00\\
 & \textbf{狀態空間 REML}                  & \textbf{0.5753} & $\mathbf{-0.1039}$ & \textbf{0.357} & 1.00\\
\midrule
\multirow{3}{*}{$2\times10^5$}
 & $\mathrm{sd}(\Delta\hat\kappa)$ & 1.0054 & $+0.3263$ & 0.282 & 1.00\\
 & 狀態空間 ML                     & 0.5392 & $-0.1400$ & 0.172 & 1.00\\
 & \textbf{狀態空間 REML}          & \textbf{0.5970} & $\mathbf{-0.0822}$ & 0.194 & 1.00\\
\midrule
\multirow{3}{*}{$10^6$}
 & $\mathrm{sd}(\Delta\hat\kappa)$ & 0.7536 & $+0.0745$ & 0.122 & 1.00\\
 & 狀態空間 ML                     & 0.6318 & $-0.0473$ & 0.107 & 1.00\\
 & \textbf{狀態空間 REML}          & \textbf{0.6551} & $\mathbf{-0.0240}$ & 0.111 & 1.00\\
\bottomrule
\end{tabular}
\label{tab:sigma}
\end{table}

表~\ref{tab:sigma} 以真值 0.6791 為基準。樸素估計量在人數五萬時的中位數
為 1.4950，偏誤高達 $+0.82$，即被灌水超過一倍；狀態空間 REML 為 0.5753，
偏誤 $-0.10$，四分位距亦由 0.896 降至 0.357。最大概似版本則如理論所預期
向下偏誤較大（$-0.28$），在三個規模下皆劣於 REML，故本文採用 REML。
須注意三種方法在人數一百萬時的差距已大幅縮小，樸素估計量的偏誤降至
$+0.07$，亦即變異成分分離的必要性隨人口規模下降。

\subsection{初始值的不確定性}

$\sigma$ 校正之後，寬度的問題大致解決，但位置的問題仍在。第參節已指出
漂移項是偏誤主導，而預測的起點 $\hat\kappa_T$ 是單一年份的估計值，其誤差
不隨配適期長度縮小。狀態空間表述在此提供了一項現成的資源：Kalman 濾波在
最後一點給出的不只是 $\kappa_{T|T}$，還有它的變異數 $P_T$，把後者納入
預測分布即可反映初始值本身的不確定性。

\begin{table}[H]
\centering
\caption{變異成分分離與初始值不確定性對預測區間的影響（$h=10$、名目 80\%、重複 100 次）}
\small
\begin{tabular}{llrc}
\toprule
$N$ & 變體 & 寬度 & 涵蓋率\\
\midrule
\multirow{4}{*}{$5\times10^4$}
 & O\quad 不可約下限（參數已知）        & 2.093 & 0.74\\
 & P\quad 現行作法（SVD 插入式）        & 8.949 & 0.27\\
 & R\quad 卜瓦松 $+$ 狀態空間 REML      & \textbf{1.610} & 0.47\\
 & RK\ \ R $+$ Kalman 初始值變異數      & 1.707 & \textbf{0.54}\\
\midrule
\multirow{4}{*}{$2\times10^5$}
 & O & 2.106 & 0.82\\
 & P & 7.616 & 0.64\\
 & R & \textbf{1.860} & \textbf{0.72}\\
 & RK & 1.901 & \textbf{0.72}\\
\midrule
\multirow{4}{*}{$10^6$}
 & O & 2.103 & 0.80\\
 & P & 6.681 & 0.93\\
 & R & \textbf{2.025} & \textbf{0.77}\\
 & RK & 2.019 & \textbf{0.77}\\
\bottomrule
\end{tabular}
\label{tab:ival}
\end{table}

表~\ref{tab:ival} 的結果可分三點讀。其一，現行作法的失效在此再次確認，
人數五萬時寬度 8.949 為不可約下限 2.093 的 4.3 倍，涵蓋率卻只有 0.27。
其二，以卜瓦松配適搭配狀態空間 REML 之後，寬度由 8.949 降至 1.610、
涵蓋率由 0.27 升至 0.47；人數二十萬時由 7.616 與 0.64 變為 1.860 與 0.72，
人數一百萬時則為 2.025 與 0.77，已與不可約下限的 2.103 與 0.80 十分接近。
須指出人數五萬時的 1.610 已低於下限 2.093，亦即 REML 在該規模略為低估
$\sigma$，區間反而稍窄，這與表~\ref{tab:sigma} 的 $-0.10$ 偏誤一致。
其三，加計初始值變異數 $P_T$ 在小人口有效而在大人口無效：人數五萬時
涵蓋率由 0.47 升至 0.54、寬度僅增加 6\%，人數二十萬與一百萬時則幾無變化。
這符合預期，因為 $P_T$ 的大小取決於 $\hat\kappa_T$ 的估計精度，人口愈大
精度愈高，該項的貢獻自然愈小。

須誠實指出殘餘的缺口。人數五萬時涵蓋率仍只有 0.54，距名目的 0.80 尚遠，
而寬度已貼合下限，因此缺口不在寬度而在位置，亦即漂移項的偏誤。這一點
無法由區間方法解決，須回到第伍節的點估計。

綜合前述，小人口區間的兩個病灶可分別處理：位置由漂移項的偏誤造成，
寬度由 $\hat\sigma$ 的灌水造成。

\begin{table}[H]
\centering
\caption{SVD 與卜瓦松最大概似的漂移項估計（真值 $-0.3678$）}
\small
\begin{tabular}{rcc}
\toprule
$N$ & SVD 的漂移項 & 卜瓦松 MLE 的漂移項\\
\midrule
50,000    & $-0.032$（偏誤 $+0.336$） & $\mathbf{-0.377}$（偏誤 $-0.010$）\\
200,000   & $-0.194$（偏誤 $+0.174$） & $\mathbf{-0.369}$（偏誤 $-0.002$）\\
1,000,000 & $-0.365$ & $-0.367$\\
\bottomrule
\end{tabular}
\label{tab:drift}
\end{table}

表~\ref{tab:drift} 顯示漂移項衰減並非小人口的固有性質，而是 SVD 對
$\log(D/E)$ 做最小平方所產生的結果。人數五萬時 SVD 估得 $-0.032$、偏誤
$+0.336$，卜瓦松最大概似則為 $-0.377$、偏誤僅 $-0.010$。卜瓦松架構不取
對數，衰減即告消失。此一發現獨立支持了 Basellini et al.（2023）
「換到卜瓦松架構後第二階段調整不再必要」的論點，並把理由由「配適總死亡數
依照模型假設成立」推進到「第一階段的估計本身即無偏」。

\subsection{平均餘命誤差的年齡歸因}

\begin{table}[H]
\centering
\caption{零歲平均餘命誤差的年齡別歸因（$N=5\times10^4$，重複 100 次；
單位為歲，負值表示該年齡使 $e_0$ 被低估）}
\small
\begin{tabular}{lrrrr}
\toprule
年齡組 & 真值死亡率 & 標準 LC 的貢獻 & 占總誤差 & Firth 的貢獻\\
\midrule
0       & $3.53\times10^{-3}$ & $-0.0028$ & 0.2\% & $-0.0097$\\
1--4    & $1.55\times10^{-4}$ & $-0.1019$ & 7.4\% & $-0.0018$\\
5--9    & $9.33\times10^{-5}$ & $-0.0750$ & 5.5\% & $-0.0086$\\
10--14  & $1.07\times10^{-4}$ & $-0.0605$ & 4.4\% & $-0.0011$\\
15--19  & $1.98\times10^{-4}$ & $-0.0435$ & 3.2\% & $-0.0090$\\
\addlinespace
50--54  & $2.37\times10^{-3}$ & $-0.0157$ & 1.1\% & $-0.0083$\\
60--64  & $4.86\times10^{-3}$ & $-0.0876$ & 6.4\% & $-0.0105$\\
65--69  & $7.47\times10^{-3}$ & $-0.1555$ & 11.3\% & $-0.0086$\\
70--74  & $1.31\times10^{-2}$ & \textbf{$-0.1856$} & \textbf{13.5\%} & $-0.0175$\\
75--79  & $2.46\times10^{-2}$ & \textbf{$-0.1855$} & \textbf{13.5\%} & $-0.0163$\\
80--84  & $4.73\times10^{-2}$ & \textbf{$-0.1872$} & \textbf{13.6\%} & $-0.0413$\\
85--89  & $8.99\times10^{-2}$ & $-0.1286$ & 9.4\% & $-0.0176$\\
95--99  & $2.59\times10^{-1}$ & $-0.0083$ & 0.6\% & $-0.0068$\\
\midrule
30 歲以下合計 & & $-0.3049$ & \textbf{20.6\%} & $-0.0464$\\
60$\sim$89 歲合計 & & $-0.9300$ & \textbf{62.8\%} & $-0.1118$\\
\bottomrule
\end{tabular}
\label{tab:e0attr}
\end{table}
\subsection{人口規模的建議門檻}

\begin{table}[H]
\centering
\caption{依事前診斷決定的估計策略（配適期 24 年）}
\small
\begin{tabular}{p{2.5cm}p{2.3cm}p{4.2cm}p{4.2cm}}
\toprule
診斷\cn{4}比值 & 約當人口 & $\alpha_x$ 與漂移項 & $\beta_x$\\
\midrule
$<0.6$
 & 2 萬以下
 & Firth 卜瓦松（必要，非選配）；診斷\cn{1}的最小死亡中位數低於 5 時尤然
 & \textbf{不予估計}；估出的方向與隨機向量無法區分\\
\addlinespace
$0.6\sim1.2$
 & 2 萬$\sim$20 萬
 & Firth 卜瓦松；方法選擇的效益在此最大
 & 卜瓦松或 Firth 搭配 Fay-Herriot 收縮；相關係數可由 0.04 提升至 0.54\\
\addlinespace
$1.2\sim2$
 & 20 萬$\sim$50 萬
 & 卜瓦松最大概似即可
 & 卜瓦松或以 $\hat\mu$ 加權的 SVD，加 Fay-Herriot 收縮\\
\addlinespace
$>2$
 & 50 萬以上
 & 各方法差異已小，SVD 亦可
 & 各方法差異已小；HeteroPCA 於此亦有效\\
\bottomrule
\end{tabular}
\label{tab:boundary}
\end{table}
\section*{參考文獻}
\addcontentsline{toc}{section}{參考文獻}

\small
\noindent\textbf{一、中文部分}
\begin{itemize}[leftmargin=2em, label={}, itemsep=3pt]
  \item[] 王信忠、金碩、余清祥（2012）。小區域死亡率推估之研究。人口學刊，
        45：121-154。doi:10.6191/JPS.2012.11
  \item[] 王信忠、余清祥、王子瑜（2017）。臺灣原住民死亡率暨生命表編撰研究。
        人口學刊，55：99-131。doi:10.6191/JPS.2017.55.03
  \item[] 余清祥、王信忠、陳譽騰（2021）。年輪變動比用於小區域人口推估的探討。
        人口學刊，63：99-133。doi:10.6191/JPS.202112\_(63).0003
  \item[] 余清祥、王信忠、呂靖翎（2025）。平均餘命與標準化死亡率之相關分析。
        人口學刊，85-116。
  \item[] 陳政勳、余清祥（2010）。小區域人口推估研究：臺北市、雲嘉兩縣、
        澎湖縣的實證分析。人口學刊，41：153-183。doi:10.6191/JPS.2010.9
\end{itemize}

\vspace{4pt}
\noindent\textbf{二、英文部分}
\begin{itemize}[leftmargin=2em, label={}, itemsep=3pt]
  \item[] Albert, A., and Anderson, J. A. 1984. ``On the Existence of Maximum
        Likelihood Estimates in Logistic Regression Models.'' \textit{Biometrika}
        71(1): 1-10.
  \item[] Alho, J. M. 1992. ``Modeling and Forecasting U.S. Mortality:
        Comment.'' \textit{Journal of the American Statistical Association}
        87(419): 673-674.
  \item[] Arriaga, E. E. 1984. ``Measuring and Explaining the Change in Life
        Expectancies.'' \textit{Demography} 21(1): 83-96.
  \item[] Baik, J., Ben Arous, G., and P\'{e}ch\'{e}, S. 2005. ``Phase Transition of
        the Largest Eigenvalue for Nonnull Complex Sample Covariance Matrices.''
        \textit{Annals of Probability} 33(5): 1643-1697.
  \item[] Basellini, U., C. G. Camarda, and H. Booth. 2023. ``Thirty Years On:
        A Review of the Lee--Carter Method for Forecasting Mortality.''
        \textit{International Journal of Forecasting} 39(3): 1033-1049.
  \item[] Basellini, U., and C. G. Camarda. 2025. ``Forecasting Cohort Mortality:
        Lee-Carter Methods and CCP-Splines.'' SocArXiv preprint, October 9.
  \item[] Benaych-Georges, F., and Nadakuditi, R. R. 2012. ``The Singular Values and
        Vectors of Low Rank Perturbations of Large Rectangular Random Matrices.''
        \textit{Journal of Multivariate Analysis} 111: 120-135.
  \item[] Bianchi, A., E. Fabrizi, N. Salvati, and N. Tzavidis. 2018.
        ``Estimation and Testing in M-quantile Regression with Applications to
        Small Area Estimation.'' \textit{International Statistical Review}
        86(3): 541-570.
  \item[] Booth, H., J. Maindonald, and L. Smith. 2002. ``Applying Lee--Carter
        under Conditions of Variable Mortality Decline.''
        \textit{Population Studies} 56(3): 325-336.
  \item[] Booth, H., L. Tickle, and L. Smith. 2005. ``Evaluation of the
        Variants of the Lee-Carter Method of Forecasting Mortality:
        A Multi-country Comparison.''
        \textit{New Zealand Population Review} 31(1): 13-34.
  \item[] Booth, H., R. J. Hyndman, L. Tickle, and P. De Jong. 2006.
        ``Lee--Carter Mortality Forecasting: A Multi-country Comparison of
        Variants and Extensions.'' \textit{Demographic Research} 15: 289-310.
  \item[] Breckling, J., and R. Chambers. 1988. ``M-quantiles.''
        \textit{Biometrika} 75(4): 761-771.
  \item[] Brouhns, N., M. Denuit, and J. K. Vermunt. 2002. ``A Poisson
        Log-bilinear Regression Approach to the Construction of Projected
        Lifetables.'' \textit{Insurance: Mathematics and Economics}
        31(3): 373-393.
  \item[] B\"uhlmann, P., and S. van de Geer. 2011. \textit{Statistics for
        High-Dimensional Data: Methods, Theory and Applications.} Springer.
  \item[] Cairns, A. J. G., D. Blake, K. Dowd, G. D. Coughlan, and
        M. Khalaf-Allah. 2011. ``Bayesian Stochastic Mortality Modelling for
        Two Populations.'' \textit{ASTIN Bulletin} 41(1): 29-59.
  \item[] Camarda, C. G., and U. Basellini. 2021. ``Smoothing, Decomposing and
        Forecasting Mortality Rates.'' \textit{European Journal of Population}
        37(3): 569-602.
  \item[] Catania, L., and A. Luati. 2020. ``Robust Estimation of a Location
        Parameter with the Integrated Hogg Function.''
        \textit{Statistics \& Probability Letters} 164: 108812.
  \item[] Chambers, R., and N. Tzavidis. 2006. ``M-quantile Models for Small
        Area Estimation.'' \textit{Biometrika} 93(2): 255-268.
  \item[] Chen, L., A. J. G. Cairns, and T. Kleinow. 2017. ``Small Population
        Bias and Sampling Effects in Stochastic Mortality Modelling.''
        \textit{European Actuarial Journal} 7(1): 193-230.
  \item[] Clogg, C. C., D. B. Rubin, N. Schenker, B. Schultz, and L. Weidman.
        1991. ``Multiple Imputation of Industry and Occupation Codes in Census
        Public-use Samples Using Bayesian Logistic Regression.''
        \textit{Journal of the American Statistical Association} 86(413): 68-78.
  \item[] Cox, D. R., and Snell, E. J. 1968. ``A General Definition of Residuals.''
        \textit{Journal of the Royal Statistical Society, Series B} 30(2): 248-265.
  \item[] Currie, I. D., M. Durban, and P. H. C. Eilers. 2004. ``Smoothing
        and Forecasting Mortality Rates.'' \textit{Statistical Modelling}
        4(4): 279-298.
  \item[] Currie, I. D. 2011. ``Modelling and Forecasting the Mortality of the
        Very Old.'' \textit{ASTIN Bulletin} 41(2): 419-427.
  \item[] Currie, I. D. 2013. ``Smoothing Constrained Generalized Linear Models
        with an Application to the Lee--Carter Model.''
        \textit{Scandinavian Actuarial Journal} 2016(4): 356-383.
  \item[] Delwarde, A., M. Denuit, and P. Eilers. 2007. ``Smoothing the
        Lee--Carter and Poisson Log-bilinear Models for Mortality Forecasting:
        A Penalized Log-likelihood Approach.''
        \textit{Statistical Modelling} 7(1): 29-48.
  \item[] Di Fonzo, T., and M. Marini. 2015. ``Reconciliation of Systems of
        Time Series according to a Growth Rates Preservation Principle.''
        \textit{Statistical Methods \& Applications} 24: 651-669.
  \item[] Djeundje, V. A. B., and I. D. Currie. 2010. ``Smoothing Dispersed
        Counts with Applications to Mortality Data.''
        \textit{Annals of Actuarial Science} 5(1): 33-52.
  \item[] Dokumentov, A., R. J. Hyndman, and L. Tickle. 2018. ``Bivariate
        Smoothing of Mortality Surfaces with Zero Counts: The Lasso2d Method.''
        \textit{Stat} 7(1): e199.
  \item[] Donoho, D. L., and Huber, P. J. 1983. ``The Notion of Breakdown Point.''
        In \textit{A Festschrift for Erich L. Lehmann}, 157-184. Belmont: Wadsworth.
  \item[] Efron, B., and Morris, C. 1973. ``Stein's Estimation Rule and Its
        Competitors---An Empirical Bayes Approach.''
        \textit{Journal of the American Statistical Association} 68(341): 117-130.
  \item[] Eilers, P. H. C., and B. D. Marx. 1996. ``Flexible Smoothing with
        B-splines and Penalties.'' \textit{Statistical Science} 11(2): 89-102.
  \item[] Fay, R. E., and Herriot, R. A. 1979. ``Estimates of Income for Small Places:
        An Application of James-Stein Procedures to Census Data.''
        \textit{Journal of the American Statistical Association} 74(366): 269-277.
  \item[] Firth, D. 1993. ``Bias Reduction of Maximum Likelihood Estimates.''
        \textit{Biometrika} 80(1): 27-38.
  \item[] Friedman, J., T. Hastie, and R. Tibshirani. 2010. ``Regularization
        Paths for Generalized Linear Models via Coordinate Descent.''
        \textit{Journal of Statistical Software} 33(1): 1-22.
  \item[] Gavish, M., and Donoho, D. L. 2017. ``Optimal Shrinkage of Singular Values.''
        \textit{IEEE Transactions on Information Theory} 63(4): 2137-2152.
  \item[] Goodman, L. A. 1979. ``Simple Models for the Analysis of Association in
        Cross-Classifications Having Ordered Categories.''
        \textit{Journal of the American Statistical Association} 74(367): 537-552.
  \item[] Hampel, F. R. 1974. ``The Influence Curve and Its Role in Robust Estimation.''
        \textit{Journal of the American Statistical Association} 69(346): 383-393.
  \item[] Hampel, F. R., Ronchetti, E. M., Rousseeuw, P. J., and Stahel, W. A. 1986.
        \textit{Robust Statistics: The Approach Based on Influence Functions}.
        New York: Wiley.
  \item[] Heinze, G., and Schemper, M. 2002. ``A Solution to the Problem of Separation
        in Logistic Regression.'' \textit{Statistics in Medicine} 21(16): 2409-2419.
  \item[] Heitjan, D. F., and Rubin, D. B. 1991. ``Ignorability and Coarse Data.''
        \textit{Annals of Statistics} 19(4): 2244-2253.
  \item[] Hoerl, A. E., and R. W. Kennard. 1970. ``Ridge Regression: Biased
        Estimation for Nonorthogonal Problems.''
        \textit{Technometrics} 12(1): 55-67.
  \item[] H\"ormann, S., \L. Kidzi\'nski, and M. Hallin. 2015. ``Dynamic
        Functional Principal Components.'' \textit{Journal of the Royal
        Statistical Society, Series B} 77(2): 319-348.
  \item[] Huber, P. J. 1964. ``Robust Estimation of a Location Parameter.''
        \textit{Annals of Mathematical Statistics} 35(1): 73-101.
  \item[] Huber, P. J., and Ronchetti, E. M. 2009. \textit{Robust Statistics},
        2nd ed. Hoboken: Wiley.
  \item[] Hunt, A., and D. Blake. 2020. ``Identifiability in Age/Period
        Mortality Models.'' \textit{Annals of Actuarial Science} 14(2): 461-499.
  \item[] Hunt, A., and D. Blake. 2021. ``On the Structure and Classification of
        Mortality Models.'' \textit{North American Actuarial Journal}
        25(sup1): S215-S234.
  \item[] Hunter, D. R., and K. Lange. 2000. ``Quantile Regression via an MM
        Algorithm.'' \textit{Journal of Computational and Graphical Statistics}
        9(1): 60-77.
  \item[] Hyndman, R. J., and M. S. Ullah. 2007. ``Robust Forecasting of
        Mortality and Fertility Rates: A Functional Data Approach.''
        \textit{Computational Statistics \& Data Analysis} 51(10): 4942-4956.
  \item[] James, W., and C. Stein. 1961. ``Estimation with Quadratic Loss.''
        \textit{Proceedings of the Fourth Berkeley Symposium on Mathematical
        Statistics and Probability} 1: 361-379.
  \item[] Jarner, S. F., and E. M. Kryger. 2011. ``Modelling Adult Mortality in
        Small Populations: The SAINT Model.''
        \textit{ASTIN Bulletin} 41(2): 377-418.
  \item[] Johnstone, I. M. 2001. ``On the Distribution of the Largest Eigenvalue in
        Principal Components Analysis.'' \textit{Annals of Statistics} 29(2): 295-327.
  \item[] Koenker, R., and Bassett, G. 1978. ``Regression Quantiles.''
        \textit{Econometrica} 46(1): 33-50.
  \item[] Koenker, R., and Park, B. J. 1996. ``An Interior Point Algorithm for Nonlinear
        Quantile Regression.'' \textit{Journal of Econometrics} 71(1-2): 265-283.
  \item[] Koenker, R. 2005. \textit{Quantile Regression}. Cambridge: Cambridge
        University Press.
  \item[] Koenker, R. 2017. ``Quantile Regression: 40 Years On.''
        \textit{Annual Review of Economics} 9: 155-176.
  \item[] Kosmidis, I., and Firth, D. 2021. ``Jeffreys-prior Penalty, Finiteness and
        Shrinkage in Binomial-response Generalized Linear Models.''
        \textit{Biometrika} 108(1): 71-82.
  \item[] Lee, R. D., and L. R. Carter. 1992. ``Modeling and Forecasting U.S.
        Mortality.'' \textit{Journal of the American Statistical Association}
        87(419): 659-671.
  \item[] Lee, R., and T. Miller. 2001. ``Evaluating the Performance of the
        Lee--Carter Method for Forecasting Mortality.''
        \textit{Demography} 38(4): 537-549.
  \item[] Lee, J. D., Y. Sun, and M. A. Saunders. 2014. ``Proximal Newton-type
        Methods for Minimizing Composite Functions.''
        \textit{SIAM Journal on Optimization} 24(3): 1420-1443.
  \item[] Li, N., R. D. Lee, and S. Tuljapurkar. 2004. ``Using the Lee--Carter
        Method to Forecast Mortality for Populations with Limited Data.''
        \textit{International Statistical Review} 72(1): 19-36.
  \item[] Li, N., and R. Lee. 2005. ``Coherent Mortality Forecasts for a Group
        of Populations: An Extension of the Lee--Carter Method.''
        \textit{Demography} 42(3): 575-594.
  \item[] Li, J. S.-H., M. R. Hardy, and K. S. Tan. 2009. ``Uncertainty in
        Mortality Forecasting: An Extension to the Classical Lee-Carter
        Approach.'' \textit{ASTIN Bulletin} 39(1): 137-164.
  \item[] Lu, Y., and D. Zhu. 2023. ``Modelling Mortality: A Bayesian
        Factor-augmented VAR (FAVAR) Approach.''
        \textit{ASTIN Bulletin} 53(1): 29-61.
  \item[] Machado, J. A. F., and Santos Silva, J. M. C. 2005. ``Quantiles for Counts.''
        \textit{Journal of the American Statistical Association} 100(472): 1226-1237.
  \item[] Meinshausen, N. 2007. ``Relaxed Lasso.''
        \textit{Computational Statistics \& Data Analysis} 52(1): 374-393.
  \item[] M\"{u}ller, C. H., and Neykov, N. 2003. ``Breakdown Points of Trimmed
        Likelihood Estimators and Related Estimators in Generalized Linear Models.''
        \textit{Journal of Statistical Planning and Inference} 116(2): 503-519.
  \item[] Newey, W. K., and Powell, J. L. 1987. ``Asymmetric Least Squares Estimation
        and Testing.'' \textit{Econometrica} 55(4): 819-847.
  \item[] Paul, D. 2007. ``Asymptotics of Sample Eigenstructure for a Large Dimensional
        Spiked Covariance Model.'' \textit{Statistica Sinica} 17(4): 1617-1642.
  \item[] Pitt, D., J. Li, and T. K. Lim. 2018. ``Smoothing Poisson Common
        Factor Model for Projecting Mortality Jointly for Both Sexes.''
        \textit{ASTIN Bulletin: The Journal of the IAA} 48(2): 509-541.
  \item[] Quenouille, M. H. 1956. ``Notes on Bias in Estimation.''
        \textit{Biometrika} 43(3-4): 353-360.
  \item[] Rabbi, A. M. F., and S. Mazzuco. 2021. ``Mortality Forecasting with
        the Lee--Carter Method: Adjusting for Smoothing and Lifespan
        Disparity.'' \textit{European Journal of Population} 37(1): 97-120.
  \item[] Rao, J. N. K., and Molina, I. 2015. \textit{Small Area Estimation}, 2nd ed.
        Hoboken: Wiley.
  \item[] Renshaw, A. E., and S. Haberman. 2003. ``Lee--Carter Mortality
        Forecasting with Age-specific Enhancement.''
        \textit{Insurance: Mathematics and Economics} 33(2): 255-272.
  \item[] Rousseeuw, P. J. 1984. ``Least Median of Squares Regression.''
        \textit{Journal of the American Statistical Association} 79(388): 871-880.
  \item[] Santolino, M. 2020. ``The Lee-Carter Quantile Mortality Model.''
        \textit{Scandinavian Actuarial Journal} 2020(7): 614-633.
  \item[] Schuette, D. R. 1978. ``A Linear Programming Approach to
        Graduation.'' \textit{Transactions of Society of Actuaries}
        30: 407-431.
  \item[] Shang, H. L., H. Booth, and R. J. Hyndman. 2011. ``Point and Interval
        Forecasts of Mortality Rates and Life Expectancy: A Comparison of Ten
        Principal Component Methods.''
        \textit{Demographic Research} 25: 173-214.
  \item[] Stein, C. 1956. ``Inadmissibility of the Usual Estimator for the Mean of a
        Multivariate Normal Distribution.'' In \textit{Proceedings of the Third
        Berkeley Symposium on Mathematical Statistics and Probability} 1: 197-206.
  \item[] Stoeldraijer, L., C. van Duin, L. van Wissen, and F. Janssen. 2018.
        ``Comparing Strategies for Matching Mortality Forecasts to the Most
        Recently Observed Data.'' \textit{Genus} 74(1): 16.
  \item[] Tibshirani, R. 1996. ``Regression Shrinkage and Selection via the
        Lasso.'' \textit{Journal of the Royal Statistical Society, Series B}
        58(1): 267-288.
  \item[] Tseng, P. 2001. ``Convergence of a Block Coordinate Descent Method
        for Nondifferentiable Minimization.'' \textit{Journal of Optimization
        Theory and Applications} 109(3): 475-494.
  \item[] Wang, H. C., C. J. Yue, and C. T. Chong. 2018. ``Mortality Models and
        Longevity Risk for Small Populations.''
        \textit{Insurance: Mathematics and Economics} 78: 351-359.
  \item[] Wei, Y., A. Pere, R. Koenker, and X. He. 2006. ``Quantile Regression
        Methods for Reference Growth Charts.'' \textit{Statistics in Medicine}
        25(8): 1369-1382.
  \item[] Wilmoth, J. R. 1993. \textit{Computational Methods for Fitting and
        Extrapolating the Lee--Carter Model of Mortality Change.} Technical
        Report, Department of Demography, University of California, Berkeley.
  \item[] Wi\'sniowski, A., P. W. F. Smith, J. Bijak, J. Raymer, and
        J. J. Forster. 2015. ``Bayesian Population Forecasting: Extending the
        Lee--Carter Method.'' \textit{Demography} 52(3): 1035-1059.
  \item[] Xu, X., and X. Chen. 2018. ``A Practical Method of Robust Estimation
        in Case of Asymmetry.'' \textit{Journal of Statistical Theory and
        Practice} 12(2): 370-396.
  \item[] Yang, S. S., J. C. Yue, and H. C. Huang. 2010. ``Modeling Longevity
        Risks Using a Principal Component Approach: A Comparison with Existing
        Stochastic Mortality Models.''
        \textit{Insurance: Mathematics and Economics} 46(1): 254-270.
  \item[] Yue, J. C., H. C. Wang, and T. Y. Wang. 2019. ``Using Graduation to
        Modify the Estimation of Lee--Carter Model for Small Populations.''
        \textit{North American Actuarial Journal}.
        doi:10.1080/10920277.2019.1650288
  \item[] Zhang, J., B. Lin, and Y. Yang. 2022. ``Maximum Nonparametric Kernel
        Likelihood Estimation for Multiplicative Linear Regression Models.''
        \textit{Statistical Papers} 63: 885-918.
  \item[] Zhang, C.-H. 2010. ``Nearly Unbiased Variable Selection under
        Minimax Concave Penalty.'' \textit{The Annals of Statistics}
        38(2): 894-942.
  \item[] Zhang, A. R., Cai, T. T., and Wu, Y. 2022. ``Heteroskedastic PCA: Algorithm,
        Optimality, and Applications.'' \textit{Annals of Statistics} 50(1): 53-80.
  \item[] Zhang, L., and Y. Zhuang. 2026. ``What KAN Mortality Say: Smooth and
        Interpretable Mortality Modeling Using Kolmogorov--Arnold Networks.''
        \textit{ASTIN Bulletin} 56(1): 32-59.
  \item[] Zou, H., and T. Hastie. 2005. ``Regularization and Variable Selection
        via the Elastic Net.'' \textit{Journal of the Royal Statistical
        Society, Series B} 67(2): 301-320.

  \item[] Zou, H. 2006. ``The Adaptive Lasso and Its Oracle Properties.''
        \textit{Journal of the American Statistical Association}
        101(476): 1418-1429.
  \item[] Zou, H., and Yuan, M. 2008. ``Composite Quantile Regression and the Oracle
        Model Selection Theory.'' \textit{Annals of Statistics} 36(3): 1108-1126.
  \item[] Zoubir, A. M., V. Koivunen, E. Ollila, and M. Muma. 2018.
        \textit{Robust Statistics for Signal Processing}. Cambridge: Cambridge
        University Press.
\end{itemize}

\end{document}
